\documentclass[10.5pt,a4paper]{article}
\usepackage[margin=2.2cm,headheight=14pt]{geometry}
\usepackage{fontspec,unicode-math}
\setmainfont{Libertinus Serif}\setsansfont{IBM Plex Sans}\setmonofont[Scale=0.84]{IBM Plex Mono}\setmathfont{Libertinus Math}
\usepackage{microtype,graphicx,booktabs,tabularx,longtable,xcolor,fvextra,caption,fancyhdr,titlesec,enumitem,array,placeins}
\usepackage{xurl}
\usepackage[colorlinks,linkcolor=teal!60!black,urlcolor=teal!60!black,citecolor=teal!60!black]{hyperref}
\captionsetup{font=small,labelfont=bf}
\renewcommand{\tabularxcolumn}[1]{>{\raggedright\arraybackslash}p{#1}}
\setlength{\emergencystretch}{3em}\setlist{itemsep=2pt,topsep=3pt}
\titleformat{\section}{\sffamily\Large\bfseries\color{teal!50!black}}{\thesection}{0.8em}{}
\titleformat{\subsection}{\sffamily\large\bfseries}{\thesubsection}{0.7em}{}
\titleformat{\paragraph}[runin]{\sffamily\bfseries}{}{}{}[.]
\pagestyle{fancy}\fancyhf{}\fancyhead[L]{\sffamily\footnotesize Kitchen food twins: semantic digital twins of food for robot manipulation}
\fancyhead[R]{\sffamily\footnotesize\nouppercase{\leftmark}}\fancyfoot[C]{\sffamily\footnotesize\thepage}
\renewcommand{\headrulewidth}{0.3pt}
\renewcommand{\floatpagefraction}{0.7}\renewcommand{\topfraction}{0.9}\renewcommand{\textfraction}{0.1}
\begin{document}
\begin{titlepage}\sffamily
\noindent\includegraphics[width=\linewidth,height=0.42\textheight,keepaspectratio]{/home/mbeetz/Documents/HouseholdTwin/report/figures/renders_food_hero.jpg}\par\vspace{0.8cm}
\noindent{\color{teal!50!black}\rule{\linewidth}{1.2pt}}\par\vspace{0.4cm}
{\Huge\bfseries Food that robots can reason about\par}\vspace{0.3cm}
{\Large Semantic digital twins of fruit, vegetables, eggs and bread\\ for robot manipulation in household kitchens: what we have so far\par}\vspace{0.3cm}
\noindent{\color{teal!50!black}\rule{\linewidth}{1.2pt}}\par\vspace{0.6cm}
{\large 16 kitchen items for the EASE and GARMI apartments, modelled part by part\par}\vspace{0.3cm}
{\normalsize [Authors and affiliations] \hfill 20 October 2026\par}\vfill
{\small\rmfamily This is a research twin built from published biology, food-science and cooking sources. It is not a
measured model of any particular piece of food. Every quantitative value carries a status label from the project's
parameter register (SOURCED, SOURCED-VERIFY, STANDARD, DERIVED, TYPICAL, CALIBRATED, ILLUSTRATIVE;
Section~\ref{sec:quality} and Appendix~\ref{app:register}). ILLUSTRATIVE values are plausible assumptions of the twin
and must not be quoted as facts. Bench results are simulations, not experiments.\par}
\end{titlepage}
\tableofcontents\clearpage
\section*{Summary}\addcontentsline{toc}{section}{Summary}
\paragraph{Purpose} We want models of everyday food that a robot can plan with, act on and learn about: geometric and
physical objects in Isaac Sim whose every part is a node of a semantic digital twin (SDT), that knows what it is
(\emph{shell}, \emph{yolk}, \emph{seed}, \emph{crust}), how the robot can come to know it (visible, only after
cutting, tactile, inferred, hidden), and how it behaves when it is gripped, dropped, squeezed, cut, peeled or cooked.
The target scenes are the EASE (IAI Bremen) and GARMI apartments, which already exist as robot-simulation scenes in
our CRAM tooling. Each item has three layers: a continuous model of a single, randomly drawn \emph{instance} (its
hidden truth), the robot's \emph{beliefs} about it (priors from a parameter register, updated by observations), and a
small \emph{discrete} model (attributes, states and actions) that plans and action designators refer to.
\paragraph{What was built} 16 items have a code module and a parameter register: the egg set (egg, apple, avocado), wave~1 (roll, banana, carrot), wave~2 (loaf, onion, orange, potato, tomato) and wave~3 (lemon, cucumber, croissant, baguette, pepper). All 16 are built as USD assets: 27 instances in \texttt{assets/food} with 712 part prims in total, each with cut states, a belief layer
(\texttt{sdt.usda}) and a hidden-truth layer (\texttt{truth.usda}). The register holds 564 entries and 205 sources, 166 of which the entries cite by key. There are 30 task designators (egg 7, roll 4, banana 3, carrot 3, loaf 2, tomato 2, potato 2, onion 2, orange 2, avocado 3), seven
Isaac Sim benches with results in \texttt{test\_reports/} (egg grasp and drop, egg opening, avocado, roll and tomato
squeeze, banana peel, carrot plank) and two CPU studies (egg cracking and boiling), and the EASE and GARMI apartments
and the GARMI robot are imported into USD.
\paragraph{One avocado, end to end} The avocado now goes through the whole loop a robot needs
(Section~\ref{sec:avocado-reasoning}): a photo becomes an SDT node; the node is a 4000-particle belief updated by
Bayes' rule from photos, squeezes, cuts, twists and the spoon; a vision-based monitor judges peeling against the
twin's expectation; an emulator halves and pits the avocado as USD states; the robot answers competence questions
from its belief, its task models and mental simulation (including whether a squeeze is worth it); and a PR2 in the
CRAM kitchen detects, squeezes, cuts round the seed, twists the halves apart and spoons out the seed, checking every
step against what it imagined -- scripted, or driven by an operator from the CRAMERA viewer.

\paragraph{From reasoning to acting and to theories} The robot then acts as an agent that enforces the effects its
mission requires and prevents the ones it forbids, deciding over the fruit's lifetime (ripening, spoilage, browning),
and its knowledge is written down and tested as an avocado manipulation microtheory -- the egg microtheory transferred
-- which the robot also learns from its own failures and which programs its reinforcement learning
(Section~\ref{sec:avocado-theory}; the full microtheory is \texttt{report/avocado\_microtheory.pdf}).
\paragraph{Headline results} (1) The egg boiling model (heat conduction through the twin's own tissue map,
doneness thresholds CALIBRATED) reproduces 6/6 doneness classes of the
hot-start timing table it was calibrated on and 3/3 of the fridge-egg
timings; this is a fit, not an independent test. (2) The egg crack study turns kitchen advice into numbers: over 400
register-sampled eggs the best strike on another egg is 0.50~m/s and on a flat
counter 0.45~m/s (strike-model constants ILLUSTRATIVE). (3) For the red tomato the PhysX FEM squeeze at 4000~Hz gives
4.1~N/mm against 4.53~N/mm from Hertz
theory and 4.5~N/mm measured (ince16,
SOURCED-VERIFY); stiffer lanes stay under-converged (green: 6.0 vs Hertz
10.2~N/mm). The same solver limit appears in the avocado and roll benches.
(4) Opening a cracked egg over a bowl in Isaac Sim keeps 99.9~\% of 7,703 albumen particles and the yolk in the bowl.
(5) The carrot bench shows in simulation why cooks cut a plank before slicing: with a slicing stroke the round root
rolls away at 1~N already, the planked half only at about 10~N. (6) The avocado's peeling monitor finds 22/24 correct verdicts (skin IoU 0.46) on synthetic
peels from RGB-D, and in the PR2 scenario the robot's ripeness belief, misled by the photo, is corrected by an
unexpected twist outcome towards the hidden truth (Section~\ref{sec:pr2}).
\paragraph{What it is not, yet} 224 of 564 register entries (40~\%) are ILLUSTRATIVE and another 108 need their source checked. Wave~3 is now built (USD assets with cut states, renders, anatomy plates) but has no task designators or
benches, and its research notes are referenced but missing. Only egg,
apple and avocado have been placed in a kitchen,
and the end-to-end GARMI egg task stopped with a numerical failure of the arm controller. The anatomy plates are
waiting for the user's review, which is a gate of the build plan. The physics benches are single runs of simplified
bodies (homogeneous FEM, rigid peel segments, kinematic shells).
\begin{figure}[htbp]\centering\includegraphics[width=\linewidth,height=0.32\textheight,keepaspectratio]{/home/mbeetz/Documents/HouseholdTwin/report/figures/renders_food_wave1_hero.jpg}\caption{Wave~1 items, path-traced in Isaac Sim~6.2 (existing render): the roll whole and as top and bottom halves, the banana whole and as a lengthwise half, and the carrot whole, as a lengthwise half and as a crown piece.}\end{figure}
\begin{table}[htbp]\centering\scriptsize\caption{What exists per item. Reg.: register entries; \% ILL.: share of ILLUSTRATIVE entries; USD: built instances in assets/food; Des.: task designators. Counted from the repository by this script.}
\begin{tabularx}{\linewidth}{@{}p{0.17\linewidth}lrrrp{0.17\linewidth}p{0.1\linewidth}rX@{}}\toprule Item & Wave & Reg. & \% ILL. & USD & Cut / derived states & Images & Des. & Bench \\\midrule
Hen egg & egg set & 36 & 33 & 8 & cutaway & plate + renders & 7 & grasp/drop, open, crack, cook \\
Apple (Gala-like) & egg set & 33 & 24 & 1 & half\_a, half\_b & plate + renders & 0 & -- \\
Hass avocado & egg set & 48 & 46 & 2 & half\_a, half\_b & plate + renders & 3 & FEM squeeze, peel monitor, emulator, PR2 (CRAM) \\
Wheat roll (Weizenbrötchen) & wave 1 & 46 & 41 & 1 & half\_a, half\_b, top\_half, bottom\_half & plate + renders & 4 & FEM squeeze \\
Cavendish banana & wave 1 & 41 & 37 & 4 & half\_a, half\_b, tip\_half, stalk\_half & plate + renders & 3 & peel \\
Nantes carrot & wave 1 & 42 & 38 & 1 & half\_a, half\_b, tip\_piece, crown\_piece & plate + renders & 3 & plank \\
Sandwich loaf (Toastbrot) & wave 2 & 26 & 42 & 1 & half\_a, half\_b, slices & plate + renders & 2 & -- \\
Yellow onion & wave 2 & 32 & 47 & 1 & half\_a, half\_b, top\_piece, root\_piece & plate + renders & 2 & -- \\
Sweet orange & wave 2 & 38 & 32 & 1 & stem\_half, navel\_half, half\_a, half\_b & plate + renders & 2 & -- \\
Potato & wave 2 & 36 & 39 & 1 & half\_a, half\_b, rose\_piece, heel\_piece & plate + renders & 2 & -- \\
Round / truss tomato & wave 2 & 37 & 46 & 1 & stem\_half, blossom\_half, half\_a, half\_b & plate + renders & 2 & FEM squeeze \\
Lemon & wave 3 & 31 & 29 & 1 & stem\_half, blossom\_half, half\_a, half\_b & plate + renders & 0 & -- \\
Greenhouse cucumber & wave 3 & 28 & 43 & 1 & half\_a, half\_b, rounds & plate + renders & 0 & -- \\
Butter croissant & wave 3 & 21 & 62 & 1 & top\_half, bottom\_half, half\_a, half\_b & plate + renders & 0 & -- \\
Baguette & wave 3 & 42 & 45 & 1 & half\_a, half\_b, top\_half, bottom\_half, crostini & plate + renders & 0 & -- \\
Bell pepper & wave 3 & 27 & 37 & 1 & half\_a, half\_b, top\_piece, bottom\_piece & plate + renders & 0 & -- \\\bottomrule\end{tabularx}\end{table}
\clearpage\section{Method: a food item as a semantic digital twin}\label{sec:method}
\paragraph{From biology to a part tree} Each item is a Python class (\texttt{food/<item>.py}) that draws one instance
from the register (\texttt{params.sample(kind, seed)}) and builds an \emph{analytic anatomy} for it: meridian curves
of revolution, offsets and lofts (\texttt{food/geometry.py}, \texttt{food/shapes.py}) and a label field that says, for
any point in space, which tissue it belongs to (shell, thick albumen, yolk; flavedo, albedo, juice vesicles; crust,
crumb). The class's \texttt{parts()} returns the anatomical partonomy as a tree of \texttt{Part} records: a USD path,
an ontology class, an observability class, a documentation string with sources, and the geometry source (a label of
the field, an analytic mesh, or a point set for things such as pores, lenticels, gas cells and seeds). The anatomy was
researched in biology and food-science sources; the user rejected two early versions (a wrong apple core, an egg-like
avocado seed), and the anatomy plates are still awaiting review (Section~\ref{sec:next}).

\paragraph{Closed meshes and cuts} \texttt{food/partgraph.py} turns every labelled part into a closed (watertight)
mesh by marching cubes on the label field, and computes volume and mass from the part's density. Cuts are plane
splits of every part mesh (manifold3d \texttt{split\_by\_plane}, nudged off the symmetry axis), so a half is again a
complete part graph with its own mass, centre of mass and inertia. Table~\ref{tab:mass} checks mass conservation.

\paragraph{USD part graph} \texttt{food/usd\_partgraph.py} writes one USD asset per instance. Every prim under
\texttt{/Food/Parts} is an SDT node with \texttt{sdt:represents} (the entity IRI of the physical part),
\texttt{sdt:class}, \texttt{sdt:partOf} and \texttt{sdt:observability}, UsdSemantics labels for perception, and
\texttt{food:mass\_g}, \texttt{food:volume\_cm3}, \texttt{food:density\_kgm3} from its closed mesh. The whole object
is one rigid body whose mass properties are summed over the parts; interior parts are invisible until a cut state
exposes them. Two layers sit over the geometry: \texttt{sdt.usda} holds the robot's beliefs (priors from the
register, observations), and \texttt{truth.usda} holds the instance's hidden parameters, which the dynamics models and
audits read but the robot never does.
\begin{figure}[htbp]\centering\includegraphics[width=\linewidth,height=0.5\textheight,keepaspectratio]{/home/mbeetz/Documents/HouseholdTwin/renders/food/scenegraph_sdt.pdf}\caption{The USD scene graphs of the first three items as SDT part graphs (existing figure from tools/anatomy\_doc.py): prim, sdt:class and part mass; the colour of the dot is the observability class.}\end{figure}
\par\medskip\noindent{\small\sffamily\bfseries Excerpt of assets/food/egg\_1/whole.usdc (read with pxr by this script)}\par\nopagebreak\vspace{2pt}\begin{Verbatim}[fontsize=\scriptsize,frame=lines,framesep=2mm,breaklines]
def Xform "Food"   # /Food
    string food:kind = 'egg'
    double food:mass_g = 65.59
    string food:shell_condition = 'intact'
    string sdt:class = 'food:Egg'
    string sdt:represents = 'https://aicor.knowledge/household/food/egg_1'
    token[] semantics:labels:class = [egg]
    token[] semantics:labels:part = [egg]
def Xform "Shell"   # /Food/Parts/Shell
    string sdt:class = 'food:Eggshell'
    string sdt:observability = 'visible'
    string sdt:partOf = 'https://aicor.knowledge/household/food/egg_1'
    string sdt:represents = 'https://aicor.knowledge/household/food/egg_1/shell'
    string sdt:usdPrimPath = '/Food/Parts/Shell'
    token[] semantics:labels:class = [egg_shell]
    token[] semantics:labels:part = [shell]
def Xform "Cuticle"   # /Food/Parts/Shell/Cuticle
    double food:density_kgm3 = 1300
    double food:mass_g = 0.09072
    double food:volume_cm3 = 0.06978
    string sdt:class = 'food:EggCuticle'
    string sdt:observability = 'visible'
    string sdt:partOf = 'https://aicor.knowledge/household/food/egg_1/shell'
    string sdt:represents = 'https://aicor.knowledge/household/food/egg_1/cuticle'
    string sdt:usdPrimPath = '/Food/Parts/Shell/Cuticle'
    token[] semantics:labels:class = [egg_cuticle]
    token[] semantics:labels:part = [cuticle]
def Xform "GerminalDisc"   # /Food/Parts/Yolk/GerminalDisc
    double food:density_kgm3 = 1030
    double food:mass_g = 0.00096
    double food:volume_cm3 = 0.000932
    string sdt:class = 'food:GerminalDisc'
    string sdt:observability = 'visible after opening'
    string sdt:partOf = 'https://aicor.knowledge/household/food/egg_1/yolk'
    string sdt:represents = 'https://aicor.knowledge/household/food/egg_1/germinal_disc'
    string sdt:usdPrimPath = '/Food/Parts/Yolk/GerminalDisc'
    token[] semantics:labels:class = [egg_germinal_disc]
    token[] semantics:labels:part = [germinal_disc]
\end{Verbatim}
\par\medskip\noindent{\small\sffamily\bfseries assets/food/egg\_1/sdt.usda (beliefs, lines 9--28)}\par\nopagebreak\vspace{2pt}\begin{Verbatim}[fontsize=\scriptsize,frame=lines,framesep=2mm,breaklines]
over "Food"
{
    custom string sdt:belief:contents = "in_shell 1.0"
    custom double sdt:belief:mass_g = 65.58658139602177
    custom string sdt:belief:mass_g_source = "weighed"
    custom string sdt:belief:shell_condition = "intact 0.97 | microcracked 0.03 (prior)"
    custom string sdt:epistemicStatus = "prior + observation"

    over "Parts"
    {
        over "Shell"
        {
            custom double sdt:belief:F_break_equatorial_N_mean = 34.92249125484492
            custom double sdt:belief:F_break_equatorial_N_p05 = 21.952609816956432
            custom double sdt:belief:F_break_equatorial_N_p95 = 48.02211933416059
            custom string sdt:epistemicStatus = "prior (register: SOURCED)"
        }

        over "AirCell"
        {
\end{Verbatim}
\par\medskip\noindent{\small\sffamily\bfseries assets/food/egg\_1/truth.usda (hidden truth, first lines)}\par\nopagebreak\vspace{2pt}\begin{Verbatim}[fontsize=\scriptsize,frame=lines,framesep=2mm,breaklines]
#usda 1.0
(
    "Simulator ground truth (hidden instance parameters): overs only, never part of the SDT; used by the dynamics models and audits."
)

over "Food"
{
    custom double truth:age_days = 11.762866264385565
    custom double truth:air_cell_height = 3.198358086590051
    custom double truth:albumen_density = 1035
    custom double truth:asymmetry_w = 1.859487302895341
    custom string truth:colour = "brown"
    custom double truth:cuticle_thickness = 9.208033800009677
    custom double truth:F_break_axial = 35.533810468062526
    custom double truth:F_break_equatorial = 29.108367303986665
    custom double truth:friction_shell_rubber = 0.5409584487874936
\end{Verbatim}
\paragraph{The parameter register} \texttt{food/params.py} holds every number the models use, per item, as
\emph{(unit, value or distribution, status, source)}. The status rule is carried over from the lab twin: no invented
physical parameter is presented as a fact. Distributions (normal with truncation, uniform, integer, choice, lognormal)
are what the hidden-world sampler draws per instance; the robot's belief starts from the same distribution as its
prior and learns the instance's values only through observations. Of the 564 entries, 441 carry a number or a distribution; the rest are model notes (formulas, tables, ranges as
text) that the code implements.
\begin{table}[htbp]\centering\footnotesize\caption{Status labels of the parameter register.}
\begin{tabularx}{\linewidth}{@{}lX@{}}\toprule Status & Meaning (food/params.py) \\\midrule
SOURCED & value or range taken from the cited source \\
SOURCED-VERIFY & value from a source whose details (species, cultivar, protocol) still need checking in the full text \\
STANDARD & a regulation or grading standard (definitional, not a measurement) \\
DERIVED & computed from other register entries by a stated physical model \\
TYPICAL & common textbook value, no citation held in this repository yet \\
CALIBRATED & free model constant fitted to a cited observation (the fit and its residual are stated) \\
ILLUSTRATIVE & twin assumption chosen to be plausible; must be calibrated before any quantitative claim \\\bottomrule\end{tabularx}\end{table}
\paragraph{Discrete schema, dynamics and beliefs} \texttt{food/schema.py} is the symbolic layer: per item, a
class, attributes with domains and observability, and actions with preconditions, effects and \emph{guards} -- the
continuous quantities, computed by the twin, that decide which effect occurs (for the egg: the grip corridor
$F_\mathrm{slip} < F < F_\mathrm{break}$ and the AICOR H8 admissibility $P(\mathrm{slip}) < \varepsilon_s$,
$P(\mathrm{rupture}) < \varepsilon_r$, $P(\mathrm{stable}) > \gamma$). The schema covers the first six items; the
wave~2 and wave~3 items expose their discrete state through \texttt{discrete()} methods of their truth classes
instead. The hidden dynamics are in \texttt{food/dynamics.py} (egg, apple, avocado) and \texttt{food/<item>\_models.py}
(the later items); the beliefs in \texttt{food/sdt.py} and the \texttt{*\_beliefs} functions. Task designators
(\texttt{food/egg\_tasks.py}, \texttt{wave1\_tasks.py}, \texttt{wave2\_tasks.py}) list goal, parameters, steps,
preconditions, effects and failure modes with cues and recoveries, each bound to the hidden model that executes it.
\begin{table}[htbp]\centering\scriptsize\caption{The discrete schema (food/schema.py).}
\begin{tabularx}{\linewidth}{@{}p{0.11\linewidth}Xp{0.3\linewidth}@{}}\toprule Class & Attributes & Actions \\\midrule
Egg & shell\_condition, contents, yolk, weight\_grade, quality\_class, shell\_colour, orientation, support, F\_break, F\_slip & Grasp, Lift, Transport, Place, Drop, CrackOnEdge, OpenUp, Candle \\
Apple & integrity, bruise, peel, core, cut\_face\_browning, firmness, orientation, seed\_positions & Grasp, Place, Drop, Halve, Core, Peel \\
Avocado & ripeness\_stage, integrity, seed\_coat, seed, skin, flesh\_condition, cut\_face\_browning & SqueezeTest, Grasp, Halve, CutRoundSeed, TwistApart, RemoveSeed, Scoop \\
BreadRoll & freshness, crust\_texture, cut\_state, toast\_level, temperature, topping, support & Grasp, SliceOpen, Reheat, ToastHalves, Store \\
Banana & ripeness\_stage, edible, peel\_state, neck, attached\_to\_hand, bruised, chilling\_injury, pulp\_browning, integrity & SeparateFromHand, SnapNeck, PeelStrip, CutRounds, Mash \\
Carrot & firmness\_class, peel\_state, trim\_state, cut\_shape, cooked\_state, surface\_blush, green\_shoulder & Scrub, Peel, TopAndTail, Plank, Cut, Snap, Boil \\\bottomrule\end{tabularx}\end{table}
\paragraph{Mass conservation} The table checks, for every built instance, that the part masses add up to the
whole and that the pieces of each cut add up to the whole (ratio 1.0000 = conserved).
{\scriptsize\begin{longtable}{@{}lrrrp{0.26\linewidth}p{0.17\linewidth}@{}}\caption{Mass checks of the built instances (assets/food/*/instance.json).\label{tab:mass}}\\\toprule Instance & whole (g) & parts (g) & diff. \% & cut pieces / whole & cut failures \\\midrule\endfirsthead
\caption[]{Mass checks of the built instances (assets/food/*/instance.json). (continued)}\\\toprule Instance & whole (g) & parts (g) & diff. \% & cut pieces / whole & cut failures \\\midrule\endhead
\texttt{apple\_\allowbreak{}3} & 217.7 & 217.7 & +0.0 & half\_a+half\_b: 1.0000 & Core/Carpel\_1/Locule, Core/Carpel\_3/Locule \\
\texttt{avocado\_\allowbreak{}4} & 197.5 & 197.5 & +0.0 & half\_a+half\_b: 1.0000 & none \\
\texttt{avocado\_\allowbreak{}photo\_\allowbreak{}burlap} & 213.0 & 213.0 & -0.0 & half\_a+half\_b: 1.0000 & none \\
\texttt{baguette\_\allowbreak{}0} & 213.6 & 216.0 & +1.1 & half\_a+half\_b: 1.0000; top\_half+bottom\_half: 1.0000; crostini: 1.0000 & none \\
\texttt{banana\_\allowbreak{}0} & 117.1 & 117.1 & +0.0 & half\_a+half\_b: 1.0000; tip\_half+stalk\_half: 1.0000 & none \\
\texttt{banana\_\allowbreak{}1} & 130.5 & 130.5 & -0.0 & half\_a+half\_b: 1.0000; tip\_half+stalk\_half: 1.0000 & none \\
\texttt{banana\_\allowbreak{}2} & 119.1 & 119.1 & +0.0 & half\_a+half\_b: 1.0000; tip\_half+stalk\_half: 1.0000 & none \\
\texttt{banana\_\allowbreak{}3} & 113.9 & 113.9 & +0.0 & half\_a+half\_b: 1.0000; tip\_half+stalk\_half: 1.0000 & none \\
\texttt{carrot\_\allowbreak{}0} & 115.0 & 115.0 & -0.0 & half\_a+half\_b: 1.0000; tip\_piece+crown\_piece: 1.0000 & none \\
\texttt{croissant\_\allowbreak{}0} & 51.3 & 51.3 & +0.0 & top\_half+bottom\_half: 1.0000; half\_a+half\_b: 1.0000 & none \\
\texttt{cucumber\_\allowbreak{}0} & 371.0 & 371.0 & +0.0 & half\_a+half\_b: 1.0000; rounds: 1.0000 & none \\
\texttt{egg\_\allowbreak{}0} & 62.0 & 62.0 & +0.0 & -- & none \\
\texttt{egg\_\allowbreak{}1} & 65.6 & 65.6 & +0.0 & -- & none \\
\texttt{egg\_\allowbreak{}2} & 61.3 & 61.3 & +0.0 & -- & none \\
\texttt{egg\_\allowbreak{}3} & 64.3 & 64.3 & +0.0 & -- & none \\
\texttt{egg\_\allowbreak{}4} & 57.1 & 57.1 & +0.0 & -- & none \\
\texttt{egg\_\allowbreak{}5} & 53.4 & 53.4 & +0.0 & -- & none \\
\texttt{egg\_\allowbreak{}6} & 72.9 & 72.9 & +0.0 & -- & none \\
\texttt{egg\_\allowbreak{}7} & 62.2 & 62.2 & +0.0 & -- & none \\
\texttt{lemon\_\allowbreak{}0} & 118.0 & 118.0 & +0.0 & stem\_half+blossom\_half: 1.0000; half\_a+half\_b: 1.0000 & none \\
\texttt{loaf\_\allowbreak{}0} & 494.6 & 501.5 & +1.4 & half\_a+half\_b: 1.0000; slices: 1.0000 & none \\
\texttt{onion\_\allowbreak{}0} & 146.5 & 146.5 & -0.0 & half\_a+half\_b: 1.0001; top\_piece+root\_piece: 1.0000 & none \\
\texttt{orange\_\allowbreak{}0} & 193.6 & 193.6 & +0.0 & stem\_half+navel\_half: 1.0000; half\_a+half\_b: 1.0000 & none \\
\texttt{pepper\_\allowbreak{}0} & 163.6 & 163.6 & +0.0 & half\_a+half\_b: 1.0000; top\_piece+bottom\_piece: 1.0000 & none \\
\texttt{potato\_\allowbreak{}0} & 126.4 & 126.4 & +0.0 & half\_a+half\_b: 1.0000; rose\_piece+heel\_piece: 1.0000 & none \\
\texttt{roll\_\allowbreak{}0} & 47.4 & 48.1 & +1.4 & half\_a+half\_b: 1.0000; top\_half+bottom\_half: 1.0000 & none \\
\texttt{tomato\_\allowbreak{}15} & 71.7 & 71.7 & +0.0 & stem\_half+blossom\_half: 1.0000; half\_a+half\_b: 1.0000 & none \\\bottomrule\end{longtable}}
Part sums deviate from the whole for \texttt{baguette\_\allowbreak{}0} (216.0 vs 213.6~g, +1.1~\%), \texttt{loaf\_\allowbreak{}0} (501.5 vs 494.6~g, +1.4~\%), \texttt{roll\_\allowbreak{}0} (48.1 vs 47.4~g, +1.4~\%): these are the two breads, whose crust and crumb meshes appear to overlap or are meshed at a different
resolution from the whole; the wave~1 checklist in \texttt{docs/wave1/README.md} reports this gate as passed, which
this check does not confirm. The apple's lengthwise cut lost two empty locules (cut failures above); their mass is
negligible.
\clearpage\section{The hen egg}\label{sec:egg}
The egg is the most developed item: anatomy, crack and cook models, a belief that computes the AICOR H8 probabilities,
seven task designators, a grasp-and-drop bench and an opening simulation in Isaac Sim. Eight instances are built
(\texttt{egg\_0}..\texttt{egg\_7}, 53.4--72.9~g); the
examples use \texttt{egg\_1} (65.6~g).
\begin{figure}[htbp]\centering\includegraphics[width=\linewidth,height=0.3\textheight,keepaspectratio]{/home/mbeetz/Documents/HouseholdTwin/report/figures/renders_food_anatomy_egg_labelled.jpg}\caption{Exploded anatomy plate of egg\_1 (existing render): outside to inside, every part is a prim of the USD part graph with its mass.}\end{figure}
\begin{figure}[htbp]\centering\begin{minipage}[t]{0.327\linewidth}\centering\includegraphics[width=\linewidth,height=0.2\textheight,keepaspectratio]{/home/mbeetz/Documents/HouseholdTwin/report/figures/renders_food_egg_cutaway.jpg}\\[-1pt]{\scriptsize cutaway state}\end{minipage}\hfill\begin{minipage}[t]{0.327\linewidth}\centering\includegraphics[width=\linewidth,height=0.2\textheight,keepaspectratio]{/home/mbeetz/Documents/HouseholdTwin/report/figures/renders_food_egg_opened_scene.jpg}\\[-1pt]{\scriptsize opened, fried}\end{minipage}\hfill\begin{minipage}[t]{0.327\linewidth}\centering\includegraphics[width=\linewidth,height=0.2\textheight,keepaspectratio]{/home/mbeetz/Documents/HouseholdTwin/report/figures/renders_food_egg_shell_fragments.jpg}\\[-1pt]{\scriptsize half shell}\end{minipage}\caption{Egg renders (existing, path-traced). The cutaway shows the yolk layers, latebra and chalazae.}\end{figure}
\begin{table}[htbp]\centering\scriptsize\caption{Part tree of egg\_1 (instance.json and whole.usdc).}
\begin{tabularx}{\linewidth}{@{}p{0.30\linewidth}p{0.27\linewidth}p{0.17\linewidth}rX@{}}\toprule Part (below /Food/Parts) & sdt:class & observability & n & mass (g) \\\midrule
\hspace*{0.0em}Shell & \texttt{food:Eggshell} & visible &  &  \\
\hspace*{0.9em}Cuticle & \texttt{food:EggCuticle} & visible &  & 0.09 \\
\hspace*{0.9em}PalisadeLayer & \texttt{food:PalisadeLayer} & hidden &  & 3.68 \\
\hspace*{0.9em}MammillaryLayer & \texttt{food:MammillaryLayer} & hidden &  & 1.52 \\
\hspace*{0.9em}Pores & \texttt{food:EggshellPore} & inferred &  & 12,731 points \\
\hspace*{0.0em}ShellMembranes & \texttt{food:ShellMembrane} & hidden &  &  \\
\hspace*{0.9em}OuterShellMembrane & \texttt{food:OuterShellMembrane} & hidden &  & 0.45 \\
\hspace*{0.9em}InnerShellMembrane & \texttt{food:InnerShellMembrane} & hidden &  & 0.15 \\
\hspace*{0.0em}AirCell & \texttt{food:AirCell} & inferred &  & 0.00 \\
\hspace*{0.0em}Albumen & \texttt{food:Albumen} & visible after opening &  &  \\
\hspace*{0.9em}OuterThinAlbumen & \texttt{food:OuterThinAlbumen} & hidden &  & 9.52 \\
\hspace*{0.9em}ThickAlbumen & \texttt{food:ThickAlbumen} & hidden &  & 22.01 \\
\hspace*{0.9em}InnerThinAlbumen & \texttt{food:InnerThinAlbumen} & hidden &  & 6.68 \\
\hspace*{0.9em}ChalaziferousLayer & \texttt{food:ChalaziferousLayer} & hidden &  & 1.18 \\
\hspace*{0.9em}Chalazae & \texttt{food:Chalazae} & visible after opening &  &  \\
\hspace*{1.8em}BluntEnd & \texttt{food:Chalaza} & visible after opening &  & 0.01 \\
\hspace*{1.8em}PointedEnd & \texttt{food:Chalaza} & visible after opening &  & 0.01 \\
\hspace*{0.0em}Yolk & \texttt{food:Yolk} & visible after opening &  &  \\
\hspace*{0.9em}VitellineMembrane & \texttt{food:VitellineMembrane} & hidden &  & 0.03 \\
\hspace*{0.9em}YolkBody & \texttt{food:YolkBody} & visible after opening &  & 20.07 \\
\hspace*{0.9em}Latebra & \texttt{food:Latebra} & hidden &  & 0.16 \\
\hspace*{0.9em}LatebraNeck & \texttt{food:LatebraNeck} & hidden &  & 0.03 \\
\hspace*{0.9em}NucleusOfPander & \texttt{food:NucleusOfPander} & hidden &  & 0.01 \\
\hspace*{0.9em}GerminalDisc & \texttt{food:GerminalDisc} & visible after opening &  & 0.00 \\\bottomrule\end{tabularx}\end{table}
\begin{table}[htbp]\centering\scriptsize\caption{Key register entries of the egg.}
\begin{tabularx}{\linewidth}{@{}p{0.19\linewidth}p{0.27\linewidth}p{0.12\linewidth}X@{}}\toprule Register key & Value / distribution & Status & Source and note \\\midrule
\texttt{length\_\allowbreak{}L} & N(57, 1.8) in [52, 63] mm & SOURCED & kop: 58.05 mm; turjaf: 56.12 mm \\
\texttt{shape\_\allowbreak{}index\_\allowbreak{}B\_\allowbreak{}over\_\allowbreak{}L} & N(0.765, 0.02) in [0.7, 0.81] & SOURCED & kop: 78.25 \%; alt: normal 72-76 \\
\texttt{shell\_\allowbreak{}thickness} & N(0.35, 0.03) in [0.28, 0.42] mm & SOURCED & kop: 0.28-0.41 mm; sb: 0.35 mm \\
\texttt{shell\_\allowbreak{}E} & U(15 .. 43) GPa & SOURCED & sev: 41-43 GPa nanoindentation; sb model 15 GPa \\
\texttt{F\_\allowbreak{}break\_\allowbreak{}equatorial} & N(35, 8) in [12, 59] N & SOURCED & kop: Hisex 36.43 N short axis, 10-59 N; icf11 34$\pm$6 N \\
\texttt{F\_\allowbreak{}break\_\allowbreak{}axial} & N(37, 9) in [12, 62] N & SOURCED & kop 34.92 N long axis; icf11 41$\pm$8 N; alt: highest along the polar axis at the sharp end \\
\texttt{shell\_\allowbreak{}contact\_\allowbreak{}stiffness} & Reissner k = 4 E t\^{}2 / (R\_local sqrt(3(1-nu\^{}2))) N/mm & DERIVED & reissner1946 + shell\_E + shell\_thickness \\
\texttt{yolk\_\allowbreak{}mass\_\allowbreak{}fraction} & N(0.3, 0.015) in [0.27, 0.32] & SOURCED & br: 27-32 \%; usda 31 \% \\
\texttt{air\_\allowbreak{}cell\_\allowbreak{}height} & 1.5 + k(T) * age; k(6 C) = 0.08, k(22 C) = 0.19 mm/d mm & SOURCED-VERIFY & gr: < 4 mm after 28 d at 6 C; > 6 mm after 25 d at 22 C; graded vs eu589/usda \\
\texttt{microcrack\_\allowbreak{}band} & 0.85 & ILLUSTRATIVE & F/F\_break in [0.85, 1) may leave an invisible micro-crack \\
\texttt{impact\_\allowbreak{}effective\_\allowbreak{}mass\_\allowbreak{}fraction} & 0.53 & CALIBRATED & c = sqrt(m\_eff)/rate fitted to cohen2025: side-on drops onto a hard surface crack $\sim$45 \% at 10 mm (model 45 \%), vertical 55 \% at 8 mm (model $\sim$57 \%); the model population is more heterogeneous than one MIT batch (8 mm side-on: model $\sim$28 \% vs 5 \%) \\
\texttt{impact\_\allowbreak{}rate\_\allowbreak{}factor} & 1.3 & ILLUSTRATIVE & dynamic/static crack force; strength rises with loading rate (only sqrt(m\_eff)/rate is identified by the calibration) \\
\texttt{crack\_\allowbreak{}strike\_\allowbreak{}model} & food/egg\_crack.py P: coupled mass 0.05 kg; membrane tears at F/Fc 1.8 (flat), 1.05 (rim), 3.0 (egg-on-egg); shatter at 2.6 / 2.2 / 2.4; shards 0.15 / 0.6 / 0 per F/Fc; yolk puncture 0.06 per shard & ILLUSTRATIVE & structure from atk\_egg, egginfo\_crack, atk\_egg\_on\_egg; numbers to calibrate \\
\texttt{cook\_\allowbreak{}yolk\_\allowbreak{}thresholds} & f65 < 0.15 runny, < 0.45 jammy, < 0.85 creamy, else firm & CALIBRATED & alphafoodie timing table (hot start, medium eggs) + Lopez-Alt fridge-egg timings \\\bottomrule\end{tabularx}\end{table}
\paragraph{States and hidden dynamics} The schema's attributes are shell\_condition, contents, yolk, weight\_grade, quality\_class, shell\_colour, orientation, support, F\_break, F\_slip. Shell condition follows the USDA vocabulary (intact, microcracked, cracked, shattered, opened); a
micro-crack is invisible and lowers the break force by the register's \texttt{microcrack\_strength\_loss}.
\texttt{EggTruth} (\texttt{food/dynamics.py}) cracks the shell when the effective contact force exceeds the
direction-dependent break force; impacts use $F = v\sqrt{k_\mathrm{eff}\, m_\mathrm{eff}}$ with the Reissner shell
stiffness in series with the surface. \texttt{food/egg\_crack.py} adds strike outcomes on a counter, a bowl rim or
another egg and the opening over a vessel; its constants are ILLUSTRATIVE. \texttt{food/egg\_cook.py} solves
axisymmetric heat conduction through the egg's own tissue map. The belief \texttt{EggBelief} (\texttt{food/sdt.py})
holds particles over crack force, shell thickness and friction and returns the grip corridor and H8 admissibility.
\begin{table}[htbp]\centering\scriptsize\caption{Egg task designators (food/egg\_tasks.py; rendered in docs/EGG\_TASKS.md).}
\begin{tabularx}{\linewidth}{@{}p{0.15\linewidth}Xp{0.36\linewidth}@{}}\toprule Designator & Goal & Twin model \\\midrule
PickUpEgg & egg held, shell\_condition unchanged & PhysX contacts + EggTruth.static\_load; EggBelief.admissible (AICOR H8) \\
TransportEgg & egg at target support, intact & EggTruth.impact (calibrated to cohen2025) \\
CrackEgg & shell cracked open, contents in the cup/bowl, yolk whole, no shell inside & food/egg\_crack.py strike + open\_up + fish\_out \\
SeparateEgg & yolk in yolk bowl (whole), white in white bowl with no yolk & EC.open\_up yolk-break model (age, drop height, shards); separation success ILLUSTRATIVE \\
BoilEgg & yolk doneness = target (runny | jammy | creamy | firm), shell intact, peels cleanly & food/egg\_cook.py EggHeat (conduction through the tissue map; thresholds CALIBRATED) \\
ScrambleEggs & curds set but glossy, no liquid egg, not browned & not yet simulated (coagulation thresholds from opentextbc) \\
FrySunnySideUp & white set and opaque incl. around the yolk, yolk whole and runny, no shell & not yet simulated \\\bottomrule\end{tabularx}\end{table}
\subsection{Bench: grasp and drop (Isaac Sim)}
\texttt{scripts/sim\_egg\_bench.py} grips egg\_1 with PhysX contacts at commanded pad forces and drops eggs onto four
kitchen surfaces; outcomes come from the hidden \texttt{EggTruth}, predictions from \texttt{EggBelief}. In the grip
sweep, the measured pad force equals the commanded one within 0.40~N; the egg slips at 0.5~N, the H8-admissible commands run from 1 to 15~N, and the shell cracks from 25~N (hidden crack force of this egg 23~N). In a validation lane eight
instances were gripped at 30~N: 6 of 8 cracked or micro-cracked, while the belief predicted a
rupture probability of about 0.34. The eight hidden crack forces (median 29.5~N) sit below the register prior's mean (35~N, SOURCED); with eight eggs this is
not conclusive, but it is worth checking whether the truth sampler and the belief draw crack forces from the same
effective distribution (the truth applies thickness scaling, which the belief may not). Drops from 1~cm already crack this egg on granite, oak board, silicone mat; on a towel it survives up to 8~cm and cracks at 15~cm (critical heights in the figure). The impact
model's effective-mass fraction is CALIBRATED to the drop statistics of Cohen et al.\ 2025 (cohen2025).
\begin{figure}[htbp]\centering\includegraphics[width=\linewidth,height=0.3\textheight,keepaspectratio]{/home/mbeetz/Documents/HouseholdTwin/report/figures/renders_food_egg_bench.jpg}\caption{Egg bench (existing figure, tools/plot\_egg\_bench.py): (a) grip sweep, belief vs simulated outcome; (b) eight instances at 30~N; (c) drops onto kitchen surfaces with peak force and impact speed.}\end{figure}
\subsection{Study: how hard to crack an egg (CPU)}
\texttt{tools/egg\_crack\_study.py} strikes 400 register-sampled eggs at 30 speeds against four tools and counts how
often the egg ends up open in the cup with a whole yolk and no shell -- directly, and with the recovery people use
(crack into a cup first, fish shards out with a half shell). The strike-outcome bands, shard rates and puncture
probabilities are ILLUSTRATIVE (register \texttt{crack\_strike\_model}), so the optimum speeds are model answers, not
measurements. This report re-ran the computation in memory with the current register (impact effective-mass fraction
0.53): the best speeds and probabilities are identical to the stored study.
\begin{table}[htbp]\centering\scriptsize\caption{Egg crack study (test\_reports/egg\_crack\_study.json) and this report's re-run.}
\begin{tabularx}{\linewidth}{@{}Xrrrrr@{}}\toprule Strike surface & best v (m/s) & P with recovery & P direct (best) & 90 \% window (m/s) & re-run v / P \\\midrule
flat counter (granite) & 0.45 & 0.88 & 0.75 & 0.40--0.50 & 0.45 / 0.88 \\
thin bowl rim (r 1.5 mm) & 0.25 & 0.76 & 0.50 & 0.20--0.25 & 0.25 / 0.76 \\
thick bowl rim (r 4 mm) & 0.35 & 0.77 & 0.49 & 0.35--0.40 & 0.35 / 0.77 \\
another egg & 0.50 & 0.85 & 0.80 & 0.45--0.65 & 0.50 / 0.85 \\\bottomrule\end{tabularx}\end{table}
\begin{figure}[htbp]\centering\includegraphics[width=\linewidth,height=0.3\textheight,keepaspectratio]{/home/mbeetz/Documents/HouseholdTwin/report/figures/renders_food_egg_crack_study.jpg}\caption{Crack study (existing figure): outcome shares (top) and task success (bottom) over strike speed per tool.}\end{figure}
\subsection{Study: boiling (CPU)}
\texttt{tools/egg\_cook\_study.py} runs \texttt{EggHeat} on egg\_4 (thermal properties TYPICAL, doneness thresholds
on the yolk fraction that reached 65~\textdegree C CALIBRATED to the alphafoodie timing table and the Lopez-Alt fridge-egg
timings). The table compares the model's doneness class with the source at each source time. Because the thresholds
were fitted to these points, agreement shows a consistent fit, not a validation; no independent timing data are held.
The cold-start case (water heated to the boil in 9~min, then off the heat, ILLUSTRATIVE profile) is not compared with
a source.
\begin{table}[htbp]\centering\scriptsize\caption{Egg boiling model vs the cooking sources it was calibrated on (test\_reports/egg\_cook\_study.json).}
\begin{tabularx}{\linewidth}{@{}rXlll@{}}\toprule min & Case & Source class & Model class & Match \\\midrule
4.5 & hot start, room-temperature egg & runny & runny & yes \\
6 & hot start, room-temperature egg & jammy & jammy & yes \\
7 & hot start, room-temperature egg & jammy & jammy & yes \\
8.5 & hot start, room-temperature egg & creamy & creamy & yes \\
10 & hot start, room-temperature egg & firm & firm & yes \\
11 & hot start, room-temperature egg & firm & firm & yes \\
6 & hot start, fridge egg & runny & runny & yes \\
8.5 & hot start, fridge egg & creamy & creamy & yes \\
11 & hot start, fridge egg & firm & firm & yes \\\bottomrule\end{tabularx}\end{table}
\begin{figure}[htbp]\centering\includegraphics[width=\linewidth,height=0.36\textheight,keepaspectratio]{/home/mbeetz/Documents/HouseholdTwin/report/figures/renders_food_egg_cook_study.jpg}\caption{Boiling study (existing figure): doneness over time with the source markers, yolk-centre temperature, and temperature sections at five times.}\end{figure}
\subsection{Simulation: opening a cracked egg (Isaac Sim)}
\texttt{scripts/sim\_egg\_open.py} replaces the cracked egg by its two shell fragments (kinematic, from the crack
model), the albumen by 7,703 PBD particles and the yolk by a rigid body, and swings the
halves apart over a glass bowl. After 3.95~s, 99.91~\% of the particles are in the bowl, none below the table, and the
yolk is in the bowl. The yolk cannot break in this simulation (rigid body); yolk breakage comes from the hidden
\texttt{open\_up} model instead. The albumen surface is reconstructed from the particles by marching cubes, because the
PhysX isosurface does not work with Fabric.
\begin{figure}[htbp]\centering\includegraphics[width=\linewidth,height=0.33\textheight,keepaspectratio]{/home/mbeetz/Documents/HouseholdTwin/report/figures/renders_food_egg_open_filmstrip.jpg}\caption{Opening simulation (existing filmstrip): halves swing apart, the albumen streams, the yolk drops whole and relaxes to its dome.}\end{figure}
\clearpage\section{Apple and avocado}
\subsection{Apple (Gala-like)}
One instance is built (\texttt{apple\_3}, 217.7~g) with two lengthwise halves.
The anatomy follows the pome: hypanthial flesh (cortex) round a core of five fused carpels, ten primary vascular
bundles, a calyx with five sepals, lenticels as a point set and five named surface regions of the cuticle.
\begin{figure}[htbp]\centering\includegraphics[width=\linewidth,height=0.28\textheight,keepaspectratio]{/home/mbeetz/Documents/HouseholdTwin/report/figures/renders_food_anatomy_apple_labelled.jpg}\caption{Exploded anatomy plate of apple\_3 (existing render).}\end{figure}
\begin{figure}[htbp]\centering\begin{minipage}[t]{0.490\linewidth}\centering\includegraphics[width=\linewidth,height=0.2\textheight,keepaspectratio]{/home/mbeetz/Documents/HouseholdTwin/report/figures/renders_food_apple_halves.jpg}\\[-1pt]{\scriptsize whole and halves}\end{minipage}\hfill\begin{minipage}[t]{0.490\linewidth}\centering\includegraphics[width=\linewidth,height=0.2\textheight,keepaspectratio]{/home/mbeetz/Documents/HouseholdTwin/report/figures/renders_food_closeup_apple_cut.jpg}\\[-1pt]{\scriptsize cut face: core and seeds}\end{minipage}\caption{Apple renders (existing, path-traced).}\end{figure}
\begin{table}[htbp]\centering\scriptsize\caption{Part tree of apple\_3 (indexed siblings merged).}
\begin{tabularx}{\linewidth}{@{}p{0.30\linewidth}p{0.27\linewidth}p{0.17\linewidth}rX@{}}\toprule Part (below /Food/Parts) & sdt:class & observability & n & mass (g) \\\midrule
\hspace*{0.0em}Skin & \texttt{food:AppleSkin} & visible &  &  \\
\hspace*{0.9em}Cuticle & \texttt{food:Cuticle} & visible &  & 0.22 \\
\hspace*{0.9em}Epidermis & \texttt{food:Epidermis} & hidden &  & 0.33 \\
\hspace*{0.9em}Hypodermis & \texttt{food:Hypodermis} & hidden &  & 0.52 \\
\hspace*{0.9em}Lenticels & \texttt{food:Lenticel} & visible &  & 1,717 points \\
\hspace*{0.0em}Stem & \texttt{food:Pedicel} & visible &  & 0.07 \\
\hspace*{0.0em}Calyx & \texttt{food:Calyx} & visible &  &  \\
\hspace*{0.9em}Sepal\_1..5 & \texttt{food:Sepal} & visible & 5 & 0.01 \\
\hspace*{0.9em}CalyxTube & \texttt{food:CalyxTube} & hidden &  & 0.01 \\
\hspace*{0.9em}StamenRemnants & \texttt{food:StamenRemnant} & visible &  & 0.00 \\
\hspace*{0.0em}Cortex & \texttt{food:Cortex} & visible after cutting &  & 203.51 \\
\hspace*{0.0em}VascularSystem & \texttt{food:VascularSystem} & visible after cutting &  &  \\
\hspace*{0.9em}SepalBundle\_1..5 & \texttt{food:SepalBundle} & visible after cutting & 5 & 0.13 \\
\hspace*{0.9em}PetalBundle\_1..5 & \texttt{food:PetalBundle} & visible after cutting & 5 & 0.13 \\
\hspace*{0.0em}Core & \texttt{food:AppleCore} & visible after cutting &  &  \\
\hspace*{0.9em}Pith & \texttt{food:CoreTissue} & visible after cutting &  & 11.64 \\
\hspace*{0.9em}Carpel\_1..5 & \texttt{food:Carpel} & visible after cutting & 5 &  \\
\hspace*{1.8em}Endocarp & \texttt{food:Endocarp} & visible after cutting & 5 & 0.66 \\
\hspace*{1.8em}Locule & \texttt{food:Locule} & visible after cutting & 5 & 0.00 \\
\hspace*{1.8em}Seed\_1..9 & \texttt{food:AppleSeed} & visible after cutting & 9 & 0.47 \\\bottomrule\end{tabularx}\end{table}
\begin{table}[htbp]\centering\scriptsize\caption{Key register entries of the apple.}
\begin{tabularx}{\linewidth}{@{}p{0.19\linewidth}p{0.27\linewidth}p{0.12\linewidth}X@{}}\toprule Register key & Value / distribution & Status & Source and note \\\midrule
\texttt{diameter\_\allowbreak{}D} & N(75, 3.5) in [66, 86] mm & SOURCED & wiki\_gala: 78 \% of Gala fruit 70-80 mm \\
\texttt{flesh\_\allowbreak{}density} & N(860, 40) in [760, 970] kg/m3 & SOURCED & wiki\_gala 0.86 g/cm3; fuji 955-968 kg/m3 (degruyter) \\
\texttt{flesh\_\allowbreak{}E} & lognormal, median 3, factor 1.7 MPa & SOURCED & fuji 5.0 MPa (consensus f4110d16); quasi-static 4.2-7.8; stored 0.5-1.7 (tjaf2005) \\
\texttt{mt\_\allowbreak{}firmness} & N(75, 8) in [50, 95] N & SOURCED & msu2025: Gala $\sim$18 lb (80 N) at harvest \\
\texttt{flesh\_\allowbreak{}shear\_\allowbreak{}strength} & N(0.4, 0.05) in [0.28, 0.5] MPa & SOURCED & fuji shear 0.45 MPa; failure stress 0.32-0.37 MPa (consensus) \\
\texttt{bruise\_\allowbreak{}energy\_\allowbreak{}threshold} & U(0.004 .. 0.012) J & SOURCED-VERIFY & no bruise below 0.25 m/s impact; bruises from 0.011 J pendulum (consensus) \\
\texttt{bruise\_\allowbreak{}volume\_\allowbreak{}per\_\allowbreak{}J} & 0.01 * V\_fruit / 0.41 J cm3/J & SOURCED-VERIFY & fsct2008 (about 1 \% of fruit volume at 410 mJ); linear in energy \\
\texttt{cutting\_\allowbreak{}force\_\allowbreak{}peak} & N(25, 5) in [12, 40] N & ILLUSTRATIVE & no apple cutting force found \\
\texttt{friction\_\allowbreak{}skin\_\allowbreak{}rubber} & U(0.7 .. 0.95) & SOURCED-VERIFY & fuji vs TPU rubber 0.88 (consensus) \\
\texttt{core\_\allowbreak{}line\_\allowbreak{}radius} & N(14, 1.5) in [11, 18] mm & SOURCED & uspto: core diameter 24.7-36.7 mm \\
\texttt{browning\_\allowbreak{}dE} & 2.38 (0 min), 2.68 (30 min), 4.05 (60 min) dE & SOURCED & pmc7022590 (Golden Delicious) \\\bottomrule\end{tabularx}\end{table}
\paragraph{States, dynamics, tasks} Schema attributes: integrity, bruise, peel, core, cut\_face\_browning, firmness, orientation, seed\_positions; actions: Grasp, Place, Drop, Halve, Core, Peel. \texttt{AppleTruth} bruises under a static load when the Hertz
maximum shear $0.31\,p_0$ exceeds the flesh shear strength, and under impact with a bruise volume linear in absorbed
energy above a threshold (fsct2008 anchor, SOURCED-VERIFY); \texttt{AppleBelief} tracks the expected bruise after
known drops and grips. There are \textbf{no apple task designators and no apple bench} yet; the apple's cutting force
is ILLUSTRATIVE (no source found).
\subsection{Hass avocado}
One instance is built (\texttt{avocado\_4}, 197.5~g); the knife goes round the
seed, so \texttt{half\_a} keeps the whole seed and \texttt{half\_b} has the seed cavity. The exocarp has five layers
including the stone-cell layer that lets ripe Hass peel cleanly, the flesh three colour bands and six vascular
strands, and the seed a two-layer testa, two cotyledons and the embryonic axis.
\begin{figure}[htbp]\centering\includegraphics[width=\linewidth,height=0.28\textheight,keepaspectratio]{/home/mbeetz/Documents/HouseholdTwin/report/figures/renders_food_anatomy_avocado_labelled.jpg}\caption{Exploded anatomy plate of avocado\_4 (existing render).}\end{figure}
\begin{figure}[htbp]\centering\begin{minipage}[t]{0.490\linewidth}\centering\includegraphics[width=\linewidth,height=0.2\textheight,keepaspectratio]{/home/mbeetz/Documents/HouseholdTwin/report/figures/renders_food_avocado_halves.jpg}\\[-1pt]{\scriptsize whole and halves}\end{minipage}\hfill\begin{minipage}[t]{0.490\linewidth}\centering\includegraphics[width=\linewidth,height=0.2\textheight,keepaspectratio]{/home/mbeetz/Documents/HouseholdTwin/report/figures/renders_food_closeup_avocado_seed.jpg}\\[-1pt]{\scriptsize half with seed}\end{minipage}\caption{Avocado renders (existing, path-traced).}\end{figure}
\begin{table}[htbp]\centering\scriptsize\caption{Part tree of avocado\_4 (indexed siblings merged).}
\begin{tabularx}{\linewidth}{@{}p{0.30\linewidth}p{0.27\linewidth}p{0.17\linewidth}rX@{}}\toprule Part (below /Food/Parts) & sdt:class & observability & n & mass (g) \\\midrule
\hspace*{0.0em}Exocarp & \texttt{food:Exocarp} & visible &  &  \\
\hspace*{0.9em}Cuticle & \texttt{food:Cuticle} & visible &  & 0.26 \\
\hspace*{0.9em}Epidermis & \texttt{food:Epidermis} & hidden &  & 0.44 \\
\hspace*{0.9em}Hypodermis & \texttt{food:Hypodermis} & hidden &  & 3.82 \\
\hspace*{0.9em}SkinParenchyma & \texttt{food:SkinParenchyma} & hidden &  & 23.02 \\
\hspace*{0.9em}StoneCellLayer & \texttt{food:StoneCellLayer} & hidden &  & 3.34 \\
\hspace*{0.9em}Lenticels & \texttt{food:Lenticel} & visible &  & 499 points \\
\hspace*{0.0em}StemButton & \texttt{food:PedicelRemnant} & visible &  & 0.04 \\
\hspace*{0.0em}Mesocarp & \texttt{food:Mesocarp} & visible after cutting &  &  \\
\hspace*{0.9em}OuterDarkGreenBand & \texttt{food:FleshBand} & visible after cutting &  & 65.66 \\
\hspace*{0.9em}MiddlePaleGreenBand & \texttt{food:FleshBand} & visible after cutting &  & 45.14 \\
\hspace*{0.9em}InnerYellowBand & \texttt{food:FleshBand} & visible after cutting &  & 28.52 \\
\hspace*{0.9em}VascularStrands & \texttt{food:VascularStrands} & visible after cutting &  &  \\
\hspace*{1.8em}Strand\_1..6 & \texttt{food:VascularStrand} & visible after cutting & 6 & 0.15 \\
\hspace*{0.0em}Endocarp & \texttt{food:Endocarp} & visible after cutting &  & 0.65 \\
\hspace*{0.0em}Seed & \texttt{food:AvocadoSeed} & visible after halving &  &  \\
\hspace*{0.9em}SeedCoat & \texttt{food:Testa} & visible after halving &  &  \\
\hspace*{1.8em}OuterTesta & \texttt{food:OuterTesta} & visible after halving &  & 0.67 \\
\hspace*{1.8em}InnerTesta & \texttt{food:InnerTesta} & hidden &  & 0.44 \\
\hspace*{0.9em}Cotyledon\_1..2 & \texttt{food:Cotyledon} & hidden & 2 & 25.32 \\
\hspace*{0.9em}EmbryonicAxis & \texttt{food:EmbryonicAxis} & hidden &  &  \\
\hspace*{1.8em}Plumule & \texttt{food:Plumule} & hidden &  & 0.00 \\
\hspace*{1.8em}Hypocotyl & \texttt{food:Hypocotyl} & hidden &  & 0.00 \\
\hspace*{1.8em}Radicle & \texttt{food:Radicle} & hidden &  & 0.00 \\\bottomrule\end{tabularx}\end{table}
\begin{table}[htbp]\centering\scriptsize\caption{Key register entries of the avocado.}
\begin{tabularx}{\linewidth}{@{}p{0.19\linewidth}p{0.27\linewidth}p{0.12\linewidth}X@{}}\toprule Register key & Value / distribution & Status & Source and note \\\midrule
\texttt{length\_\allowbreak{}L} & N(98, 6) in [80, 120] mm & SOURCED-VERIFY & specialty: 7-15 cm; no Hass mm data held \\
\texttt{diameter\_\allowbreak{}D} & N(66, 3.5) in [56, 76] mm & TYPICAL & chosen so mass $\sim$ 200 g (uspto mean 204 g) \\
\texttt{skin\_\allowbreak{}thickness} & N(1.6, 0.15) in [1.2, 2] mm & SOURCED & uspto\_pp14239: Hass skin 1.6 mm \\
\texttt{seed\_\allowbreak{}length} & N(40, 2.5) in [34, 47] mm & DERIVED & seed mass 28.4-30.6 g (uspto) at 1.1 g/cm3 \\
\texttt{ripeness\_\allowbreak{}stage} & U(1 .. 4.6) 1..5 & STANDARD & avo\_au\_stages; continuous hidden value \\
\texttt{ripening\_\allowbreak{}rate} & N(0.55, 0.1) in [0.3, 0.8] stage/d & SOURCED-VERIFY & ashton2006: firm-ripe $\sim$day 8, soft-ripe $\sim$9.5 at 20 C; landahl: 304 $\to$ 2 N in 6 d at 18 C \\
\texttt{penetrometer\_\allowbreak{}force} & stage 1..5 = 118, 46, 15, 8.3, 4.9 N (geometric band centres), log-linear N & SOURCED & avo\_au\_stages: hard >10, rubbery 2.2-10, breaking 1.1-2.1, firm-ripe 0.7-1.0, ripe 0.4-0.6 kgf (11.1 mm tip, 8 mm) \\
\texttt{flesh\_\allowbreak{}E} & E = 2.29 (F\_pen / 118 N)\^{}0.837 $\to$ 2.29 .. 0.16 MPa MPa & SOURCED & farinas2021: 2.29$\pm$0.75 $\to$ 0.16$\pm$0.08 MPa over ripening (endpoints); mapping ILLUSTRATIVE \\
\texttt{flesh\_\allowbreak{}poisson} & 0.45 & ILLUSTRATIVE & near-incompressible soft tissue \\
\texttt{bruise\_\allowbreak{}force\_\allowbreak{}ratio} & 0.6 & SOURCED-VERIFY & hia2018: a 10 N thumb press bruises firm-ripe fruit (visible within 48 h); ratio of thumb pressure to penetrometer pressure at stage 4 \\\bottomrule\end{tabularx}\end{table}
\paragraph{States, dynamics, tasks} Schema attributes: ripeness\_stage, integrity, seed\_coat, seed, skin, flesh\_condition, cut\_face\_browning; actions: SqueezeTest, Grasp, Halve, CutRoundSeed, TwistApart, RemoveSeed, Scoop. Ripeness advances with time on the Avocados Australia scale
(STANDARD); flesh bruises when contact pressure exceeds a fraction of the penetrometer pressure of the stage; halving
around the seed separates cleanly only from stage~3. \texttt{AvocadoBelief} updates the ripeness stage from squeeze
stiffness and skin darkness. The halving and pitting designators, the probabilistic node and everything the robot
does with them follow in Section~\ref{sec:avocado-reasoning}.
\subsection{Bench: ripeness by squeezing (Isaac Sim FEM)}
\texttt{scripts/sim\_avocado\_squeeze.py} squeezes a PhysX volume-deformable avocado (avocado\_4's outer surface,
uniform flesh, seed and skin not separate) between a fixed and a force-controlled plate at six ripeness stages; the
flesh modulus per stage maps the penetrometer force to the Farinas~2021 modulus range (endpoints SOURCED, mapping
ILLUSTRATIVE). The secant stiffness table feeds the belief's stiffness model. As a cross-check this report computes
the project's own Hertz model (\texttt{food.sdt.hertz\_squeeze\_stiffness}, two contacts on a sphere of the fruit's
diameter 65.4~mm, $\nu$~0.45) at 4~N. The FEM is stiffer than this sphere model by a nearly constant factor of 1.67--1.69 for stages 3--5. The avocado is not a sphere, so the factor is probably a shape effect (not established; in the tomato bench the FEM lies below Hertz instead), and the check confirms the trend over ripeness, not the absolute value. The factor falls to 1.46 at stage 2 and 1.13 at stage 1, and a separate probe run gave 11.9~N/mm for the stage-1 modulus against 7.0~N/mm in the final run (the probe's solver settings are not recorded in its log): the hard lanes are under-converged, as in the tomato and roll benches. An earlier 1000~Hz run (\texttt{avocado\_squeeze\_1000hz.json})
is invalid: its rigid control lane moved 9~mm, so all its stiffnesses were set to a placeholder.
\begin{table}[htbp]\centering\scriptsize\caption{Avocado squeeze stiffness at 4~N (test\_reports/avocado\_squeeze.json; Hertz column DERIVED in this report).}
\begin{tabularx}{\linewidth}{@{}rrrrr@{}}\toprule stage & E (MPa) & k FEM (N/mm) & k Hertz (N/mm) & FEM / Hertz \\\midrule
1 & 2.290 & 7.03 & 6.21 & 1.13 \\
2 & 1.041 & 5.36 & 3.67 & 1.46 \\
3 & 0.407 & 3.31 & 1.96 & 1.69 \\
4 & 0.248 & 2.37 & 1.41 & 1.67 \\
4.5 & 0.199 & 2.06 & 1.22 & 1.69 \\
5 & 0.160 & 1.77 & 1.05 & 1.68 \\\bottomrule\end{tabularx}\end{table}
With the stiffness table, one simulated squeeze of a hidden stage-4.2 fruit at 4~N (8~\% noise) moves the flat prior to a stage estimate of 4.37~$\pm$~0.31, and a second squeeze at 2~N to 4.29~$\pm$~0.23 with $P(\text{ready, stage} \ge 4) = 0.89$ (recomputed here with the plot script's seeds). The bruise-safe plate
force falls from 8~N (the largest tested) for hard fruit to 1~N for ripe fruit (panel d), consistent with the
Hort~Innovation warning that a 10~N thumb press bruises firm-ripe fruit (hia2018, SOURCED).
\begin{figure}[htbp]\centering\includegraphics[width=\linewidth,height=0.25\textheight,keepaspectratio]{/home/mbeetz/Documents/HouseholdTwin/report/figures/renders_food_avocado_squeeze.jpg}\caption{Avocado squeeze study (existing figure, tools/plot\_avocado\_squeeze.py): FEM force-compression curves, stiffness model, belief update, bruise-safe force.}\end{figure}
\clearpage\section{Wave 1: bread roll, banana, carrot}
Wave~1 was built with the twelve-step pipeline of the briefing document (scope, research notes, part tree, register,
geometry, USD, hidden models, beliefs, designators, bench, renders, review). The research notes are in
\texttt{docs/wave1/}; the summary page \texttt{docs/wave1/README.md} is generated by \texttt{tools/wave1\_doc.py}.
\begin{table}[htbp]\centering\scriptsize\caption{Calibration anchors of the wave-1 hidden-state models (from docs/wave1/README.md).}
\begin{tabularx}{\linewidth}{@{}p{0.2\linewidth}Xp{0.27\linewidth}p{0.18\linewidth}@{}}\toprule Model & Anchor & Twin & Status \\\midrule
Roll crust water / crispness & hirte\_2012: 9.0 \% @ 0.5 h, 13.9 \% @ 24 h, critical 10.1 \% @ 2.1 h & w = 9.73 t\^{}0.112 (10.6 \% @ 2.1 h) & CALIBRATED, residual +0.5 points \\
Roll crumb staling & lagrain\_2012 E 20 $\to$ 157 kPa (5 h .. 7 d) & log-time interpolation & SOURCED-VERIFY \\
Carrot boiling & carrot\_boil: 6 mm rounds 4-5 min crisp-tender, +1-2 min soft & crisp-tender 4.8 min, tender 6.9 min; whole 31 mm 16 min (source 10-15) & CALIBRATED thresholds on peng14 kinetics \\
Banana ripening & hou15 1 $\to$ 7 in 7 d (20 C, ethylene); ringer21 hue per stage & 0.6 stage/d, hue table & SOURCED-VERIFY \\\bottomrule\end{tabularx}\end{table}
\begin{table}[htbp]\centering\scriptsize\caption{Wave-1 beliefs scored on simulated observations (from docs/wave1/README.md; 200 simulated instances each).}
\begin{tabularx}{\linewidth}{@{}p{0.1\linewidth}Xp{0.4\linewidth}@{}}\toprule Item & Observations & Result \\\midrule
Roll & gripper squeeze stiffness (sd 15 \%), crackle count (Poisson) & freshness class 66 \% (squeeze k + crackles, 200 rolls) \\
Banana & camera peel hue (sd 3 deg) & stage RMS error 0.74 (one hue reading, 200 fingers) \\
Carrot & bend test (5 \% misread) & limp class 96 \% (one bend test, 200 roots) \\\bottomrule\end{tabularx}\end{table}
\FloatBarrier\subsection{Wheat roll (Weizenbrötchen)}\label{sec:roll}
A roll is a baked foam, not a plant organ: a dense browned crust round a crumb of gas cells. The twin is the oval single-score Weizenbrötchen; the score (Ausbund) is a separate crust part, the gas cells a point set that paints the cut faces. Built instance: \texttt{roll\_\allowbreak{}0}, 47.4~g; states: whole, half\_a (23.5 g), half\_b (23.9 g), top\_half (19.1 g), bottom\_half (28.3 g). Research notes: \texttt{docs/wave1/ROLL.md}.
\begin{figure}[htbp]\centering\begin{minipage}[t]{0.490\linewidth}\centering\includegraphics[width=\linewidth,height=0.2\textheight,keepaspectratio]{/home/mbeetz/Documents/HouseholdTwin/report/figures/renders_food_anatomy_roll_labelled.jpg}\\[-1pt]{\scriptsize anatomy plate}\end{minipage}\hfill\begin{minipage}[t]{0.490\linewidth}\centering\includegraphics[width=\linewidth,height=0.2\textheight,keepaspectratio]{/home/mbeetz/Documents/HouseholdTwin/report/figures/renders_food_wave1_roll_sliced.jpg}\\[-1pt]{\scriptsize sliced open}\end{minipage}\caption{Roll: exploded anatomy plate and cut halves (existing renders).}\end{figure}
\begin{table}[htbp]\centering\scriptsize\caption{Part tree of roll\_0.}
\begin{tabularx}{\linewidth}{@{}p{0.30\linewidth}p{0.27\linewidth}p{0.17\linewidth}rX@{}}\toprule Part (below /Food/Parts) & sdt:class & observability & n & mass (g) \\\midrule
\hspace*{0.0em}Crust & \texttt{food:BreadCrust} & visible &  &  \\
\hspace*{0.9em}UpperCrust & \texttt{food:UpperCrust} & visible &  & 3.77 \\
\hspace*{0.9em}ScoreCrust & \texttt{food:ScoreCrust} & visible &  & 0.29 \\
\hspace*{0.9em}BottomCrust & \texttt{food:BottomCrust} & visible &  & 3.41 \\
\hspace*{0.0em}Crumb & \texttt{food:BreadCrumb} & visible after cutting &  &  \\
\hspace*{0.9em}TransitionZone & \texttt{food:CrustCrumbTransition} & visible after cutting &  & 4.67 \\
\hspace*{0.9em}CrumbCore & \texttt{food:CrumbCore} & visible after cutting &  & 35.92 \\
\hspace*{0.9em}GasCells & \texttt{food:GasCell} & visible after cutting &  & 275,295 points \\
\hspace*{0.0em}Dusting & \texttt{food:FlourDusting} & visible &  &  \\\bottomrule\end{tabularx}\end{table}
\begin{table}[htbp]\centering\scriptsize\caption{Key register entries of the roll.}
\begin{tabularx}{\linewidth}{@{}p{0.19\linewidth}p{0.27\linewidth}p{0.12\linewidth}X@{}}\toprule Register key & Value / distribution & Status & Source and note \\\midrule
\texttt{height} & N(52, 3) in [44, 60] mm & SOURCED & hirte\_2012 $\sim$50 mm; musiienko\_2025 50-55 mm \\
\texttt{crust\_\allowbreak{}thickness\_\allowbreak{}top} & N(1.3, 0.2) in [0.9, 2] mm & SOURCED & hirte\_2012: cut-off crust 1.0-1.5 mm \\
\texttt{crumb\_\allowbreak{}density} & N(180, 15) in [150, 220] kg/m3 & DERIVED & roll specific volume 5.9-6.4 ml/g (kurek\_2016) $\to$ whole $\sim$0.16-0.17 g/cm3, minus the crust share; crumb 0.088-0.252 (zghal\_1999, SOURCED-VERIFY) \\
\texttt{crumb\_\allowbreak{}porosity} & U(0.73 .. 0.81) & SOURCED & hirte\_2012 uCT \\
\texttt{crumb\_\allowbreak{}E\_\allowbreak{}kPa\_\allowbreak{}by\_\allowbreak{}age} & 20 @ 5, 70 @ 48, 125 @ 96, 157 @ 168 kPa @ h & SOURCED-VERIFY & lagrain\_2012 pan bread, read off figures; log-time interpolation \\
\texttt{crumb\_\allowbreak{}poisson} & 0.103 & SOURCED & lagrain\_2012 \\
\texttt{crust\_\allowbreak{}water\_\allowbreak{}model} & w(t) = 9.73 t\_h\^{}0.112 (40 \% RH, 22 C) \% & CALIBRATED & power law through hirte\_2012 9.0 \% @ 0.5 h and 13.9 \% @ 24 h; predicts 10.6 \% @ 2.1 h (measured: 10.1 \% critical reached at 2.1 h) $\to$ residual +0.5 points \\
\texttt{crust\_\allowbreak{}water\_\allowbreak{}critical} & 10.1 \% & SOURCED & hirte\_2012: half the sensory crispness lost \\
\texttt{crust\_\allowbreak{}break\_\allowbreak{}force} & 3 N & SOURCED & hirte\_2012: $\ge$ 3 N with a 1 x 10 mm blade \\
\texttt{squeeze\_\allowbreak{}k\_\allowbreak{}model} & k = 0.29 E / 20 kPa (30 mm pads across the width, 1 $\to$ 3 N) N/mm & DERIVED & Isaac FEM (test\_reports/roll\_squeeze\_summary.json): 20 kPa lane converged (0.276 $\to$ 0.29 N/mm, 1000 $\to$ 2000 Hz); stiffer lanes under-converged $\to$ linear-elastic scaling; crust excluded (GAP) \\
\texttt{saw\_\allowbreak{}setpoint} & -30 N, 0.02 m/s, $\pm$0.02 m strokes (baguette) N, m/s, m & SOURCED & luo\_2026 (robot) \\\bottomrule\end{tabularx}\end{table}
\paragraph{States and dynamics} Schema attributes: freshness, crust\_texture, cut\_state, toast\_level, temperature, topping, support. Crust softening is water migration from the crumb (power law CALIBRATED through two hirte\_2012 points),
crumb staling a log-time interpolation of the Lagrain~2012 modulus (SOURCED-VERIFY, pan bread).\begin{table}[htbp]\centering\scriptsize\caption{Task designators for the wheat roll (weizenbrötchen). Demo: symbolic run on seed 0 by this report (default parameters; a smoke test, not a validation).}
\begin{tabularx}{\linewidth}{@{}p{0.16\linewidth}Xp{0.3\linewidth}l@{}}\toprule Designator & Goal & Twin model & Demo \\\midrule
PickRoll & roll held, not dented & RollTruth.grasp\_corridor & ok \\
SliceRollOpen & cut\_state = sliced\_open (or hinged), crumb not squashed & RollTruth.saw & ok \\
ReheatRoll & crust\_texture = crisp, temperature = warm & RollTruth.reheat & ok \\
ToastRollHalves & toast\_level = golden on the cut faces & RollTruth.toast\_faces & ok \\\bottomrule\end{tabularx}\end{table}
\paragraph{Bench: staling squeeze (Isaac Sim FEM)} \texttt{scripts/sim\_roll\_squeeze.py} squeezes a
crumb-only FEM roll (crust not separate: GAP) between 30~mm pads at four crumb moduli. Only the fresh lane converges
(+4.5~\% between solver settings); the stiffer lanes still rise with the solver rate and are lower bounds. The
register therefore uses the converged lane scaled linearly in $E$ (\texttt{squeeze\_k\_model}, DERIVED).
\begin{table}[htbp]\centering\scriptsize\caption{Roll squeeze, incremental stiffness 1--3~N in N/mm (test\_reports/roll\_squeeze\_summary.json).}
\begin{tabularx}{\linewidth}{@{}rrrrrr@{}}\toprule E (kPa) & age (h) & k 1000 Hz & k 2000 Hz & change \% & k model (N/mm) \\\midrule
20 & 5 & 0.276 & 0.289 & +4.5 & 0.29 \\
70 & 48 & 0.497 & 0.631 & +21.2 & 1.01 \\
125 & 96 & 0.546 & 0.797 & +31.5 & 1.81 \\
157 & 168 & 0.535 & 0.876 & +38.9 & 2.28 \\\bottomrule\end{tabularx}\end{table}
\FloatBarrier\subsection{Cavendish banana}\label{sec:banana}
A banana finger is a parthenocarpic berry of three carpels. The peel (exocarp and outer mesocarp) is four tissue layers with latex vessels; the pulp three carpel segments round the placental axis with ovule remnants and the 'strings'. Four finger instances are built. Built instance: \texttt{banana\_\allowbreak{}0}, 117.1~g; states: whole, half\_a (58.3 g), half\_b (58.8 g), tip\_half (59.1 g), stalk\_half (58.0 g). Research notes: \texttt{docs/wave1/BANANA.md}.
\begin{figure}[htbp]\centering\begin{minipage}[t]{0.490\linewidth}\centering\includegraphics[width=\linewidth,height=0.2\textheight,keepaspectratio]{/home/mbeetz/Documents/HouseholdTwin/report/figures/renders_food_anatomy_banana_labelled.jpg}\\[-1pt]{\scriptsize anatomy plate}\end{minipage}\hfill\begin{minipage}[t]{0.490\linewidth}\centering\includegraphics[width=\linewidth,height=0.2\textheight,keepaspectratio]{/home/mbeetz/Documents/HouseholdTwin/report/figures/renders_food_wave1_banana_half.jpg}\\[-1pt]{\scriptsize lengthwise half}\end{minipage}\caption{Banana: exploded anatomy plate and lengthwise half (existing renders).}\end{figure}
\begin{table}[htbp]\centering\scriptsize\caption{Part tree of banana\_0.}
\begin{tabularx}{\linewidth}{@{}p{0.30\linewidth}p{0.27\linewidth}p{0.17\linewidth}rX@{}}\toprule Part (below /Food/Parts) & sdt:class & observability & n & mass (g) \\\midrule
\hspace*{0.0em}Peel & \texttt{food:BananaPeel} & visible &  &  \\
\hspace*{0.9em}Epidermis & \texttt{food:Epidermis} & visible &  & 0.58 \\
\hspace*{0.9em}Chlorenchyma & \texttt{food:Chlorenchyma} & hidden &  & 4.58 \\
\hspace*{0.9em}Parenchyma & \texttt{food:PeelParenchyma} & visible after peeling &  & 21.43 \\
\hspace*{0.9em}InnerAerenchyma & \texttt{food:Aerenchyma} & visible after peeling &  & 4.76 \\
\hspace*{0.9em}Laticifers & \texttt{food:Laticifer} & hidden &  & 1,500 points \\
\hspace*{0.0em}Stalk & \texttt{food:Pedicel} & visible &  & 1.34 \\
\hspace*{0.0em}CrownCushion & \texttt{food:CrownRemnant} & visible &  & 0.21 \\
\hspace*{0.0em}BlossomTip & \texttt{food:PerianthRemnant} & visible &  & 0.01 \\
\hspace*{0.0em}Pulp & \texttt{food:BananaPulp} & visible after peeling &  &  \\
\hspace*{0.9em}Segment\_1..3 & \texttt{food:PulpSegment} & visible after peeling & 3 & 83.00 \\
\hspace*{0.9em}CentralAxis & \texttt{food:Placenta} & visible after cutting &  & 1.05 \\
\hspace*{0.9em}Ovules & \texttt{food:AbortedOvule} & visible after cutting &  & 258 points \\
\hspace*{0.9em}PhloemStrings & \texttt{food:PhloemBundle} & visible after peeling &  &  \\
\hspace*{1.8em}String\_1..3 & \texttt{food:PhloemBundle} & visible after peeling & 3 & 0.12 \\\bottomrule\end{tabularx}\end{table}
\begin{table}[htbp]\centering\scriptsize\caption{Key register entries of the banana.}
\begin{tabularx}{\linewidth}{@{}p{0.19\linewidth}p{0.27\linewidth}p{0.12\linewidth}X@{}}\toprule Register key & Value / distribution & Status & Source and note \\\midrule
\texttt{length\_\allowbreak{}outer} & N(210, 15) in [170, 250] mm & SOURCED & convex-face length incl. pedicel (eu1333\_2011 definition): 190 (mal22), 233 (ciro05) \\
\texttt{calibre} & N(35, 1.5) in [30, 40] mm & SOURCED & ciro05 34/36/34 mm; eu1333 min 27 mm \\
\texttt{peel\_\allowbreak{}thickness} & N(2.6, 0.25) in [2, 3.2] mm & SOURCED & mal22 2.56 mm (Grand Naine); amn12 1.1-3.8 mm, thinning on ripening \\
\texttt{ripeness\_\allowbreak{}stage} & U(2 .. 7) 1..7 & STANDARD & markon / ma22 seven-stage colour scale \\
\texttt{peel\_\allowbreak{}hue\_\allowbreak{}by\_\allowbreak{}stage} & [118, 98.7, 97.3, 92.7, 89.4, 87.5, 82] deg & SOURCED-VERIFY & ringer21 R2..R7; stage 1 from col20 Gros Michel day 0 (h $\sim$122) $\to$ 118 \\
\texttt{pulp\_\allowbreak{}E\_\allowbreak{}compression} & 1.75 (day 1) $\to$ 1.09 (day 12) MPa & SOURCED & ciro05 cylinders 20 x 40 mm \\
\texttt{peel\_\allowbreak{}fracture} & 53.0 $\to$ 2.1 on ripening N & SOURCED-VERIFY & khanal23 (test geometry not in abstract) \\
\texttt{finger\_\allowbreak{}drop\_\allowbreak{}force} & $\sim$49 on day 1, > 2x lower by day 3 N & SOURCED-VERIFY & imsabai; mbeguie09 \\
\texttt{peel\_\allowbreak{}bond\_\allowbreak{}ratios} & (0.45, 0.3, 1.5) & ILLUSTRATIVE & peel-to-pulp adhesion / ridge seam / strip tear strength as multiples of the peel fracture force (khanal23 53 $\to$ 2.1 N, SOURCED-VERIFY); adhesion > seam so the peel splits along the ridges (khanal23) instead of sliding off as a sleeve; GAP: no peeling-force study \\
\texttt{stage\_\allowbreak{}rate\_\allowbreak{}20C} & N(0.6, 0.15) in [0.3, 1] stage/d & SOURCED-VERIFY & hou15: 1 $\to$ 7 in 7 d with ethylene (0.86/d); ma22 day bands $\sim$0.4/d; 18 C RS6 at 9 d \\
\texttt{ridges} & one of 4, 5 & ILLUSTRATIVE & 'slightly ridged' (mal22), 'nearly 4-sided'; count GAP \\\bottomrule\end{tabularx}\end{table}
\paragraph{States and dynamics} Schema attributes: ripeness\_stage, edible, peel\_state, neck, attached\_to\_hand, bruised, chilling\_injury, pulp\_browning, integrity. Ripening on the seven-stage colour scale (STANDARD) at a SOURCED-VERIFY rate; peel hue per stage
SOURCED-VERIFY; the peeling bonds are ILLUSTRATIVE ratios of the peel fracture force, because no published peeling
force was found (GAP).\begin{table}[htbp]\centering\scriptsize\caption{Task designators for the cavendish banana. Demo: symbolic run on seed 0 by this report (default parameters; a smoke test, not a validation).}
\begin{tabularx}{\linewidth}{@{}p{0.16\linewidth}Xp{0.3\linewidth}l@{}}\toprule Designator & Goal & Twin model & Demo \\\midrule
SeparateBanana & one finger free, peel at the pedicel intact & BananaTruth.separate\_from\_hand & ok \\
PeelBanana & peel\_state = off, pulp unbruised, strings mostly on the peel & BananaTruth.snap\_neck, peel\_strip & ok \\
SliceBananaRounds & integrity = rounds of $\sim$6 mm & BananaTruth.cut\_rounds, browning\_class & ok \\\bottomrule\end{tabularx}\end{table}
\paragraph{Bench: peeling (Isaac Sim)} \texttt{scripts/sim\_banana\_peel.py} models the peel as 4 strips of 12 rigid segments joined by breakable
adhesion, seam and hinge joints, and pulls strip~1 from the snapped neck with a spring finger. The force level
inherits the ILLUSTRATIVE bond ratios, so only the trend is meaningful: green fruit needs higher peaks and its strip
may tear short; ripe fruit peels cleanly at a fraction of a newton. The stage-4 strip tore below its nominal strength
(joint forces of light segments are spiky), and the outcome flips between runs. The rendered filmstrip shows the strips
only faintly.
\begin{table}[htbp]\centering\scriptsize\caption{Banana peel bench (test\_reports/banana\_peel.json).}
\begin{tabularx}{\linewidth}{@{}rlrrrrX@{}}\toprule stage & adh. / seam / tear (N) & peak (N) & median (N) & strip 1 & others & outcome \\\midrule
2 & 13.9 / 9.3 / 46.4 & 14.2 & 2.67 & 11/12 & 0 & strip peeled off \\
4 & 4.7 / 3.2 / 15.8 & 7.1 & 1.26 & 5/12 & 0 & strip torn short \\
6 & 1.6 / 1.1 / 5.4 & 2.5 & 0.43 & 12/12 & 2 & strip peeled off \\\bottomrule\end{tabularx}\end{table}
\FloatBarrier\subsection{Nantes carrot}\label{sec:carrot}
A carrot is a storage taproot: periderm, orange secondary phloem, vascular cambium, paler secondary xylem (the core) and the primary xylem; the crown carries petiole stubs, lenticel bands and lateral-root scars are point sets. Built instance: \texttt{carrot\_\allowbreak{}0}, 115.0~g; states: whole, half\_a (57.4 g), half\_b (57.6 g), tip\_piece (52.4 g), crown\_piece (62.6 g). Research notes: \texttt{docs/wave1/CARROT.md}.
\begin{figure}[htbp]\centering\begin{minipage}[t]{0.490\linewidth}\centering\includegraphics[width=\linewidth,height=0.2\textheight,keepaspectratio]{/home/mbeetz/Documents/HouseholdTwin/report/figures/renders_food_anatomy_carrot_labelled.jpg}\\[-1pt]{\scriptsize anatomy plate}\end{minipage}\hfill\begin{minipage}[t]{0.490\linewidth}\centering\includegraphics[width=\linewidth,height=0.2\textheight,keepaspectratio]{/home/mbeetz/Documents/HouseholdTwin/report/figures/renders_food_wave1_carrot_crown_cut.jpg}\\[-1pt]{\scriptsize crown piece}\end{minipage}\caption{Carrot: exploded anatomy plate and crown piece (existing renders).}\end{figure}
\begin{table}[htbp]\centering\scriptsize\caption{Part tree of carrot\_0.}
\begin{tabularx}{\linewidth}{@{}p{0.30\linewidth}p{0.27\linewidth}p{0.17\linewidth}rX@{}}\toprule Part (below /Food/Parts) & sdt:class & observability & n & mass (g) \\\midrule
\hspace*{0.0em}Periderm & \texttt{food:Periderm} & visible &  & 2.51 \\
\hspace*{0.0em}SecondaryPhloem & \texttt{food:SecondaryPhloem} & visible after peeling &  & 101.41 \\
\hspace*{0.0em}VascularCambium & \texttt{food:VascularCambium} & visible after cutting &  & 0.76 \\
\hspace*{0.0em}SecondaryXylem & \texttt{food:SecondaryXylem} & visible after cutting &  & 10.09 \\
\hspace*{0.0em}PrimaryXylem & \texttt{food:PrimaryXylem} & visible after cutting &  & 0.14 \\
\hspace*{0.0em}Crown & \texttt{food:Crown} & visible &  &  \\
\hspace*{0.9em}Shoulder &  &  &  &  \\
\hspace*{0.9em}PetioleStubs & \texttt{food:PetioleStubs} & visible &  &  \\
\hspace*{1.8em}Stub\_1..7 & \texttt{food:PetioleStub} & visible & 7 & 0.09 \\
\hspace*{0.0em}RootTip &  &  &  &  \\
\hspace*{0.0em}LateralRootScars & \texttt{food:LateralRootScar} & visible &  & 48 points \\
\hspace*{0.0em}LenticelBands & \texttt{food:LenticelBand} & visible &  & 843 points \\\bottomrule\end{tabularx}\end{table}
\begin{table}[htbp]\centering\scriptsize\caption{Key register entries of the carrot.}
\begin{tabularx}{\linewidth}{@{}p{0.19\linewidth}p{0.27\linewidth}p{0.12\linewidth}X@{}}\toprule Register key & Value / distribution & Status & Source and note \\\midrule
\texttt{length} & N(190, 15) in [150, 230] mm & SOURCED & con22: 182.7 / 210.8 mm (Nantes) \\
\texttt{shoulder\_\allowbreak{}diameter} & N(31, 2) in [25, 38] mm & SOURCED & con22: 30.3 / 31.4 mm \\
\texttt{xylem\_\allowbreak{}radius\_\allowbreak{}fraction} & N(0.33, 0.04) in [0.25, 0.42] & SOURCED-VERIFY & singh17 $\sim$0.38 over mixed genotypes; Nantes bred for a small core; vob24 trend crown $\to$ tip; peng14 4-7 mm core \\
\texttt{E\_\allowbreak{}phloem} & N(4.6, 0.15) in [4.33, 4.85] MPa & SOURCED-VERIFY & vob24 after Stopa 2017 \\
\texttt{poisson} & 0.475 & SOURCED-VERIFY & vob24 \\
\texttt{bend\_\allowbreak{}Fmax} & 48.6 N & SOURCED-VERIFY & jah18 three-point bending (span: see source) \\
\texttt{blade\_\allowbreak{}shear} & 41.1 N & SOURCED-VERIFY & jah18 straight blade 1.4 mm, 30 deg \\
\texttt{limp\_\allowbreak{}rwc} & 0.95 & SOURCED-VERIFY & herp99\_calbo00: firmness from turgor $\sim$0 below RWC 0.95 \\
\texttt{softening\_\allowbreak{}kinetics} & 2nd-order fractional conversion; k(80 C) 0.035, k(110 C) 1.453; Ea 138.9 kJ/mol; Finf/F0 0.02 1/min & SOURCED & peng14 \\
\texttt{peel\_\allowbreak{}depth} & 2 mm & SOURCED & con22 hand peeler; mass 12.1-14.6 \% \\
\texttt{boil\_\allowbreak{}rounds\_\allowbreak{}6mm} & 4-5 crisp-tender, +1-2 soft min & SOURCED-VERIFY & carrot\_boil \\\bottomrule\end{tabularx}\end{table}
\paragraph{States and dynamics} Schema attributes: firmness\_class, peel\_state, trim\_state, cut\_shape, cooked\_state, surface\_blush, green\_shoulder. Boiling uses the Peng~2014 softening kinetics (SOURCED) with a conduction lag and CALIBRATED doneness
thresholds; limpness follows turgor and relative water content (SOURCED-VERIFY).\begin{table}[htbp]\centering\scriptsize\caption{Task designators for the nantes carrot. Demo: symbolic run on seed 0 by this report (default parameters; a smoke test, not a validation).}
\begin{tabularx}{\linewidth}{@{}p{0.16\linewidth}Xp{0.3\linewidth}l@{}}\toprule Designator & Goal & Twin model & Demo \\\midrule
PeelCarrot & peel\_state = peeled, $\sim$2 mm removed & CarrotTruth.peel\_stroke & ok \\
CutCarrotRounds & cut\_shape = rounds (5-6 mm) & CarrotTruth.cut & ok \\
BoilCarrot & cooked\_state = crisp\_tender (or tender) & carrot\_models.boil (peng14 + conduction lag) & ok \\\bottomrule\end{tabularx}\end{table}
\paragraph{Bench: why cooks plank a carrot (Isaac Sim)} \texttt{scripts/sim\_carrot\_plank.py} presses a knife
onto carrot\_0 up to the cutting force (44~N, the jah18 blade force scaled to this
root, SOURCED-VERIFY), straight down or with a slicing stroke, on the whole root and on a planked half. A straight
press acts through the axis and hardly moves the round root; a stroke rolls it away at once. The planked half holds
during strokes up to about 10~N, then the blade friction yaws it -- a claw-grip hand is needed. Escape criterion: roll > 10 deg or slide > 5 mm before F\_cut.
\begin{table}[htbp]\centering\scriptsize\caption{Carrot plank bench (test\_reports/carrot\_plank.json).}
\begin{tabularx}{\linewidth}{@{}llrrr@{}}\toprule carrot & knife & roll at F\textsubscript{cut} (deg) & slide (mm) & escapes at (N) \\\midrule
whole & press & 5.1 & 1.5 & -- \\
whole & stroke & 26.1 & 7.2 & 1 \\
planked & press & 0.1 & 0.0 & -- \\
planked & stroke & 105.1 & 389.5 & 10 \\\bottomrule\end{tabularx}\end{table}
\begin{figure}[htbp]\centering\includegraphics[width=\linewidth,height=0.25\textheight,keepaspectratio]{/home/mbeetz/Documents/HouseholdTwin/report/figures/banana_carrot_benches.pdf}\caption{Banana and carrot bench traces drawn by this report from the stored logs: (a) finger force during peeling per ripeness stage; (b) roll angle of the carrot over knife force (log scale above 1 deg).}\end{figure}
\clearpage\section{Wave 2: loaf, onion, orange, potato, tomato}
The five wave-2 items are built as USD assets with cut states and have exploded anatomy plates, research notes in
\texttt{docs/wave2/} and ten task designators (\texttt{food/wave2\_tasks.py}). Path-traced renders of each item, whole and cut, were added on 20~October. They have no entries in
\texttt{food/schema.py} (their discrete state comes from the truth classes' \texttt{discrete()}), and only the
tomato has a physics bench.
\begin{figure}[htbp]\centering\includegraphics[width=\linewidth,height=0.22\textheight,keepaspectratio]{/home/mbeetz/Documents/HouseholdTwin/report/figures/renders_food_wave2_group.jpg}\caption{Wave~2 items, path-traced: sandwich loaf, onion, orange, potato and tomato.}\end{figure}
\FloatBarrier\subsection{Sandwich loaf (Toastbrot)}\label{sec:loaf}
Industrial German Toastbrot, baked in a lidded tin: thin soft crust on the tin faces and the flat top, square cut ends (most packs have no heel slices), pre-sliced into 20 slices. Discrete attributes exposed by the truth class: freshness, mould\_risk, slices\_left. Built instance: \texttt{loaf\_\allowbreak{}0}, 494.6~g; states: whole, half\_a (247.1 g), half\_b (247.5 g), slices (494.6 g). Research notes: \texttt{docs/wave2/LOAF.md}.
\begin{figure}[htbp]\centering\begin{minipage}[t]{0.490\linewidth}\centering\includegraphics[width=\linewidth,height=0.19\textheight,keepaspectratio]{/home/mbeetz/Documents/HouseholdTwin/report/figures/renders_food_anatomy_loaf_labelled.jpg}\\[-1pt]{\scriptsize exploded anatomy plate}\end{minipage}\hfill\begin{minipage}[t]{0.490\linewidth}\centering\includegraphics[width=\linewidth,height=0.19\textheight,keepaspectratio]{/home/mbeetz/Documents/HouseholdTwin/report/figures/renders_food_item_loaf.jpg}\\[-1pt]{\scriptsize whole and cut half on the worktop}\end{minipage}\caption{loaf\_0: anatomy plate and path-traced render (Isaac Sim 6.2).}\end{figure}
\begin{table}[htbp]\centering\scriptsize\caption{Part tree of loaf\_0 (indexed siblings merged).}
\begin{tabularx}{\linewidth}{@{}p{0.30\linewidth}p{0.27\linewidth}p{0.17\linewidth}rX@{}}\toprule Part (below /Food/Parts) & sdt:class & observability & n & mass (g) \\\midrule
\hspace*{0.0em}Crust & \texttt{food:BreadCrust} & visible &  &  \\
\hspace*{0.9em}TopCrust & \texttt{food:TopCrust} & visible &  & 8.70 \\
\hspace*{0.9em}TinCrust & \texttt{food:TinCrust} & visible &  & 8.20 \\
\hspace*{0.0em}Crumb & \texttt{food:BreadCrumb} & visible after cutting &  & 484.62 \\
\hspace*{0.9em}GasCells & \texttt{food:GasCell} & visible after cutting &  & 6,217,490 points \\\bottomrule\end{tabularx}\end{table}
\begin{table}[htbp]\centering\scriptsize\caption{Key register entries of the sandwich loaf (toastbrot).}
\begin{tabularx}{\linewidth}{@{}p{0.19\linewidth}p{0.27\linewidth}p{0.12\linewidth}X@{}}\toprule Register key & Value / distribution & Status & Source and note \\\midrule
\texttt{length} & N(250, 6) in [235, 265] mm & SOURCED-VERIFY & toast\_pack $\sim$0.25 m \\
\texttt{crumb\_\allowbreak{}density} & N(190, 10) in [170, 215] kg/m3 & SOURCED & nuss90 0.19 g/cm3 \\
\texttt{slices} & 20 & SOURCED-VERIFY & toast\_pack \\
\texttt{heels} & one of 0, 2 & SOURCED & test25\_toast: most packs no heels; organic 2 \\
\texttt{crust\_\allowbreak{}thickness\_\allowbreak{}top} & U(0.8 .. 1.3) mm & ILLUSTRATIVE & thin soft lidded crust; GAP \\
\texttt{firmness\_\allowbreak{}by\_\allowbreak{}day} & 1.98 / 3.62 / 4.18 / 5.64 at day 1 / 5 / 9 / 13 (25 C, bagged) N & SOURCED & sehn20 lab white pan loaf \\
\texttt{tensile\_\allowbreak{}failure} & 0.31-0.39 (12 mm neck) at 9.2-9.7 mm N & SOURCED & nuss90 \\
\texttt{toast\_\allowbreak{}time} & 80 .. > 150 to 'best' (per toaster) s & SOURCED & toaster26 \\
\texttt{fridge\_\allowbreak{}staling\_\allowbreak{}factor} & 6 & TYPICAL & toast\_misc (popular, weak) \\\bottomrule\end{tabularx}\end{table}
\begin{table}[htbp]\centering\scriptsize\caption{Task designators for the sandwich loaf (toastbrot). Demo: symbolic run on seed 0 by this report (default parameters; a smoke test, not a validation).}
\begin{tabularx}{\linewidth}{@{}p{0.16\linewidth}Xp{0.3\linewidth}l@{}}\toprule Designator & Goal & Twin model & Demo \\\midrule
TakeToastSlice & one slice out of the bag, untorn; bag closed again & LoafTruth.take\_slice; Isaac: loaf 'slices' state (one rigid body per slice) & ok \\
ToastSlice & toast\_level = golden & LoafTruth.toast & ok \\\bottomrule\end{tabularx}\end{table}
\FloatBarrier\subsection{Yellow onion}\label{sec:onion}
Swollen leaf bases (fleshy scales) on a compressed stem, the basal plate, under dry papery outer skins; each scale is a revolved region converging on the plate. Discrete attributes exposed by the truth class: cut\_state, peel\_state, sprouting, trim\_state. Built instance: \texttt{onion\_\allowbreak{}0}, 146.5~g; states: whole, half\_a (73.5 g), half\_b (73.0 g), top\_piece (56.1 g), root\_piece (90.4 g). Research notes: \texttt{docs/wave2/ONION.md}.
\begin{figure}[htbp]\centering\begin{minipage}[t]{0.490\linewidth}\centering\includegraphics[width=\linewidth,height=0.19\textheight,keepaspectratio]{/home/mbeetz/Documents/HouseholdTwin/report/figures/renders_food_anatomy_onion_labelled.jpg}\\[-1pt]{\scriptsize exploded anatomy plate}\end{minipage}\hfill\begin{minipage}[t]{0.490\linewidth}\centering\includegraphics[width=\linewidth,height=0.19\textheight,keepaspectratio]{/home/mbeetz/Documents/HouseholdTwin/report/figures/renders_food_item_onion.jpg}\\[-1pt]{\scriptsize whole and cut half on the worktop}\end{minipage}\caption{onion\_0: anatomy plate and path-traced render (Isaac Sim 6.2).}\end{figure}
\begin{table}[htbp]\centering\scriptsize\caption{Part tree of onion\_0 (indexed siblings merged).}
\begin{tabularx}{\linewidth}{@{}p{0.30\linewidth}p{0.27\linewidth}p{0.17\linewidth}rX@{}}\toprule Part (below /Food/Parts) & sdt:class & observability & n & mass (g) \\\midrule
\hspace*{0.0em}DryScales & \texttt{food:DryScaleGroup} & visible &  &  \\
\hspace*{0.9em}Tunic\_1..2 & \texttt{food:DryScale} & visible after peeling & 2 & 1.49 \\
\hspace*{0.0em}FleshyScales & \texttt{food:FleshyScaleGroup} & visible after peeling &  &  \\
\hspace*{0.9em}Scale\_1..7 & \texttt{food:FleshyScale} & visible after cutting & 7 & 138.95 \\
\hspace*{0.0em}ShootApex & \texttt{food:ShootApex} & visible after cutting &  & 4.91 \\
\hspace*{0.0em}BasalPlate & \texttt{food:BasalPlate} & visible &  & 0.88 \\
\hspace*{0.0em}Roots & \texttt{food:RootCrown} & visible &  &  \\
\hspace*{0.9em}Root\_1..25 & \texttt{food:Root} & visible & 25 & 0.03 \\
\hspace*{0.0em}Neck & \texttt{food:Neck} & visible &  & 0.23 \\\bottomrule\end{tabularx}\end{table}
\begin{table}[htbp]\centering\scriptsize\caption{Key register entries of the yellow onion.}
\begin{tabularx}{\linewidth}{@{}p{0.19\linewidth}p{0.27\linewidth}p{0.12\linewidth}X@{}}\toprule Register key & Value / distribution & Status & Source and note \\\midrule
\texttt{diameter} & N(65, 7) in [50, 85] mm & SOURCED & hb66\_onion medium 51-83 mm; bah04 51.7 \\
\texttt{n\_\allowbreak{}fleshy} & one of 7, 8 & SOURCED-VERIFY & kim\_heb20 7.6 $\pm$ 0.2 \\
\texttt{scale\_\allowbreak{}thickness\_\allowbreak{}outer} & U(3.5 .. 4.5) mm & ILLUSTRATIVE & kim\_heb20 5.5 mm (bulk mean), wu25 $\sim$3 mm; constrained by count x radius \\
\texttt{n\_\allowbreak{}dry} & one of 2, 3 & ILLUSTRATIVE & GAP (ordinal skin-thickness classes only) \\
\texttt{flesh\_\allowbreak{}density} & 1040 kg/m3 & SOURCED-VERIFY & bah04 bulb 1.04 g/cm3 \\
\texttt{bulb\_\allowbreak{}E} & 2.03-5.45, nu 0.26-0.45 MPa & SOURCED-VERIFY & jafari16 \\
\texttt{tissue\_\allowbreak{}E} & 2.2 skin on / 1.4 off MPa & SOURCED & wu25 \\
\texttt{crushing\_\allowbreak{}load} & 400-780 N & SOURCED-VERIFY & bah04 \\
\texttt{lf\_\allowbreak{}window} & peak $\sim$1, gone $\sim$10 min & SOURCED & lokke12 \\
\texttt{dice\_\allowbreak{}spacing} & mince $\le$3, fine 3-6, medium 6-13, coarse 13-19 mm & SOURCED & onion\_howto \\\bottomrule\end{tabularx}\end{table}
\begin{table}[htbp]\centering\scriptsize\caption{Task designators for the yellow onion. Demo: symbolic run on seed 0 by this report (default parameters; a smoke test, not a validation).}
\begin{tabularx}{\linewidth}{@{}p{0.16\linewidth}Xp{0.3\linewidth}l@{}}\toprule Designator & Goal & Twin model & Demo \\\midrule
DiceOnion & even dice of the target size, root end kept until last & OnionTruth.dice, onion\_models.lf\_dose & ok \\
SauteOnion & translucent | golden | caramelised as asked & onion\_models.saute\_state & ok \\\bottomrule\end{tabularx}\end{table}
\FloatBarrier\subsection{Sweet orange}\label{sec:orange}
A hesperidium: flavedo with oil glands, spongy albedo, segment membranes (endocarp) with juice vesicles, three vascular bundles per segment and the central axis; navel or Valencia. Discrete attributes exposed by the truth class: juice, peel\_state, segments. Built instance: \texttt{orange\_\allowbreak{}0}, 193.6~g; states: whole, stem\_half (88.7 g), navel\_half (104.9 g), half\_a (96.9 g), half\_b (96.7 g). Research notes: \texttt{docs/wave2/ORANGE.md}.
\begin{figure}[htbp]\centering\begin{minipage}[t]{0.490\linewidth}\centering\includegraphics[width=\linewidth,height=0.19\textheight,keepaspectratio]{/home/mbeetz/Documents/HouseholdTwin/report/figures/renders_food_anatomy_orange_labelled.jpg}\\[-1pt]{\scriptsize exploded anatomy plate}\end{minipage}\hfill\begin{minipage}[t]{0.490\linewidth}\centering\includegraphics[width=\linewidth,height=0.19\textheight,keepaspectratio]{/home/mbeetz/Documents/HouseholdTwin/report/figures/renders_food_item_orange.jpg}\\[-1pt]{\scriptsize whole and cut half on the worktop}\end{minipage}\caption{orange\_0: anatomy plate and path-traced render (Isaac Sim 6.2).}\end{figure}
\begin{table}[htbp]\centering\scriptsize\caption{Part tree of orange\_0 (indexed siblings merged).}
\begin{tabularx}{\linewidth}{@{}p{0.30\linewidth}p{0.27\linewidth}p{0.17\linewidth}rX@{}}\toprule Part (below /Food/Parts) & sdt:class & observability & n & mass (g) \\\midrule
\hspace*{0.0em}Peel & \texttt{food:CitrusPeel} & visible &  &  \\
\hspace*{0.9em}Flavedo & \texttt{food:Flavedo} & visible &  & 17.88 \\
\hspace*{0.9em}Albedo & \texttt{food:Albedo} & visible after peeling &  & 36.09 \\
\hspace*{0.9em}OilGlands & \texttt{food:OilGland} & visible &  & 8,666 points \\
\hspace*{0.0em}StemButton & \texttt{food:Calyx} & visible &  & 0.07 \\
\hspace*{0.0em}Core & \texttt{food:Columella} & visible after cutting &  & 2.13 \\
\hspace*{0.0em}Segments & \texttt{food:SegmentGroup} & visible after peeling &  &  \\
\hspace*{0.9em}Segment\_1..10 & \texttt{food:Segment} & visible after peeling & 10 &  \\
\hspace*{1.8em}Membrane & \texttt{food:SegmentMembrane} & visible after peeling & 10 & 4.32 \\
\hspace*{1.8em}Vesicles & \texttt{food:JuiceVesicles} & visible after cutting & 10 & 133.01 \\
\hspace*{1.8em}Seed\_1..4 & \texttt{food:CitrusSeed} & visible after cutting & 4 & 0.13 \\\bottomrule\end{tabularx}\end{table}
\begin{table}[htbp]\centering\scriptsize\caption{Key register entries of the sweet orange.}
\begin{tabularx}{\linewidth}{@{}p{0.19\linewidth}p{0.27\linewidth}p{0.12\linewidth}X@{}}\toprule Register key & Value / distribution & Status & Source and note \\\midrule
\texttt{cultivar} & one of navel, valencia & ILLUSTRATIVE & German supermarket: navel (winter) / Valencia (summer) \\
\texttt{diameter} & N(75, 6) in [61, 88] mm & SOURCED & ucr\_citrus Washington navel 61-88 mm \\
\texttt{peel\_\allowbreak{}thickness} & N(5, 0.6) in [3.5, 6.4] mm & SOURCED & ucr\_citrus rind 4.1-6.4 / 3.5-6.2 mm \\
\texttt{segments} & one of 9, 10, 11, 12 & SOURCED-VERIFY & citrus\_patents 9.6-11 \\
\texttt{flavedo\_\allowbreak{}E} & 103 MPa & SOURCED & smith18 (30-300 across species) \\
\texttt{albedo\_\allowbreak{}E} & 0.79 MPa & SOURCED & smith18 ($\pm$0.58) \\
\texttt{gland\_\allowbreak{}burst\_\allowbreak{}pressure} & 30-140 kPa & SOURCED & smith18 \\
\texttt{juice\_\allowbreak{}fraction} & U(45 .. 54) \% & SOURCED & ucr\_citrus navel 45-54 (Valencia 37-56) \\
\texttt{peel\_\allowbreak{}strip\_\allowbreak{}force} & navel 4, Valencia 8 per 16 mm strip N & ILLUSTRATIVE & ranking from ritenour13; GAP \\
\texttt{bitter\_\allowbreak{}juice\_\allowbreak{}h} & 3 h & SOURCED-VERIFY & edis891 'within several hours' (navel) \\\bottomrule\end{tabularx}\end{table}
\begin{table}[htbp]\centering\scriptsize\caption{Task designators for the sweet orange. Demo: symbolic run on seed 0 by this report (default parameters; a smoke test, not a validation).}
\begin{tabularx}{\linewidth}{@{}p{0.16\linewidth}Xp{0.3\linewidth}l@{}}\toprule Designator & Goal & Twin model & Demo \\\midrule
JuiceOrange & juice in the glass, few seeds & OrangeTruth.juice & ok \\
PeelOrange & peel off in few pieces, segments intact & OrangeTruth.peel, bend\_peel & ok \\\bottomrule\end{tabularx}\end{table}
The hand-juicing model gives 63~mL for orange\_0 (notes; source 60--80~mL by hand, TYPICAL). Peeling
force in newtons is a GAP; the register's strip forces are ILLUSTRATIVE.
\FloatBarrier\subsection{Potato}\label{sec:potato}
The tuber is a stem: periderm with lenticels, eyes on a spiral (denser at the rose end), cortex, vascular ring, perimedullary starch store and a branched pith. Discrete attributes exposed by the truth class: cooked\_state, greening, peel\_state, sweetened. Built instance: \texttt{potato\_\allowbreak{}0}, 126.4~g; states: whole, half\_a (63.5 g), half\_b (63.0 g), rose\_piece (56.1 g), heel\_piece (70.3 g). Research notes: \texttt{docs/wave2/POTATO.md}.
\begin{figure}[htbp]\centering\begin{minipage}[t]{0.490\linewidth}\centering\includegraphics[width=\linewidth,height=0.19\textheight,keepaspectratio]{/home/mbeetz/Documents/HouseholdTwin/report/figures/renders_food_anatomy_potato_labelled.jpg}\\[-1pt]{\scriptsize exploded anatomy plate}\end{minipage}\hfill\begin{minipage}[t]{0.490\linewidth}\centering\includegraphics[width=\linewidth,height=0.19\textheight,keepaspectratio]{/home/mbeetz/Documents/HouseholdTwin/report/figures/renders_food_item_potato.jpg}\\[-1pt]{\scriptsize whole and cut half on the worktop}\end{minipage}\caption{potato\_0: anatomy plate and path-traced render (Isaac Sim 6.2).}\end{figure}
\begin{table}[htbp]\centering\scriptsize\caption{Part tree of potato\_0 (indexed siblings merged).}
\begin{tabularx}{\linewidth}{@{}p{0.30\linewidth}p{0.27\linewidth}p{0.17\linewidth}rX@{}}\toprule Part (below /Food/Parts) & sdt:class & observability & n & mass (g) \\\midrule
\hspace*{0.0em}Periderm & \texttt{food:Periderm} & visible &  & 0.69 \\
\hspace*{0.9em}Lenticels & \texttt{food:Lenticel} & visible &  & 1,415 points \\
\hspace*{0.0em}Eyes & \texttt{food:EyeGroup} & visible &  &  \\
\hspace*{0.9em}Eye\_1..12 & \texttt{food:Eye} & visible & 12 & 0.05 \\
\hspace*{0.0em}StolonScar & \texttt{food:StolonScar} & visible &  & 0.02 \\
\hspace*{0.0em}Cortex & \texttt{food:Cortex} & visible after peeling &  & 53.75 \\
\hspace*{0.0em}VascularRing & \texttt{food:VascularRing} & visible after cutting &  & 3.24 \\
\hspace*{0.0em}PerimedullaryZone & \texttt{food:PerimedullaryZone} & visible after cutting &  & 58.14 \\
\hspace*{0.0em}Pith & \texttt{food:Pith} & visible after cutting &  & 10.53 \\\bottomrule\end{tabularx}\end{table}
\begin{table}[htbp]\centering\scriptsize\caption{Key register entries of the potato.}
\begin{tabularx}{\linewidth}{@{}p{0.19\linewidth}p{0.27\linewidth}p{0.12\linewidth}X@{}}\toprule Register key & Value / distribution & Status & Source and note \\\midrule
\texttt{cultivar} & one of Linda, Belana, Gala & STANDARD & bsl21 \\
\texttt{length} & N(85, 10) in [60, 110] mm & SOURCED-VERIFY & 75 / 65 / 99 mm (Astrix / Diamont / Santana) \\
\texttt{width} & N(55, 4) in [45, 65] mm & SOURCED-VERIFY & 48 / 45 / 58 mm \\
\texttt{eyes} & integer 8 .. 14 & SOURCED-VERIFY & Norgold 10-12 buds; 6-9 apical (mag25) \\
\texttt{phyllotaxis} & 0.4 turn & ILLUSTRATIVE & spiral arrangement sourced; the 2/5 fraction is a GAP \\
\texttt{periderm\_\allowbreak{}thickness} & N(65, 10) in [46, 84] um & SOURCED & mc20 phellem \\
\texttt{density} & N(1070, 12) in [1050, 1100] kg/m3 & SOURCED-VERIFY & 1053 (kk18); SG 1.070-1.103 \\
\texttt{E} & U(3.1 .. 5.6) MPa & SOURCED-VERIFY & potato\_mech \\
\texttt{softening} & first order, k(100 C) 0.28-0.40, Ea 103-145 1/min, kJ/mol & SOURCED & potato\_cook (Harada) \\
\texttt{boil\_\allowbreak{}times} & diced 15-20, large whole up to 60, whole 20-25 (cold start) min & SOURCED & gritzer \\
\texttt{peel\_\allowbreak{}depth} & 1.1 mm & SOURCED-VERIFY & potato\_peel ($\sim$9 \% waste) \\\bottomrule\end{tabularx}\end{table}
\begin{table}[htbp]\centering\scriptsize\caption{Task designators for the potato. Demo: symbolic run on seed 0 by this report (default parameters; a smoke test, not a validation).}
\begin{tabularx}{\linewidth}{@{}p{0.16\linewidth}Xp{0.3\linewidth}l@{}}\toprule Designator & Goal & Twin model & Demo \\\midrule
BoilPotatoes & done: a skewer passes without resistance & potato\_models.boil (conduction + first-order softening, calibrated to 2 cm cubes) & ok \\
PeelPotato & peeled, eyes and green parts removed & potato\_models.peel\_loss\_fraction & ok \\\bottomrule\end{tabularx}\end{table}
\paragraph{Model checks reported in the notes} \texttt{docs/wave2/POTATO.md} reports the boiling
model (conduction lag plus first-order softening, k100 and Ea SOURCED) calibrated to 2~cm cubes done at 17.5~min
(sources 15--20) and checked on 3~cm cubes (22~min) and a large 75~mm tuber (50~min; sources up to 60); a whole
55~mm tuber takes 34~min against 20--25 in German guides, so the model runs long there. The register's 1.1~mm peel
removes 10.0~\% of the twin's mass against about 9~\% measured -- an independent check of the geometry.
\FloatBarrier\subsection{Round / truss tomato}\label{sec:tomato}
Skin, outer pericarp wall, septa, columella with placenta, locules filled with gel round the seeds; 2--4 locules; truss tomatoes keep calyx and pedicel. Discrete attributes exposed by the truth class: bruised, cut\_state, flavour, ripeness\_class, skin. Built instance: \texttt{tomato\_\allowbreak{}15}, 71.7~g; states: whole, stem\_half (35.6 g), blossom\_half (36.0 g), half\_a (35.7 g), half\_b (35.9 g). Research notes: \texttt{docs/wave2/TOMATO.md}.
\begin{figure}[htbp]\centering\begin{minipage}[t]{0.490\linewidth}\centering\includegraphics[width=\linewidth,height=0.19\textheight,keepaspectratio]{/home/mbeetz/Documents/HouseholdTwin/report/figures/renders_food_anatomy_tomato_labelled.jpg}\\[-1pt]{\scriptsize exploded anatomy plate}\end{minipage}\hfill\begin{minipage}[t]{0.490\linewidth}\centering\includegraphics[width=\linewidth,height=0.19\textheight,keepaspectratio]{/home/mbeetz/Documents/HouseholdTwin/report/figures/renders_food_item_tomato.jpg}\\[-1pt]{\scriptsize whole and cut half on the worktop}\end{minipage}\caption{tomato\_15: anatomy plate and path-traced render (Isaac Sim 6.2).}\end{figure}
\begin{table}[htbp]\centering\scriptsize\caption{Part tree of tomato\_15 (indexed siblings merged).}
\begin{tabularx}{\linewidth}{@{}p{0.30\linewidth}p{0.27\linewidth}p{0.17\linewidth}rX@{}}\toprule Part (below /Food/Parts) & sdt:class & observability & n & mass (g) \\\midrule
\hspace*{0.0em}Skin & \texttt{food:TomatoSkin} & visible &  & 1.18 \\
\hspace*{0.0em}Pericarp & \texttt{food:Pericarp} & visible after cutting &  &  \\
\hspace*{0.9em}OuterWall & \texttt{food:OuterPericarp} & visible after cutting &  & 43.54 \\
\hspace*{0.9em}Septa & \texttt{food:SeptumGroup} & visible after cutting &  &  \\
\hspace*{1.8em}Septum\_1..3 & \texttt{food:Septum} & visible after cutting & 3 & 2.60 \\
\hspace*{0.9em}Columella & \texttt{food:Columella} & visible after cutting &  & 3.23 \\
\hspace*{0.0em}Locules & \texttt{food:LoculeGroup} & visible after cutting &  &  \\
\hspace*{0.9em}Locule\_1..3 & \texttt{food:Locule} & visible after cutting & 3 &  \\
\hspace*{1.8em}Placenta & \texttt{food:Placenta} & visible after cutting & 3 & 4.82 \\
\hspace*{1.8em}Gel & \texttt{food:LocularGel} & visible after cutting & 3 & 15.73 \\
\hspace*{1.8em}Seeds & \texttt{food:SeedGroup} & visible after cutting & 3 &  \\
\hspace*{2.7em}Seed\_1..90 & \texttt{food:TomatoSeed} & visible after cutting & 90 & 0.34 \\
\hspace*{0.0em}Calyx & \texttt{food:Calyx} & visible &  &  \\
\hspace*{0.9em}Sepal\_1..5 & \texttt{food:Sepal} & visible & 5 & 0.11 \\
\hspace*{0.0em}StemScar & \texttt{food:StemScar} & visible &  & 0.03 \\
\hspace*{0.0em}Pedicel & \texttt{food:Pedicel} & visible &  & 0.08 \\\bottomrule\end{tabularx}\end{table}
\begin{table}[htbp]\centering\scriptsize\caption{Key register entries of the round / truss tomato.}
\begin{tabularx}{\linewidth}{@{}p{0.19\linewidth}p{0.27\linewidth}p{0.12\linewidth}X@{}}\toprule Register key & Value / distribution & Status & Source and note \\\midrule
\texttt{diameter} & N(63, 4) in [52, 76] mm & SOURCED & fdc\_tomato 61-66 mm for 91-123 g \\
\texttt{locules} & one of 2, 3, 4 & SOURCED & ucd\_rost: small types 2, globe 3-5 \\
\texttt{pericarp\_\allowbreak{}thickness} & N(6.5, 0.8) in [5, 9] mm & SOURCED-VERIFY & meta19 5.3 mm (small); 7-12 germplasm \\
\texttt{plate\_\allowbreak{}firmness} & initial G/P/R 15.1 / 6.0 / 4.5; average 13.2 / 4.7 / 3.1 N/mm & SOURCED-VERIFY & ince16 \\
\texttt{skin\_\allowbreak{}rupture\_\allowbreak{}2mm} & G/P/R 6.92 / 3.99 / 2.74 N & SOURCED-VERIFY & ince16 \\
\texttt{skin\_\allowbreak{}E} & 252-293 (red); failure $\le$ 25-28 MPa at 9-10 \% MPa & SOURCED & bn05 \\
\texttt{damage\_\allowbreak{}strain} & 0.03 & SOURCED-VERIFY & tomato\_mech (Li 2017 FEM) \\
\texttt{skin\_\allowbreak{}hue\_\allowbreak{}by\_\allowbreak{}stage} & [118, 110, 96, 71, 52, 43] deg & DERIVED & h = atan2(1, a*/b*) at the band centres of ince16; ripe $\sim$40 deg (Radzevicius 2009, SV) \\
\texttt{chilling\_\allowbreak{}threshold} & 12.5 C & SOURCED & wflo\_tomato 12.5-12.8 \\
\texttt{flesh\_\allowbreak{}density} & 1000 kg/m3 & ILLUSTRATIVE & GAP (cherry 1194, suspect) \\\bottomrule\end{tabularx}\end{table}
\begin{table}[htbp]\centering\scriptsize\caption{Task designators for the round / truss tomato. Demo: symbolic run on seed 0 by this report (default parameters; a smoke test, not a validation).}
\begin{tabularx}{\linewidth}{@{}p{0.16\linewidth}Xp{0.3\linewidth}l@{}}\toprule Designator & Goal & Twin model & Demo \\\midrule
SliceTomato & slices 6-13 mm, not squashed, gel kept & TomatoTruth.cut & ok \\
PeelTomato & skin off, flesh firm & TomatoTruth.blanch & ok \\\bottomrule\end{tabularx}\end{table}
\paragraph{Bench: plate firmness (Isaac Sim FEM)} \texttt{scripts/sim\_tomato\_squeeze.py} is an
independent consistency check of the register: lanes take the whole-fruit Hertz modulus of one study (tomato\_mech,
0.58~MPa small red .. 1.97~MPa small green) and compare the FEM plate stiffness with Hertz theory and with the plate
firmness measured by another study (ince16). Analytically the red fruit matches; the green one comes out lower than
measured. The FEM converges for the soft lane only at 4000~Hz / 250 iterations; at 1000~Hz all lanes gave about
1.8~N/mm whatever the modulus.
\begin{table}[htbp]\centering\scriptsize\caption{Tomato plate stiffness 4--20~N in N/mm (test\_reports/tomato\_squeeze*.json). Measured values ince16, SOURCED-VERIFY.}
\begin{tabularx}{\linewidth}{@{}rrrrrX@{}}\toprule E (MPa) & FEM 1000 Hz & FEM 4000 Hz & Hertz & FEM/Hertz & measured initial / average \\\midrule
0.58 & 1.76 & 4.10 & 4.53 & 0.90 & red: 4.5 / 3.1 \\
1 & 1.81 & 5.23 & 6.51 & 0.80 & -- \\
1.4 & 1.76 & 5.63 & 8.15 & 0.69 & -- \\
1.97 & 1.79 & 6.00 & 10.23 & 0.59 & green: 15.1 / 13.2 \\\bottomrule\end{tabularx}\end{table}
\begin{figure}[htbp]\centering\includegraphics[width=\linewidth,height=0.25\textheight,keepaspectratio]{/home/mbeetz/Documents/HouseholdTwin/report/figures/squeeze_validation.pdf}\caption{Squeeze benches compared by this report: (a) tomato FEM at two solver settings vs Hertz and the measured plate firmness; (b) avocado FEM vs the project's Hertz model and an unrecorded-settings probe run; (c) roll FEM at two solver settings vs the register's linear model. Common finding: soft lanes converge, stiff lanes need a higher solver rate than used.}\end{figure}
\clearpage\section{Wave 3: lemon, cucumber, croissant, baguette, pepper}
The five wave-3 items exist as code (\texttt{food/<item>.py} with geometry, part tree and asset spec;
\texttt{food/<item>\_models.py} with hidden models and beliefs) and as registers. On 20~October they were
added to the asset builder and built (seed~0: USD part graph, cut states, belief and truth layers), rendered on the
worktop and as exploded anatomy plates. They still have no task designators and no benches. Their module docstrings refer to research
notes in \texttt{docs/wave3/}, which \textbf{does not exist} in the repository. The part trees below come from
instantiating each item (seed~0) on the CPU.
\begin{figure}[htbp]\centering\includegraphics[width=\linewidth,height=0.24\textheight,keepaspectratio]{/home/mbeetz/Documents/HouseholdTwin/report/figures/renders_food_wave3_group.jpg}\caption{Wave~3 items (seed~0), path-traced: lemon, cucumber, croissant, baguette and pepper.}\end{figure}
\begin{table}[htbp]\centering\scriptsize\caption{Wave-3 status (checked by this report; cache in report/cache/wave3\_parts.json).}
\begin{tabularx}{\linewidth}{@{}p{0.16\linewidth}lrrlp{0.2\linewidth}X@{}}\toprule Item & Module & Reg. & \% ILL. & instantiates & USD states & images \\\midrule
Lemon & \texttt{food/lemon.py} & 31 & 29 & yes, 43 parts & stem\_half, blossom\_half, half\_a, half\_b & plate + render \\
Greenhouse cucumber & \texttt{food/cucumber.py} & 28 & 43 & yes, 8 parts & half\_a, half\_b, rounds & plate + render \\
Butter croissant & \texttt{food/croissant.py} & 21 & 62 & yes, 3 parts & top\_half, bottom\_half, half\_a, half\_b & plate + render \\
Baguette & \texttt{food/baguette.py} & 42 & 45 & yes, 9 parts & half\_a, half\_b, top\_half, bottom\_half, crostini & plate + render \\
Bell pepper & \texttt{food/pepper.py} & 27 & 37 & yes, 12 parts & half\_a, half\_b, top\_piece, bottom\_piece & plate + render \\\bottomrule\end{tabularx}\end{table}
\FloatBarrier\subsection{Lemon}
A hesperidium like the orange (the class derives from it) with a prolate outline, the stylar nipple at the blossom end and a short neck at the stem end; thicker, paler albedo.
\begin{figure}[htbp]\centering\begin{minipage}[t]{0.490\linewidth}\centering\includegraphics[width=\linewidth,height=0.19\textheight,keepaspectratio]{/home/mbeetz/Documents/HouseholdTwin/report/figures/renders_food_anatomy_lemon_labelled.jpg}\\[-1pt]{\scriptsize exploded anatomy plate}\end{minipage}\hfill\begin{minipage}[t]{0.490\linewidth}\centering\includegraphics[width=\linewidth,height=0.19\textheight,keepaspectratio]{/home/mbeetz/Documents/HouseholdTwin/report/figures/renders_food_item_lemon.jpg}\\[-1pt]{\scriptsize whole and cut half on the worktop}\end{minipage}\caption{lemon\_0: anatomy plate and path-traced render (Isaac Sim 6.2).}\end{figure}
\begin{table}[htbp]\centering\scriptsize\caption{Part tree of lemon seed 0 from \texttt{parts()}.}
\begin{tabularx}{\linewidth}{@{}p{0.32\linewidth}p{0.3\linewidth}Xr@{}}\toprule Part (below /Food/Parts) & class & observability & n \\\midrule
\hspace*{0.0em}Peel & \texttt{food:CitrusPeel} & visible &  \\
\hspace*{0.9em}Flavedo & \texttt{food:Flavedo} & visible &  \\
\hspace*{0.9em}Albedo & \texttt{food:Albedo} & visible after peeling &  \\
\hspace*{0.9em}OilGlands & \texttt{food:OilGland} & visible &  \\
\hspace*{0.0em}StemButton & \texttt{food:Calyx} & visible &  \\
\hspace*{0.0em}Core & \texttt{food:Columella} & visible after cutting &  \\
\hspace*{0.0em}Segments & \texttt{food:SegmentGroup} & visible after peeling &  \\
\hspace*{0.9em}Segment\_1..10 & \texttt{food:Segment} & visible after peeling & 10 \\
\hspace*{1.8em}Membrane & \texttt{food:SegmentMembrane} & visible after peeling & 10 \\
\hspace*{1.8em}Vesicles & \texttt{food:JuiceVesicles} & visible after cutting & 10 \\
\hspace*{1.8em}Seed\_1..6 & \texttt{food:CitrusSeed} & visible after cutting & 6 \\\bottomrule\end{tabularx}\end{table}
\begin{table}[htbp]\centering\scriptsize\caption{Key register entries of the lemon.}
\begin{tabularx}{\linewidth}{@{}p{0.19\linewidth}p{0.27\linewidth}p{0.12\linewidth}X@{}}\toprule Register key & Value / distribution & Status & Source and note \\\midrule
\texttt{diameter} & N(61, 3) in [55, 67] mm & SOURCED & agu20 55-67 mm \\
\texttt{shape\_\allowbreak{}index} & U(1.3 .. 1.55) & SOURCED & length / diameter from agu20 ED/PD 0.6-0.8 (body without nipple) \\
\texttt{peel\_\allowbreak{}thickness} & N(5.6, 0.5) in [4.9, 7.1] mm & SOURCED & agu20, jen22\_lemon \\
\texttt{juice\_\allowbreak{}fraction} & U(30 .. 45) \% & SOURCED & agu20 31.7-44.3 (EU min 20) \\
\texttt{puncture\_\allowbreak{}force} & 41.2 N & SOURCED & ben26 5 mm probe \\
\texttt{cutting\_\allowbreak{}force} & 36.3 N & SOURCED & ben26 craft knife \\
\texttt{compression\_\allowbreak{}force} & 165.8 at 16 mm N & SOURCED & ben26 \\
\texttt{nipple\_\allowbreak{}height} & U(3 .. 7) mm & ILLUSTRATIVE & stylar mammilla; GAP \\\bottomrule\end{tabularx}\end{table}
\FloatBarrier\subsection{Greenhouse cucumber}
A pepo of three fused carpels: skin, firm mesocarp, a seed cavity with three placentae carrying immature seeds round a small core; bent onto a gentle arc.
\begin{figure}[htbp]\centering\begin{minipage}[t]{0.490\linewidth}\centering\includegraphics[width=\linewidth,height=0.19\textheight,keepaspectratio]{/home/mbeetz/Documents/HouseholdTwin/report/figures/renders_food_anatomy_cucumber_labelled.jpg}\\[-1pt]{\scriptsize exploded anatomy plate}\end{minipage}\hfill\begin{minipage}[t]{0.490\linewidth}\centering\includegraphics[width=\linewidth,height=0.19\textheight,keepaspectratio]{/home/mbeetz/Documents/HouseholdTwin/report/figures/renders_food_item_cucumber.jpg}\\[-1pt]{\scriptsize whole and cut half on the worktop}\end{minipage}\caption{cucumber\_0: anatomy plate and path-traced render (Isaac Sim 6.2).}\end{figure}
\begin{table}[htbp]\centering\scriptsize\caption{Part tree of cucumber seed 0 from \texttt{parts()}.}
\begin{tabularx}{\linewidth}{@{}p{0.32\linewidth}p{0.3\linewidth}Xr@{}}\toprule Part (below /Food/Parts) & class & observability & n \\\midrule
\hspace*{0.0em}Skin & \texttt{food:CucumberSkin} & visible &  \\
\hspace*{0.0em}Flesh & \texttt{food:Mesocarp} & visible after cutting &  \\
\hspace*{0.0em}SeedCavity & \texttt{food:SeedCavity} & visible after cutting &  \\
\hspace*{0.9em}Placenta\_1..3 & \texttt{food:Placenta} & visible after cutting & 3 \\
\hspace*{0.9em}Core & \texttt{food:CentralCore} & visible after cutting &  \\
\hspace*{0.9em}Seeds & \texttt{food:CucumberSeed} & visible after cutting &  \\\bottomrule\end{tabularx}\end{table}
\begin{table}[htbp]\centering\scriptsize\caption{Key register entries of the greenhouse cucumber.}
\begin{tabularx}{\linewidth}{@{}p{0.19\linewidth}p{0.27\linewidth}p{0.12\linewidth}X@{}}\toprule Register key & Value / distribution & Status & Source and note \\\midrule
\texttt{length} & N(320, 25) in [250, 380] mm & SOURCED & unece\_cuc 25-30 cm minimum; sakata11 \\
\texttt{diameter} & N(42, 3) in [36, 50] mm & TYPICAL & derived: 400 g / 32 cm / 950 kg/m3 $\to$ $\sim$41 mm \\
\texttt{skin\_\allowbreak{}thickness} & U(1 .. 1.6) mm & SOURCED-VERIFY & cuc\_misc peelable exocarp 1-2 mm \\
\texttt{cavity\_\allowbreak{}fraction} & U(0.45 .. 0.55) & ILLUSTRATIVE & seed cavity / fruit radius; GAP \\
\texttt{puncture\_\allowbreak{}3mm} & flesh 8.7, placenta 2.8, skin 9.7 N & SOURCED & sakata11 \\
\texttt{E\_\allowbreak{}apparent} & 0.85-1.39 (slice) | 2.5 (DiSECt) MPa & SOURCED-VERIFY & jah16; disect\_cuc \\
\texttt{firmness\_\allowbreak{}loss\_\allowbreak{}9d} & 49 \% & SOURCED-VERIFY & jah16 \\
\texttt{chilling} & < 10 C after 1-3 d C, d & SOURCED & ucd\_cuc \\\bottomrule\end{tabularx}\end{table}
\FloatBarrier\subsection{Butter croissant}
A rolled triangle of laminated dough: a glazed crust, a laminated crumb whose cut faces show the spiral of dough sheets and voids, and loose flakes.
\begin{figure}[htbp]\centering\begin{minipage}[t]{0.490\linewidth}\centering\includegraphics[width=\linewidth,height=0.19\textheight,keepaspectratio]{/home/mbeetz/Documents/HouseholdTwin/report/figures/renders_food_anatomy_croissant_labelled.jpg}\\[-1pt]{\scriptsize exploded anatomy plate}\end{minipage}\hfill\begin{minipage}[t]{0.490\linewidth}\centering\includegraphics[width=\linewidth,height=0.19\textheight,keepaspectratio]{/home/mbeetz/Documents/HouseholdTwin/report/figures/renders_food_item_croissant.jpg}\\[-1pt]{\scriptsize whole and cut half on the worktop}\end{minipage}\caption{croissant\_0: anatomy plate and path-traced render (Isaac Sim 6.2).}\end{figure}
\begin{table}[htbp]\centering\scriptsize\caption{Part tree of croissant seed 0 from \texttt{parts()}.}
\begin{tabularx}{\linewidth}{@{}p{0.32\linewidth}p{0.3\linewidth}Xr@{}}\toprule Part (below /Food/Parts) & class & observability & n \\\midrule
\hspace*{0.0em}Crust & \texttt{food:GlazedCrust} & visible &  \\
\hspace*{0.0em}LaminatedCrumb & \texttt{food:LaminatedCrumb} & visible after cutting &  \\
\hspace*{0.0em}Flakes & \texttt{food:Flake} & visible &  \\\bottomrule\end{tabularx}\end{table}
\begin{table}[htbp]\centering\scriptsize\caption{Key register entries of the butter croissant.}
\begin{tabularx}{\linewidth}{@{}p{0.19\linewidth}p{0.27\linewidth}p{0.12\linewidth}X@{}}\toprule Register key & Value / distribution & Status & Source and note \\\midrule
\texttt{length} & N(135, 8) in [115, 160] mm & ILLUSTRATIVE & finished dimensions GAP; chosen for $\sim$180 cm3 = 57 g x 3.13 cm3/g (ste25, USDA medium) \\
\texttt{mass\_\allowbreak{}check} & 42 / 57 / 67 (small / medium / large) g & SOURCED & croissant\_misc (USDA) \\
\texttt{rolling\_\allowbreak{}turns} & U(3.5 .. 4) & SOURCED & croissant\_misc \\
\texttt{butter\_\allowbreak{}layers} & 27 & SOURCED & ste25: 3 single folds \\
\texttt{crumb\_\allowbreak{}density} & 300 kg/m3 & DERIVED & 1 / 3.13 cm3/g (ste25) = 0.32 g/cm3 whole \\
\texttt{crust\_\allowbreak{}thickness} & U(0.8 .. 1.4) mm & ILLUSTRATIVE & GAP \\
\texttt{firmness} & 5.16 (1 h) $\to$ 19.2 (24 h), TPA 25 mm 50 \% N & SOURCED & ste25 \\
\texttt{flakes} & integer 40 .. 120 & ILLUSTRATIVE & loose flakes; GAP \\\bottomrule\end{tabularx}\end{table}
\FloatBarrier\subsection{Baguette}
The roll's baked-foam model on a long pointed footprint with five to seven overlapping diagonal scores (grignes) and a thicker crust.
\begin{figure}[htbp]\centering\begin{minipage}[t]{0.490\linewidth}\centering\includegraphics[width=\linewidth,height=0.19\textheight,keepaspectratio]{/home/mbeetz/Documents/HouseholdTwin/report/figures/renders_food_anatomy_baguette_labelled.jpg}\\[-1pt]{\scriptsize exploded anatomy plate}\end{minipage}\hfill\begin{minipage}[t]{0.490\linewidth}\centering\includegraphics[width=\linewidth,height=0.19\textheight,keepaspectratio]{/home/mbeetz/Documents/HouseholdTwin/report/figures/renders_food_item_baguette.jpg}\\[-1pt]{\scriptsize whole and cut half on the worktop}\end{minipage}\caption{baguette\_0: anatomy plate and path-traced render (Isaac Sim 6.2).}\end{figure}
\begin{table}[htbp]\centering\scriptsize\caption{Part tree of baguette seed 0 from \texttt{parts()}.}
\begin{tabularx}{\linewidth}{@{}p{0.32\linewidth}p{0.3\linewidth}Xr@{}}\toprule Part (below /Food/Parts) & class & observability & n \\\midrule
\hspace*{0.0em}Crust & \texttt{food:BreadCrust} & visible &  \\
\hspace*{0.9em}UpperCrust & \texttt{food:UpperCrust} & visible &  \\
\hspace*{0.9em}ScoreCrust & \texttt{food:ScoreCrust} & visible &  \\
\hspace*{0.9em}BottomCrust & \texttt{food:BottomCrust} & visible &  \\
\hspace*{0.0em}Crumb & \texttt{food:BreadCrumb} & visible after cutting &  \\
\hspace*{0.9em}TransitionZone & \texttt{food:CrustCrumbTransition} & visible after cutting &  \\
\hspace*{0.9em}CrumbCore & \texttt{food:CrumbCore} & visible after cutting &  \\
\hspace*{0.9em}GasCells & \texttt{food:GasCell} & visible after cutting &  \\
\hspace*{0.0em}Dusting & \texttt{food:FlourDusting} & visible &  \\\bottomrule\end{tabularx}\end{table}
\begin{table}[htbp]\centering\scriptsize\caption{Key register entries of the baguette.}
\begin{tabularx}{\linewidth}{@{}p{0.19\linewidth}p{0.27\linewidth}p{0.12\linewidth}X@{}}\toprule Register key & Value / distribution & Status & Source and note \\\midrule
\texttt{length} & N(600, 30) in [540, 660] mm & SOURCED & wp\_baguette 55-65 cm \\
\texttt{width} & N(55, 4) in [45, 62] mm & SOURCED & wp\_baguette 4-6 cm \\
\texttt{n\_\allowbreak{}scores} & one of 5, 6, 7 & SOURCED & wp\_grigne: 5 classic; 5-7 typical for 55-65 cm \\
\texttt{crust\_\allowbreak{}thickness\_\allowbreak{}top} & N(1.9, 0.3) in [1.4, 3.5] mm & SOURCED & dv2012 1.7 mean (zou2024: 3-4, SV) \\
\texttt{crumb\_\allowbreak{}density} & N(200, 20) in [160, 250] kg/m3 & DERIVED & $\sim$250 g / $\sim$0.9 L bulk (dimensions, wp\_baguette) minus a denser crust \\
\texttt{crumb\_\allowbreak{}porosity} & U(0.75 .. 0.85) & ILLUSTRATIVE & open irregular crumb; tomography GAP \\
\texttt{cell\_\allowbreak{}diameter\_\allowbreak{}median} & U(1.3 .. 1.9) mm & ILLUSTRATIVE & 'irregular voids of various sizes' \\
\texttt{crust\_\allowbreak{}L} & N(52, 3) in [44, 58] L* & SOURCED-VERIFY & zou2024 powder 57.5; intact darker \\\bottomrule\end{tabularx}\end{table}
\FloatBarrier\subsection{Bell pepper}
A hollow berry: pericarp wall round an air cavity, 2--4 lobes, ribs (septa), a placenta hanging from the stem end with the seeds, peduncle and calyx.
\begin{figure}[htbp]\centering\begin{minipage}[t]{0.490\linewidth}\centering\includegraphics[width=\linewidth,height=0.19\textheight,keepaspectratio]{/home/mbeetz/Documents/HouseholdTwin/report/figures/renders_food_anatomy_pepper_labelled.jpg}\\[-1pt]{\scriptsize exploded anatomy plate}\end{minipage}\hfill\begin{minipage}[t]{0.490\linewidth}\centering\includegraphics[width=\linewidth,height=0.19\textheight,keepaspectratio]{/home/mbeetz/Documents/HouseholdTwin/report/figures/renders_food_item_pepper.jpg}\\[-1pt]{\scriptsize whole and cut half on the worktop}\end{minipage}\caption{pepper\_0: anatomy plate and path-traced render (Isaac Sim 6.2).}\end{figure}
\begin{table}[htbp]\centering\scriptsize\caption{Part tree of pepper seed 0 from \texttt{parts()}.}
\begin{tabularx}{\linewidth}{@{}p{0.32\linewidth}p{0.3\linewidth}Xr@{}}\toprule Part (below /Food/Parts) & class & observability & n \\\midrule
\hspace*{0.0em}Pericarp & \texttt{food:Pericarp} & visible &  \\
\hspace*{0.9em}Skin & \texttt{food:PepperSkin} & visible &  \\
\hspace*{0.9em}Wall & \texttt{food:PepperWall} & visible after cutting &  \\
\hspace*{0.0em}Ribs & \texttt{food:SeptumGroup} & visible after cutting &  \\
\hspace*{0.9em}Rib\_1..3 & \texttt{food:Septum} & visible after cutting & 3 \\
\hspace*{0.0em}Placenta & \texttt{food:Placenta} & visible after cutting &  \\
\hspace*{0.0em}Seeds & \texttt{food:PepperSeed} & visible after cutting &  \\
\hspace*{0.0em}Stem & \texttt{food:StemGroup} & visible &  \\
\hspace*{0.9em}Peduncle & \texttt{food:Peduncle} & visible &  \\
\hspace*{0.9em}Calyx & \texttt{food:Calyx} & visible &  \\\bottomrule\end{tabularx}\end{table}
\begin{table}[htbp]\centering\scriptsize\caption{Key register entries of the bell pepper.}
\begin{tabularx}{\linewidth}{@{}p{0.19\linewidth}p{0.27\linewidth}p{0.12\linewidth}X@{}}\toprule Register key & Value / distribution & Status & Source and note \\\midrule
\texttt{diameter} & N(78, 4) in [68, 90] mm & SOURCED & he22 77.6 mm \\
\texttt{length} & N(83, 6) in [70, 100] mm & SOURCED & he22 81.3-85.3 mm \\
\texttt{lobes} & one of 3, 4 & SOURCED & sha26: 2-4, mostly 3 \\
\texttt{wall\_\allowbreak{}thickness} & N(7.2, 0.6) in [5.5, 8.5] mm & SOURCED & he22 7.5 red / 6.9 orange \\
\texttt{seeds} & integer 40 .. 140 & SOURCED-VERIFY & pepper\_misc 38-142 \\
\texttt{rib\_\allowbreak{}depth} & U(6 .. 12) mm & ILLUSTRATIVE & GAP \\
\texttt{firmness\_\allowbreak{}whole} & 43.5 N & SOURCED & pepper\_misc (8 mm probe) \\
\texttt{weight\_\allowbreak{}loss} & 2.9-4.8 at 17 C / 55 \% RH; shrivel from 6-8 \% loss g/d & SOURCED & alt20 \\
\texttt{red\_\allowbreak{}LCh} & (38.8, 41.9, 41.6) L*,C*,h & SOURCED & bar20 \\\bottomrule\end{tabularx}\end{table}
\clearpage\section[An avocado as the robot's belief state]{An avocado as the robot's belief state: from perception to action}\label{sec:avocado-reasoning}
The sections before describe food items as part graphs with beliefs and truths. This section follows one item, the
Hass avocado, through the whole loop a robot needs: perceive it, keep a probabilistic belief about its hidden
properties, ask and answer questions about it, predict what an action will do, act, and check the outcome against the
prediction. Everything here is code in this repository and runs on the twin (\texttt{avocado\_4} as the hidden truth);
nothing is validated on a real fruit or a real robot.
\begin{figure}[htbp]\centering\includegraphics[width=\linewidth,height=0.26\textheight,keepaspectratio]{/home/mbeetz/Documents/HouseholdTwin/report/figures/avocado_pipeline.pdf}\caption{The loop of this section. The SDT node (USD part graph, particle belief, discrete state) is the robot's belief state about the avocado; perception updates it by Bayes' rule, reasoning queries it, action outcomes are compared with what the node predicted and fed back.}\end{figure}
\subsection{From a photo to an SDT node}
\texttt{scripts/avocado\_from\_image.py} (\texttt{food/avocado\_vision.py}) estimates an avocado's node from one
photo: colour clustering picks the fruit-like blob, GrabCut refines it and its convex hull restores the shaded flank;
the white balance comes from the fruit's own specular highlights (dichromatic reflection: highlights carry the
illuminant's colour); the silhouette gives the aspect ratio and, by fitting the twin's meridian, the neck taper; the
skin's CIELAB lightness and chroma enter a ripeness likelihood (hue only as green / not green, because the register's
hue chart for ripe Hass is doubtful); skin relief and the stem button are read off as well. Every estimate is
labelled OBSERVED, ESTIMATED or PRIOR, and the observations go into the belief of the next subsection, from whose
posterior the instance is then built.
\begin{table}[htbp]\centering\scriptsize\caption{Estimates from the burlap photo (assets/food/avocado\_photo\_burlap/observation.json).}
\begin{tabularx}{\linewidth}{@{}p{0.17\linewidth}rrllX@{}}\toprule Estimate & value & sd & unit & status & method \\\midrule
\texttt{aspect\_\allowbreak{}L\_\allowbreak{}over\_\allowbreak{}D} & 1.56 & 0.047 & 1 & OBSERVED & silhouette length / maximum width along the principal axis (side view assumed) \\
\texttt{neck} & 0.025 & 0.03 & 1 & OBSERVED & least-squares fit of the twin's meridian to the silhouette profile (rms 0.029) \\
\texttt{length\_\allowbreak{}L} & 98 & 6 & mm & PRIOR & register prior (no scale reference in the image); pass --length-mm to measure it \\
\texttt{diameter\_\allowbreak{}D} & 62.9 & 4.3 & mm & ESTIMATED & length / observed aspect ratio \\
\texttt{white\_\allowbreak{}balance\_\allowbreak{}gains\_\allowbreak{}BGR} & 1.06, 1.02, 0.93 & -- & - & OBSERVED & from 840 specular-highlight pixels on the skin (dichromatic model) \\
\texttt{skin\_\allowbreak{}Lab\_\allowbreak{}L} & 12.7 & -- & L* & OBSERVED & median of the skin (highlights / shadow out) \\
\texttt{skin\_\allowbreak{}chroma} & 7.8 & -- & C* & OBSERVED & median chroma of the skin \\
\texttt{skin\_\allowbreak{}hue} & 47.9 & -- & deg & OBSERVED & median hue angle of the skin \\
\texttt{ripeness\_\allowbreak{}stage} & 3.93 & 0.73 & 1..5 & ESTIMATED & posterior over the Avocados Australia chart: skin L*, C* against the twin's colour model, hue as green / not green (likelihoods ILLUSTRATIVE) \\
\texttt{skin\_\allowbreak{}relief\_\allowbreak{}index} & 0.596 & -- & 1 & OBSERVED & std of high-passed luminance / mean luminance on the skin \\
\texttt{pebble\_\allowbreak{}scale} & 1.2 & 0.2 & mm & ESTIMATED & half the dominant texture wavelength (6.0 px) at the image scale, clipped to the register range 1.2-2.0 mm (ILLUSTRATIVE mapping) \\
\texttt{stem\_\allowbreak{}button} & present & -- & - & OBSERVED & lighter, browner region at the narrow end near the axis \\
\texttt{stem\_\allowbreak{}button\_\allowbreak{}length} & 6.2 & 1 & mm & ESTIMATED & visible length along the axis / 2 (side view of a clipped button, ILLUSTRATIVE) \\
\texttt{pose} & standing on the blossom end & -- & - & OBSERVED & direction of the principal axis from the wide to the narrow end in the image \\
\texttt{interior} & not observable & -- & - & PRIOR & seed, flesh bands, bruising: register priors until the fruit is cut or squeezed \\\bottomrule\end{tabularx}\end{table}
\begin{figure}[htbp]\centering\begin{minipage}[t]{0.490\linewidth}\centering\includegraphics[width=\linewidth,height=0.24\textheight,keepaspectratio]{/home/mbeetz/Documents/HouseholdTwin/report/figures/renders_vision_avocado_photo_burlap_analysis.jpg}\\[-1pt]{\scriptsize analysis sheet of a real photo}\end{minipage}\hfill\begin{minipage}[t]{0.490\linewidth}\centering\includegraphics[width=\linewidth,height=0.24\textheight,keepaspectratio]{/home/mbeetz/Documents/HouseholdTwin/report/figures/cram_camera_avocado_4_counter.jpg}\\[-1pt]{\scriptsize the PR2's head-camera view (Isaac Sim render of avocado\_4)}\end{minipage}\caption{Photo input. Left: the analysis of a user photo (segmentation, white balance, profile fit, skin colour). Right: the image the PR2 scenario uses as its camera view, rendered from the hidden-truth instance.}\end{figure}
On the rendered camera view the shape estimates are close (aspect ratio and length within a few
per cent of the instance), but the ripeness reading is biased: the rendered skin is darker than the colour chart
expects for the true stage, and the photo moves the belief to stage 3.97~$\pm$~0.53 while the truth is 2.96. Skin colour is a weak ripeness cue
(avo\_au\_stages: firmness is the better index); the likelihood is not calibrated on renders. Section~\ref{sec:pr2}
shows how the robot recovers from this.
\subsection{The node as a probability distribution: Bayesian filtering}
The avocado node is a particle filter (\texttt{food/avocado\_belief.py}): 4000 particles, each a complete candidate
avocado drawn from the register (length, diameter, neck, skin thickness, tubercle size, seed length, width and
offset, ripeness stage and ripening rate, skin friction, stem button, bruise flag, which side the seed coat sticks
to). Derived per particle: mass (body of revolution of the twin's meridian plus the seed, calibrated on the built
twins), aspect ratio, penetrometer force and flesh modulus (register ripeness model), squeeze stiffness (FEM table of
the previous section), seed and flesh mass. \emph{Prediction} is ripening at each particle's own rate; a squeeze
above a particle's bruise threshold flips its bruise flag. \emph{Measurement models} (noise ILLUSTRATIVE unless
stated): photo shape and skin colour, stem button, calliper, scale, gripper squeeze, the knife's stop depths round the
seed, the twist outcome and torque, the spoon's lever force, the extracted seed's length, the seed coat seen in the
cavity, a bruise seen after cutting. Resampling is systematic with a regularising jitter when the effective sample
size halves. The node is written into USD as \texttt{sdt\_belief.usda}: per variable mean, sd, 5/50/95~\% quantiles,
histogram and support as attributes of the part prims, plus the observation history and the particle file.
\begin{table}[htbp]\centering\scriptsize\caption{Validation of the filter on avocado\_4 with simulated noisy observations over four days (scripts/avocado\_belief\_demo.py; truth: mass 197~g, seed 40.0~mm).}
\begin{tabularx}{\linewidth}{@{}Xrrrrr@{}}\toprule Step & t (h) & ripeness belief & truth & mass (g) & seed length (mm) \\\midrule
prior & 0 & 2.75 $\pm$ 1.04 & 2.96 & 215 $\pm$ 26 & 40.1 $\pm$ 2.4 \\
photo: shape & 0 & 2.74 $\pm$ 1.05 & 2.96 & 215 $\pm$ 26 & 40.0 $\pm$ 2.4 \\
photo: skin colour & 0 & 2.99 $\pm$ 0.80 & 2.96 & 216 $\pm$ 25 & 40.0 $\pm$ 2.4 \\
photo: stem button & 0 & 2.99 $\pm$ 0.80 & 2.96 & 216 $\pm$ 25 & 40.0 $\pm$ 2.4 \\
scale: mass & 0 & 2.97 $\pm$ 0.79 & 2.96 & 197 $\pm$ 6 & 40.1 $\pm$ 2.4 \\
photo after 48 h: skin colour & 48 & 4.07 $\pm$ 0.65 & 4.11 & 197 $\pm$ 6 & 40.1 $\pm$ 2.4 \\
squeeze at 72 h (4 N) & 72 & 4.66 $\pm$ 0.38 & 4.68 & 197 $\pm$ 6 & 40.1 $\pm$ 2.4 \\
cut at 96 h: seed length & 96 & 4.89 $\pm$ 0.21 & 5.00 & 197 $\pm$ 6 & 39.2 $\pm$ 1.3 \\\bottomrule\end{tabularx}\end{table}
\begin{figure}[htbp]\centering\begin{minipage}[t]{0.490\linewidth}\centering\includegraphics[width=\linewidth,height=0.24\textheight,keepaspectratio]{/home/mbeetz/Documents/HouseholdTwin/report/figures/assets_food_avocado_4_belief_timeline.jpg}\\[-1pt]{\scriptsize belief over time}\end{minipage}\hfill\begin{minipage}[t]{0.490\linewidth}\centering\includegraphics[width=\linewidth,height=0.24\textheight,keepaspectratio]{/home/mbeetz/Documents/HouseholdTwin/report/figures/assets_food_avocado_4_belief_posteriors.jpg}\\[-1pt]{\scriptsize posteriors}\end{minipage}\caption{The avocado node's belief while it ripens and is observed (assets/food/avocado\_4/belief/).}\end{figure}
\subsection{Watching the robot peel: a vision-based peeling monitor}
\texttt{food/peel\_monitor.py} judges a peeled avocado half (lying cut face down, one oblique RGB-D view) against the
twin's \emph{expectation}: the ideal peeled surface of this instance, registered to the observed depth (trimmed ICP
and a height refinement over $x$, $y$ and yaw). Every visible surface point is classed as peeled well, skin remaining
(pebbled texture, darker, or standing proud of the expected surface), flesh missing (deeper than expected, with a
slope-dependent threshold and hysteresis over the torn area) or not judged. From the classes come the measures
(peeled share, skin area and largest patch, missing area, lost volume, deepest gouge, coverage), a score from 0 to
100 and a verdict (success / rework / not judged) with reasons, written into the node. The test set is synthetic:
\texttt{food/avocado\_peel.py} turns peeling outcomes into USD states (skin patches kept on the flesh, gouges cut out
of it by mesh booleans, ground truth recorded) and Isaac Sim renders them with depth and segmentation, the
expectation at a perturbed pose (2--4~mm, 2--4$^\circ$) so that the monitor must register it itself. The thresholds
are calibrated on these renders only.
\begin{table}[htbp]\centering\scriptsize\caption{Peeling monitor on 24 synthetic peels (6 clean, 6 with skin patches, 6 with gouges, 6 mixed; peel\_set/eval/summary.json). IoU, precision and recall are means over the samples that contain the defect; err: mean absolute area error.}
\begin{tabularx}{\linewidth}{@{}lrrrrrrrrr@{}}\toprule Input & verdicts & reworks found & false alarms & IoU skin & P / R skin & err (mm2) & IoU missing & P / R missing & err (mm2) \\\midrule
RGB-D & 22/24 & 14/15 & 1 & 0.46 & 0.53 / 0.69 & 15 & 0.53 & 0.83 / 0.59 & 40 \\
RGB only & 20/24 & 13/15 & 2 & 0.32 & 0.35 / 0.75 & 38 & 0.00 & 0.00 / 0.03 & 220 \\\bottomrule\end{tabularx}\end{table}
Depth carries the gouges: from colour alone the missing flesh is not found at all, while the skin patches,
being darker and pebbled, are found from colour too. While building the set, three rendering faults of edited meshes
had to be fixed before the evaluation meant anything: per-corner normals that were not trimmed with the faces, texture
coordinates lost in the mesh booleans, and material subsets that became empty -- each made the renderer drop the
mesh's material, so the skin patches came out pale.
\begin{figure}[htbp]\centering\includegraphics[width=\linewidth,height=0.55\textheight,keepaspectratio]{/home/mbeetz/Documents/HouseholdTwin/report/figures/assets_food_avocado_4_peel_set_eval_panels.jpg}\caption{Peeling monitor examples (left: render; middle: ground truth; right: the monitor's classes, red skin remaining, orange flesh missing, green peeled well).}\end{figure}
\begin{figure}[htbp]\centering\includegraphics[width=\linewidth,height=0.22\textheight,keepaspectratio]{/home/mbeetz/Documents/HouseholdTwin/report/figures/assets_food_avocado_4_peel_set_eval_scatter.jpg}\caption{Skin and missing areas: monitor against ground truth, RGB-D and RGB only.}\end{figure}
\subsection{Halving and pitting: an emulator}
The emulator (\texttt{food/avocado\_emulator.py}, mechanics in \texttt{food/avocado\_process.py}) is a state machine
on the node: whole $\to$ scored / partly scored $\to$ halved / halved with torn flesh $\to$ seed out, the seed coat on
the seed or left in the cavity. Each action checks its precondition, runs quasi-static mechanics on the hidden
instance and writes the new USD states, and returns what a robot would sense (\emph{measured}) apart from the truth:
\begin{table}[htbp]\centering\scriptsize\caption{Emulator actions.}
\begin{tabularx}{\linewidth}{@{}p{0.18\linewidth}Xp{0.27\linewidth}@{}}\toprule Action & Mechanics (constants ILLUSTRATIVE) & USD state \\\midrule
\texttt{cut\_\allowbreak{}round\_\allowbreak{}seed} & puncture at the stem end; force = skin resistance $\times$ skin thickness + flesh cutting toughness(stage) $\times$ depth; the blade stops on the seed coat (stop-depth profile); a force limit can stall the blade & kerf carved from skin and flesh down to the seed coat; cut faces labelled \\
\texttt{twist\_\allowbreak{}apart} & torque = uncut flesh bridges (polar moment $\times$ flesh shear strength) + seed adhesion over half the seed; ripe fruit releases at the seed, unripe tears ($p_\text{clean}$(stage)) & two halves; torn flesh as holes in the cavity wall \\
\texttt{remove\_\allowbreak{}seed} & spoon lever (0.3 $\times$ adhesion $\times$ area), knife twist (flagged hazardous for a hand-held half, farley2019), squeeze-out (ripe only); the spoon may gouge; the seed coat stays in the cavity when it adheres to the flesh & pitted half with seed-coat sheets and gouges; the seed \\\bottomrule\end{tabularx}\end{table}
\begin{table}[htbp]\centering\scriptsize\caption{Avocado task designators (food/avocado\_tasks.py); CutRoundSeed is coraplex's CuttingAction with the technique HALVING.}
\begin{tabularx}{\linewidth}{@{}p{0.15\linewidth}Xp{0.25\linewidth}@{}}\toprule Designator & Goal & Failures (cue, recovery) \\\midrule
\texttt{CutRoundSeed} & integrity = scored: a lengthwise cut through skin and flesh all round the seed & partial\_cut, no\_entry \\
\texttt{TwistApart} & integrity = halved; the seed stays in one half & stuck, torn \\
\texttt{RemoveSeed} & seed = out; cavity clean (no seed coat, no gouge) & stuck, testa\_left, gouge \\\bottomrule\end{tabularx}\end{table}
\begin{table}[htbp]\centering\scriptsize\caption{Three emulator runs on avocado\_4 (process/summary.json; truth: stage 2.96, forced to 1.8 in the unripe run; seed 40.0~mm; seed coat adheres to the flesh). The belief starts from the register prior and is updated only from the measured part of each outcome.}
\begin{tabularx}{\linewidth}{@{}lllrrr@{}}\toprule Scenario & Step & Outcome & ripeness belief & seed (mm) & P(coat on flesh) \\\midrule
ripe & \texttt{CutRoundSeed} & scored & 2.81 $\pm$ 1.03 & 40.5 $\pm$ 2.1 & 0.49 \\
ripe & \texttt{TwistApart} & halved\_torn & 1.98 $\pm$ 0.60 & 40.6 $\pm$ 2.1 & 0.47 \\
ripe & \texttt{RemoveSeed} & pitted\_testa\_left & 2.92 $\pm$ 0.36 & 41.4 $\pm$ 0.9 & 0.97 \\
unripe & \texttt{CutRoundSeed} & partly\_scored & 2.80 $\pm$ 1.03 & 40.5 $\pm$ 2.1 & 0.49 \\
unripe & \texttt{TwistApart} & stuck & 2.80 $\pm$ 1.03 & 40.5 $\pm$ 2.1 & 0.49 \\
unripe & \texttt{CutRoundSeed} & scored & 2.80 $\pm$ 1.02 & 40.6 $\pm$ 2.0 & 0.49 \\
unripe & \texttt{TwistApart} & halved\_torn & 1.98 $\pm$ 0.59 & 40.6 $\pm$ 2.0 & 0.45 \\
unripe & \texttt{RemoveSeed} & pitted\_testa\_left & 1.91 $\pm$ 0.45 & 38.2 $\pm$ 0.8 & 0.96 \\
knife & \texttt{CutRoundSeed} & scored & 2.81 $\pm$ 1.03 & 40.5 $\pm$ 2.1 & 0.49 \\
knife & \texttt{TwistApart} & halved\_torn & 1.98 $\pm$ 0.60 & 40.6 $\pm$ 2.1 & 0.47 \\
knife & \texttt{RemoveSeed} & pitted\_testa\_left & 2.93 $\pm$ 0.36 & 39.8 $\pm$ 0.9 & 0.97 \\\bottomrule\end{tabularx}\end{table}
\begin{figure}[htbp]\centering\begin{minipage}[t]{0.327\linewidth}\centering\includegraphics[width=\linewidth,height=0.17\textheight,keepaspectratio]{/home/mbeetz/Documents/HouseholdTwin/report/figures/assets_food_avocado_4_process_ripe_render_scored.jpg}\\[-1pt]{\scriptsize cut round the seed}\end{minipage}\hfill\begin{minipage}[t]{0.327\linewidth}\centering\includegraphics[width=\linewidth,height=0.17\textheight,keepaspectratio]{/home/mbeetz/Documents/HouseholdTwin/report/figures/assets_food_avocado_4_process_ripe_render_halves.jpg}\\[-1pt]{\scriptsize halves (torn twist)}\end{minipage}\hfill\begin{minipage}[t]{0.327\linewidth}\centering\includegraphics[width=\linewidth,height=0.17\textheight,keepaspectratio]{/home/mbeetz/Documents/HouseholdTwin/report/figures/assets_food_avocado_4_process_unripe_render_pitted.jpg}\\[-1pt]{\scriptsize pitted, seed, other half}\end{minipage}\caption{Emulator states, path-traced in Isaac Sim from the USD states the emulator writes.}\end{figure}
\subsection{Competence questions and mental simulation}
With the node as its belief state, the robot answers competence questions from three kinds of knowledge
(\texttt{cram/avocado\_sdt.py}): the \emph{belief} (what is it, how ripe and when edible, where is the seed and how
deep will the blade go), the \emph{task models} (what can I do now, what am I doing and what can go wrong), and
\emph{mental simulation}: the halving and pitting mechanics run on avocados drawn from the belief. An imagined
avocado needs only the geometry the mechanics use, which takes 1~ms to build instead of 1~s for the full twin, so a
question can roll out 150 imagined avocados in seconds. Answers carry the evidence they rest on.

\paragraph{Should it squeeze first? Value of information} The goal is ``guacamole now, or store the avocado two
days''. Utility = quality of the pitted halves (penalties for torn flesh, seed coat left, gouges) $\times$ how well
the flesh mashes at preparation time (0 below stage 2.5, 1 from stage 4; ILLUSTRATIVE), waiting costs 0.05. The robot
first picks the squeeze force: the strongest of 4, 3 and 2~N whose bruise risk under the belief is below 2~\% (bruise
threshold 0.6 $\times$ penetrometer force, register). Then a preposterior analysis imagines squeeze readings for the
imagined avocados (the belief's own measurement model; sensor error grows as the force falls), updates on each and
asks whether the decision changes and by how much the expected utility rises.
{\scriptsize\begin{longtable}{@{}p{0.2\linewidth}p{0.12\linewidth}p{0.62\linewidth}@{}}\caption{The robot's competence answers in the PR2 scenario (cram/results/pr2\_run.json, recorded 2026-10-20 19:40).}\\\toprule Question & Answered from & Answer \\\midrule\endfirsthead
\caption[]{The robot's competence answers in the PR2 scenario (cram/results/pr2\_run.json, recorded 2026-10-20 19:40). (continued)}\\\toprule Question & Answered from & Answer \\\midrule\endhead
What is the object? & belief & a Hass avocado (Avocado < Fruit < Produce < Food), 217 $\pm$ 26 g, 95 x 67 mm, ripeness stage 4.0 $\pm$ 0.5 (nearly ripe); parts: exocarp\_skin: tubercles, stem\_button; mesocarp\_flesh: fibre\_strands; seed: seed\_coat (testa), cotyledons x2, embryo\_axis \\
How ripe is it, and when can it be eaten? & belief & P(ready to eat, stage $\ge$ 4.5) = 0.13; P(still hard, stage < 2) = 0.00; ready in about 0.7 days (median over the belief) \\
Where is the seed, how deep will the knife go? & belief & blade stops on the seed coat at a depth (90 \% interval) of side 13..19 mm, top 32..41 mm, bottom 15..22 mm \\
What can I do with it now? & task model & state integrity=whole, seed=in: CutRoundSeed: integrity = scored: a lengthwise cut through skin and flesh all round the seed \\
Should I squeeze it before deciding? & mental simulation & yes: without squeezing I would choose 'store 2 days, then guacamole' (expected utility 0.77; guacamole now (halve, spoon out the seed): 0.68, store 2 days, then guacamole: 0.77); a 3 N squeeze (P(bruise) 0.000) changes that choice in 8\% of the imagined readings and is worth +0.012 \\
What did the squeeze tell me? & execution monitoring & as expected: 3.40 N/mm (imagined 2.05..3.97); ripeness belief ripeness\_stage 3.97 $\pm$ 0.53; penetrometer\_N 9.32 $\pm$ 5.28 $\to$ ripeness\_stage 3.79 $\pm$ 0.52; penetrometer\_N 10.12 $\pm$ 3.89 \\
Guacamole now, or store it? & mental simulation & store 2 days, then guacamole (expected utility guacamole now: 0.64, store 2 days, then guacamole: 0.77) \\
What happens if I halve and pit it now? & mental simulation & halve now, spoon out the seed: P(success) 1.00, P(clean halves) 0.92, P(torn) 0.08, P(seed coat left) 0.47, P(gouge) 0.19, quality 0.80 | halve now, knife-twist the seed out: P(success) 1.00, P(clean halves) 0.92, P(torn) 0.08, P(seed coat left) 0.47, P(gouge) 0.00, quality 0.85, HAZARD | wait 2 days, halve, spoon: P(success) 1. ... \\
Should I rather wait? & mental simulation & waiting 2 days raises the expected quality from 0.80 to 0.86 (P(clean halves) 0.92 $\to$ 1.00) \\
How should I do it (parameters)? & mental simulation & halve now, spoon out the seed; twist torque limit 0.30 N m (imagined p95 0.20), seed removal limit 6.5 (imagined p95 4.3); keep the blade on the seed coat \\
What am I doing with it? & task model & CutRoundSeed: integrity = scored: a lengthwise cut through skin and flesh all round the seed; steps: place the fruit on its side on the board $\to$ pierce at the stem end until the blade meets the seed $\to$ keep the edge on the seed coat $\to$ turn the fruit once round under the blade $\to$ withdraw; possible failures: partial\_cut: watch f ... \\\bottomrule\end{longtable}}
\begin{table}[htbp]\centering\scriptsize\caption{Imagined outcomes of the plans the robot compared (after the photo and the squeeze).}
\begin{tabularx}{\linewidth}{@{}Xrrrrrrlr@{}}\toprule Plan & wait (h) & n & P(clean) & P(torn) & P(coat left) & P(gouge) & hazard & quality \\\midrule
halve now, spoon out the seed & 0 & 150 & 0.92 & 0.08 & 0.47 & 0.19 & no & 0.80 \\
halve now, knife-twist the seed out & 0 & 150 & 0.92 & 0.08 & 0.47 & 0.00 & yes & 0.85 \\
wait 2 days, halve, spoon & 48 & 150 & 1.00 & 0.00 & 0.43 & 0.13 & no & 0.86 \\\bottomrule\end{tabularx}\end{table}
\subsection{The PR2 in the CRAM kitchen}\label{sec:pr2}
\texttt{cram/avocado\_pr2\_demo.py} puts the node into the CRAM framework (\texttt{\textasciitilde/WORK}, imported, not
changed): the apartment kitchen of the CRAM resources with the PR2, an avocado on a 4~cm cutting board on the
counter, a knife in the PR2's right gripper and a spoon on the counter. \texttt{Avocado(Fruit)} is a semantic
annotation of the semantic digital twin that carries the robot's node; the twin's meshes are exported as coloured
part bodies (skin, three flesh bands, endocarp, seed coat) for each state (\texttt{cram/export\_meshes.py}).
The PR2 drives to the counter, looks at the board and runs coraplex's \texttt{DetectAction}; the detection creates the
node from the register prior and the head-camera image is folded in. The robot then answers the questions above,
decides whether to squeeze, squeezes with the left gripper (fingers set to the believed diameter minus force over
believed stiffness; the kinematic simulation has no force control, so the reading comes from the fruit model), and
halves and pits the avocado with its own motions (\texttt{cram/avocado\_motions.py}): the left hand holds the neck,
the knife held flat in the horizontal cut plane moves its tip once round the seed outline of the belief means
(\texttt{CutRoundSeedAction}, a coraplex \texttt{CuttingAction} with its own tool path), the knife is laid down, the
left hand pinches the lower half, the right hand twists the upper half off about the cut-plane normal and puts it
down, then takes the spoon, slides it down beside the seed and under it, levers it out and drops it on the board.
Forces and outcomes come from the mechanics run on the hidden true avocado; after each step the measured outcome is
placed in the imagined distribution (as expected; unexpected if $P<0.2$; surprising if $P<0.05$ or outside the
95~\% band), the belief is updated and the world shows the new pieces.

Making the motions reachable shaped the scene: from the first stance the PR2 could not hold the knife flat close to
its body at board height (it stands 15~cm further back now); a flat knife at counter height drove the forearm into
the counter (hence the board); the cut-plane alignment pointed the gripper back at the robot until one axis was
flipped; and the spoon scoops from the far side of the seed towards the robot, lifting up and back out of the cavity.
\begin{table}[htbp]\centering\scriptsize\caption{Execution monitoring in the PR2 scenario (cram/results/pr2\_run.json).}
\begin{tabularx}{\linewidth}{@{}p{0.13\linewidth}p{0.22\linewidth}p{0.22\linewidth}lX@{}}\toprule Step & Imagined & Measured & Verdict & Belief after \\\midrule
\texttt{SqueezeTest} & stiffness 2.05..3.97 N/mm (95 \%), mean 2.51 & stiffness 3.40 N/mm at 3 N & as expected & ripeness\_stage 3.79 $\pm$ 0.52; penetrometer\_N 10.12 $\pm$ 3.89 \\
\texttt{CutRoundSeed} & P(scored) = 1.00; F\_puncture\_N imagined 6.24..14.65 & scored, F\_puncture\_N 13.10 & as expected & ripeness\_stage 3.80 $\pm$ 0.53; seed\_length 40.45 $\pm$ 2.18; testa\_adheres\_to\_flesh 0.48 $\pm$ 0.50 \\
\texttt{TwistApart} & P(halved\_torn) = 0.07; T\_peak\_Nm imagined 0.07..0.15 & halved\_torn, T\_peak\_Nm 0.18; flesh tears and stays on the seed & unexpected & ripeness\_stage 3.00 $\pm$ 0.38; seed\_length 40.43 $\pm$ 2.09; testa\_adheres\_to\_flesh 0.44 $\pm$ 0.50 \\
\texttt{RemoveSeed} & P(pitted\_testa\_left) = 0.49; F\_peak\_N imagined 2.65..5.71 & pitted\_testa\_left, F\_peak\_N 3.37; spoon took flesh from the cavity rim & as expected & ripeness\_stage 3.16 $\pm$ 0.35; seed\_length 41.45 $\pm$ 0.85; testa\_adheres\_to\_flesh 0.96 $\pm$ 0.19 \\\bottomrule\end{tabularx}\end{table}
\begin{figure}[htbp]\centering\includegraphics[width=\linewidth,height=0.26\textheight,keepaspectratio]{/home/mbeetz/Documents/HouseholdTwin/report/figures/avocado_pr2_belief.pdf}\caption{The robot's ripeness belief through the PR2 scenario: the photo makes the avocado look riper than it is, the squeeze (a weak cue) corrects little, the plan predicts clean halves, the twist tears the flesh -- not as expected -- and the belief falls towards the truth.}\end{figure}
At the end the node holds ripeness 3.16~$\pm$~0.35 (truth 2.96), seed length 41.4~$\pm$~0.9~mm (truth 40.0~mm) and $P(\text{seed coat sticks to the flesh}) = 0.96$ (truth: flesh), from 18 answers, 3 imagined plans and 4 monitored steps. The decision itself is honest about its context: the robot would store the avocado for two days;
the scenario asks for now, so it proceeds and expects torn flesh with low probability -- which is exactly where it is
wrong, because the photo misled it.
\subsection{An interactive kitchen in CRAMERA}
\texttt{cram/avocado\_kitchen.py} keeps the same kitchen live in CRAMERA, the CRAM viewer (3D scene, EQL query panel,
plan graph), and lets an operator drive it from the query panel. Commands and questions are ready-made queries
(buttons); typed requests are matched to them:
\begin{table}[htbp]\centering\scriptsize\caption{The operator's controls in the CRAMERA query panel.}
\begin{tabularx}{\linewidth}{@{}p{0.2\linewidth}X@{}}\toprule Kind & Entries \\\midrule
commands (queued, executed with PR2 motions) & Detect the avocado; Squeeze the avocado; Wait two days; Cut round the seed; Twist the halves apart; Spoon out the seed; Reset the kitchen \\
questions answered at once & What is the object? How ripe is it? Where is the seed? What can you do now? What are you doing? \\
questions answered by mental simulation & Should you squeeze it first? Guacamole now, or store it? What happens if you halve it now? \\
views & the robot's answers, its belief, what it imagined, how the execution went, the command log, the commands with their availability \\
typed phrasings (examples) & squeeze it; halve it; open the avocado; remove the seed; is it ripe; what can you do; start over \\\bottomrule\end{tabularx}\end{table}
A command is checked against the robot's and the avocado's state and refused with the reason (``the halves
have to be apart first'', ``I have not detected an avocado yet''), otherwise queued and executed by a worker while
the viewer follows the robot. ``Wait two days'' ripens the hidden avocado and the belief; a twist after waiting
usually comes out clean. Two details of CRAMERA shaped the design: the viewer builds every preset's code to word it
in English, so a command fires only when the query is evaluated (the command sits in an attribute read by the query's
filter); and the world records every annotation it adds or removes as JSON, so the node is kept beside the
annotation, not in it. Running beside another CRAMERA session needs its own bridge port, data directory and viewer.
No screenshot of the 3D view is included: the headless browser available here has no WebGL.
\paragraph{What this section does not show} No real avocado, no real robot. The kinematic simulation has no
contacts: knife, fingers and spoon pass through the fruit, and the mechanics supply forces and outcomes. The
mechanics' constants, the utility and the photo's colour likelihood are ILLUSTRATIVE; the peeling monitor's
thresholds are calibrated on the twin's own renders. The robot's seed-centre estimate is the twin's nominal position
(no sideways offset), and torn flesh is not shown on the seed piece.
\clearpage\section[From reasoning to acting, and from acting to theories]{From reasoning to acting, and from
acting to theories}\label{sec:avocado-theory}
Section~\ref{sec:avocado-reasoning} gave the robot a belief about one avocado and the means to ask questions about it.
This section takes three further steps. First, the robot becomes an \emph{agent} that knows which effects its task
must bring about and which it must not cause, predicts both over the fruit's lifetime, and chooses, executes and
audits plans accordingly. Second, the knowledge it acts on is written down as a \emph{microtheory} -- the egg
microtheory of the AICOR working report transferred to the avocado -- and tested like a scientific theory against
held-out simulated fruit, with ablations and misspecified plants. Third, the robot \emph{learns} that theory: it
grows from a deliberately weak seed by its own failures, and reinforcement-learning problems are generated from its
successive versions. The full microtheory is a separate document, \texttt{report/avocado\_microtheory.pdf}
(15 pages); this section summarises it. Everything is simulation evidence on partly ILLUSTRATIVE laws.
\subsection{An agent that enforces and prevents effects}
\texttt{cram/avocado\_agent.py} turns the node's knowledge into decisions. A \emph{mission} names what is to be
served when (guacamole now; half now and half for tomorrow's breakfast; guacamole in two days) and bounds every
effect: a goal must hold with at least its probability, a prevented effect may occur with at most its probability,
and a bound of zero is a safety rule that no utility can trade against (Table~\ref{tab:effects}). The plan space
combines when to cut (now or at the serving time), where the whole fruit waits (counter or fridge), how the seed comes
out (spoon or knife twist), whether seed-coat sheets are lifted, and how a kept half is stored (six treatments from
nothing to vacuum-packed in the fridge). For each plan the agent runs mental simulations over 80 avocados drawn from
its belief: the lifetime model (\texttt{food/avocado\_lifetime.py}: ripening at the fruit's rate with a $Q_{10}$ of
2.5 between counter and fridge, a ripe window of 24--48~h, decline, spoilage, browning of a cut face as
$\mathrm{d}E=k\sqrt{t\,f_\text{treatment}}$) advances each one to the cut, the processing mechanics halve and pit it
with the repairs the agent would make, and the servings are scored. A geometric check (\texttt{plan\_hold}) places
the holding gripper at least 10~mm from the blade's plane, pinching the fruit's top or, if that fails, pressing it
from above.
\begin{table}[htbp]\centering\scriptsize\caption{Effects and bounds of the standard missions (cram/avocado\_agent.py).\label{tab:effects}}
\begin{tabularx}{\linewidth}{@{}lllX@{}}\toprule Effect & Kind & Bound & Meaning \\\midrule
halves\_apart & enforce & $\ge 0.95$ & the two halves are separated \\
seed\_out & enforce & $\ge 0.95$ & the seed is out of the fruit \\
cavity\_clean & enforce & $\ge 0.90$ & no seed-coat sheets left in the cavity \\
edible\_at\_service & enforce & $\ge 0.70$ & every serving mashes well and is not past its best (eating quality $\ge$ 0.6) \\
fresh\_at\_service & enforce & $\ge 0.80$ & no serving's cut face is visibly brown \\
flesh\_intact & enforce & $\ge 0.30$ & no torn flesh, no gouges \\
hazard & prevent (safety rule) & $\le 0.00$ & a knife strike into the seed of a held half (pitting injuries, farley2019) \\
blade\_near\_hand & prevent (safety rule) & $\le 0.00$ & the holding gripper within 10 mm of the blade's plane \\
bruise & prevent & $\le 0.02$ & a squeeze bruises the flesh \\
spoiled & prevent & $\le 0.05$ & the fruit spoils before it is eaten \\
browned & prevent & $\le 0.20$ & a cut face is visibly brown when served \\
testa\_left & prevent & $\le 0.10$ & seed-coat sheets stay in the cavity \\
torn & prevent & $\le 0.60$ & the twist tears the flesh (it clings to the seed) \\
gouge & prevent & $\le 0.60$ & the spoon takes flesh \\\bottomrule\end{tabularx}\end{table}
\begin{table}[htbp]\centering\scriptsize\caption{The agent's decisions for the avocado the PR2 sees on the counter (belief after the head-camera photo: ripeness\_stage 3.97 $\pm$ 0.53; seed\_length 40.03 $\pm$ 2.46; testa\_adheres\_to\_flesh 0.51 $\pm$ 0.50). Every knife-twist plan is rejected by the safety rule (P(hazard) = 1).}
\begin{tabularx}{\linewidth}{@{}p{0.17\linewidth}p{0.3\linewidth}lX@{}}\toprule Mission & Chosen plan & feasible & Why \\\midrule
Guacamole now & cut now; seed out with the spoon; peel off seed-coat sheets & 1/4 & meets every goal and stays within every risk bound; expected quality 0.75 (best of 1 feasible plans out of 4) \\
Half now, half for tomorrow's breakfast & cut now; seed out with the spoon; peel off seed-coat sheets; keep the other half: vacuum+fridge & 3/24 & meets every goal and stays within every risk bound; expected quality 0.68 (best of 3 feasible plans out of 24) \\
Guacamole in two days & keep it whole in the counter for 48 h, then cut; seed out with the spoon; peel off seed-coat sheets & 2/12 & meets every goal and stays within every risk bound; expected quality 0.89 (best of 2 feasible plans out of 12) \\\bottomrule\end{tabularx}\end{table}
\paragraph{Lifetime-based decisions} The same fruit leads to different plans for different missions. Its
lifetime under the belief: counter: edible from 0.0 d to 3.8 d, P(spoiled) > 5 \% from 4.5 d; fridge: edible from 0.0 d to 12.0 d, no spoilage risk within 12 d. For guacamole in two days the agent keeps
the fruit whole on the counter and cuts at the serving time (expected quality higher than cutting now, because the
fruit ripens into its best window). For half now and half tomorrow the kept half is what decides: only treatments
that slow browning enough keep its cut face fresh after 15~hours (Table~\ref{tab:keep}).
\begin{table}[htbp]\centering\scriptsize\caption{Half now, half tomorrow: the plans that cut now, use the spoon and lift the seed coat, by storage of the kept half.\label{tab:keep}}
\begin{tabularx}{\linewidth}{@{}llllX@{}}\toprule Kept half stored as & browning factor & P(browned) & quality & Feasible / violated bound \\\midrule
none & 1.00 & 1.00 & 0.37 & fresh\_at\_service 0.00 < 0.80; browned 1.00 > 0.20 \\
lemon & 0.45 & 1.00 & 0.46 & fresh\_at\_service 0.00 < 0.80; browned 1.00 > 0.20 \\
wrap & 0.35 & 1.00 & 0.50 & fresh\_at\_service 0.00 < 0.80; browned 1.00 > 0.20 \\
wrap+fridge & 0.12 & 0.00 & 0.60 & yes \\
lemon+wrap+fridge & 0.06 & 0.00 & 0.65 & yes \\
vacuum+fridge & 0.03 & 0.00 & 0.68 & yes \\\bottomrule\end{tabularx}\end{table}
\paragraph{Execution with a goal watch and an effect audit} In the CRAMERA kitchen
(\texttt{cram/avocado\_kitchen.py}) a mission is a command: the agent decides, explains the choice and the rejected
alternatives, and the PR2 executes the plan with the motions of Section~\ref{sec:pr2} (with new ones for holding,
lifting the seed coat with the spoon tip and squeezing lemon on the kept half). After every step a goal watch compares
the observed effects with the predicted ones and reassesses the remaining risk (e.g.\ a torn twist raises the
gouge risk of the spoon); after the mission an effect audit lists each goal and prevented effect as achieved,
violated or still pending (the kept half until breakfast). Advancing the kitchen's clock ripens, spoils and browns
the hidden fruit, so a decision to wait can be checked against what happens.
\subsection{An avocado manipulation microtheory}
The egg microtheory distinguishes grasping an egg intact from cracking it open: the same fracture is a failure in one
task and the goal in the other. For the avocado the event is the separation of the flesh: forbidden while the fruit
is squeezed to judge its ripeness (Task A), required in the cut plane and at the twist when it is halved and pitted
(Task B), and forbidden everywhere else (bruise, torn cavity wall, gouge). The microtheory
(\texttt{theory/}) has the egg's parts: a research SDT whose hidden parameters are drawn once per avocado and whose
events are patches over relational facts; reduced-order laws for squeezing, bruising, cutting, twisting, seed removal,
peeling, ripening and browning; a probabilistic hybrid model whose jump probabilities are marginals over the hidden
parameters; an observation model and a particle belief; and phase-indexed requirements over the patch trace. Its two
central predictions are a squeeze \emph{corridor} between an informative and a bruising force, which is empty for ripe
fruit, and a ripeness \emph{band in time} for the cut, between torn flesh and overripeness
(Figure~\ref{fig:mt}).
\begin{figure}[htbp]\centering\begin{minipage}[t]{0.490\linewidth}\centering\includegraphics[width=\linewidth,height=0.2\textheight,keepaspectratio]{/home/mbeetz/Documents/HouseholdTwin/report/figures/avo_mt_corridor.pdf}\\[-1pt]{\scriptsize Task A: squeeze corridor (model lines, simulation markers; green: safe and informative)}\end{minipage}\hfill\begin{minipage}[t]{0.490\linewidth}\centering\includegraphics[width=\linewidth,height=0.2\textheight,keepaspectratio]{/home/mbeetz/Documents/HouseholdTwin/report/figures/avo_mt_band.pdf}\\[-1pt]{\scriptsize Task B: success by day of the cut (simulated, model dashed)}\end{minipage}\caption{The two central predictions of the avocado microtheory, checked against held-out simulated avocados (details in avocado\_microtheory.pdf).\label{fig:mt}}\end{figure}
\begin{table}[htbp]\centering\scriptsize\caption{Held-out evaluation: the same 90 avocados (30 unripe, 30 breaking, 30 ripe) for every policy; discordant pairs and exact McNemar tests in the microtheory document.}
\begin{tabularx}{\linewidth}{@{}p{0.18\linewidth}p{0.3\linewidth}lX@{}}\toprule Task & Policy & success & Paired against (only policy : only baseline) \\\midrule
A: assess ripeness by touch, intact & closed-loop squeeze under a 1 \textbackslash{}\% bruise chance constraint & 82/90 & thumb 10 N 46:0; photo only 14:2; open-loop prior plan 3:1 \\
B: halve and pit within 5 days & belief band: daily squeeze, cut on the planned day, repairs, seed-coat lifting & 88/90 & cut now 48:0; wait 2 d 21:2; wait 4 d 20:2; photo-only band 2:2 \\
B, ablations & remove one distinction from the robot's model & -- & no seed-coat side 49:0; ready / not ready 1:2; no repair 0:0; population ripening rate 0:0 \\\bottomrule\end{tabularx}\end{table}
\paragraph{What it shows} The belief regulates better than fixed policies (the thumb press bruises every ripe
fruit; cutting at once tears unripe fruit; waiting four days lets ripe fruit go past its best). Only one removed
distinction is load-bearing on these avocados -- which side the seed coat sticks to; the others make no detectable
difference, which needs a discriminating test, not a conclusion. An error that makes the flesh weaker than believed
cannot be corrected in time (the closed loop then succeeds on 4/30 ripe fruit: the bruise is
the first evidence), while a camera bias is corrected by touch. A replay audit of 1,350 episodes (8,699 steps) finds no mismatch between the relational action axioms and the
simulator's patches; 22 pairs of steps from the same relational state reach different
branches -- the hidden attributes (ripeness, the seed-coat side) decide, so a proof needs them.
\subsection{Learning the theory from failures}
Following \emph{Manipulation as Theory Construction}, the robot's competence is a theory
$\mathcal{T}_k=\langle\Sigma_k,A_k,M_k,C_k,E_k\rangle$: concepts that are perceivable through the SDT, actionable through
designators and traceable; assertions; body-world-system families (scopes); scoped claims
$\Gamma=\mathrm{Competent}(F_\text{task},F_\text{BWS},A_\text{assess})$; and evidence (\texttt{theory/package.py}). The
seed $\mathcal{T}_0$ knows only how to halve and pit a ripe Hass avocado. Five requests expand the task or the family;
each runs the executive derived from the current theory, the failures become a localized gap, competing hypotheses are
tested by discriminating experiments in the research SDT on development avocados, a patch is fitted, and the
integrator admits, restricts or versions it; the revised claim is assessed on the held-out avocados
(\texttt{theory/learning.py}, Table~\ref{tab:revision}).
\begin{table}[htbp]\centering\scriptsize\caption{The revision path of the avocado theory (simulation evidence).\label{tab:revision}}
\begin{tabularx}{\linewidth}{@{}lp{0.25\linewidth}p{0.26\linewidth}p{0.2\linewidth}X@{}}\toprule  & Request & Gap: theory before & Patch & Held-out after \\\midrule
T0 & halve and pit a ripe Hass & -- & -- & 30/30 verified \\
T1 & Halve this firm avocado. & RipenessDependentSeparation: 10/60 & admit; +RipenessStage, TornFlesh & 88/90 verified \\
T2 & Pit it; the cavity must be clean. & SeedCoatUnrepresented: 39/90 & admit; +SeedCoatSide & 88/90 verified \\
T3 & Use the new camera (lighter images). & ColourRipenessMisspecified: 27/30 & restrict + version; +TactileRipeness & 30/30 verified \\
T4 & Keep the other half for tomorrow's breakfast. & BrowningUnrepresented: 0/90 & admit; +Browning & 88/90 verified \\
T5 & Halve and pit this other cultivar. & CultivarOutsideFamily: 41/60 & restrict & 41/60 refuted \\\bottomrule\end{tabularx}\end{table}
The patches recover the plant's laws from the experiments: the clean-release law's centre at stage 2.72 (plant 2.8; its width is overestimated on a coarse grid), the share of
fruit whose seed coat sticks to the flesh 0.53, the peel repair's success
0.82 per pass (plant 0.85) and the new camera's lightness bias +10.4 (plant +10). The last request -- a cultivar whose seed clings ten times harder -- is answered by
\emph{restriction}: the only remedy, the knife twist, is forbidden, so the claim is retracted for that family; the
evidence grows while the warranted scope contracts. At runtime a \emph{commitment gate} separates warranted
commitments from uncertain, unsupported, counterevidenced and out-of-scope ones:
\begin{table}[htbp]\centering\scriptsize\caption{The commitment gate of the final theory.}
\begin{tabularx}{\linewidth}{@{}p{0.3\linewidth}lX@{}}\toprule Situation & State & Reason \\\midrule
ripe fruit, photo + touch & warranted & claim Γ2 clean cavity verified (88/90), instance in scope (P = 0.97) \\
photo only & uncertain & the scope requires tactile ripeness evidence; only photo \\
knife twist proposed & unsupported & forbidden effect predicted: hazard P = 1.00 \\
sticky cultivar & counterevidenced & claim Γ2 on a sticky cultivar retracted \\
not an avocado? & out-of-scope & P(instance in Hass, any ripeness) = 0.30 \\
browning not represented & unsupported & concepts not operationally covered: Bruising \\\bottomrule\end{tabularx}\end{table}
\subsection{Reinforcement learning programmed by the theory}
Reinforcement learning lives inside the theory (\texttt{theory/prl.py}): the learning problem is generated from a
theory version. Its state features are the theory's perceivable concepts (believed stage and spread, day, a torn
twist, the seed coat seen), its actions the designators' parameters (cut or wait, spoon or knife twist, peel passes
only if the theory has the seed-coat concept), its reward the claim's KPIs, and its hard constraints the forbidden
effects, as a shield that masks every action whose forbidden effect the theory predicts. A tabular Monte-Carlo learner
trains on 5,600 research-SDT episodes and is evaluated greedily on the held-out avocados (Table~\ref{tab:prl}).
\begin{table}[htbp]\centering\scriptsize\caption{Programmable RL on the 90 held-out avocados of Task B.\label{tab:prl}}
\begin{tabularx}{\linewidth}{@{}Xll@{}}\toprule Program / policy & success & hazards \\\midrule
RL from T1 (no seed-coat concept), shielded & 38/90 & 0 \\
RL from T3, shielded & 88/90 & 0 \\
RL from T3, hazard as a soft penalty 0.05 & 0/90 & 90 \\
hand-coded executive of the same theory (belief band) & 88/90 & 0 \\\bottomrule\end{tabularx}\end{table}
\begin{figure}[htbp]\centering\includegraphics[width=\linewidth,height=0.2\textheight,keepaspectratio]{/home/mbeetz/Documents/HouseholdTwin/report/figures/avo_mt_rl.pdf}\caption{Left: success during training. Right: the cut / wait rule learned from $\mathcal{T}_3$ with the shield, by believed stage and day -- the ripeness band, learned.}\end{figure}
The learner inherits the theory's distinctions and its safety rules: from $\mathcal{T}_3$ with the shield it
learns the band and matches the hand-coded executive; from $\mathcal{T}_1$, which has no seed-coat concept, no amount
of training cleans a cavity whose dirt it cannot represent; and with the hazard as a small penalty instead of a shield
it buys the forbidden knife strike in every episode. What the theory does not represent the learner cannot learn, and
what the theory forbids a reward must not be allowed to buy.

\paragraph{What this section does not show} No physical validation: the plant and the robot share the laws except
where a misspecification is injected; contact is quasi-static mechanics, not contact simulation; the photo in the
microtheory is a symbolic observation from the colour law; the learning loop's agents are scripted functions, and the
RL is tabular on a handful of features. The 30 held-out avocados per class are reused across tables (paired, not
independent).
\clearpage\section{The kitchens: EASE and GARMI apartments}\label{sec:kitchen}
\texttt{scripts/import\_apartments.py} imports the EASE and GARMI apartments and the GARMI robot from the CRAM viz
scene bundles (read in place from \texttt{\textasciitilde/ResearchSpace}, which stays read-only) into USD with Isaac
Sim's URDF importer: doors and drawers become joints, collisions are convex decompositions. The GARMI URDF needed its
gazebo, ros2\_control and transmission blocks stripped before the importer would accept it.
\begin{table}[htbp]\centering\footnotesize\caption{Imported scenes, counted in the source URDFs by this report.}
\begin{tabularx}{\linewidth}{@{}Xrrr@{}}\toprule Scene & links & joints & movable joints \\\midrule
EASE (IAI) apartment & 118 & 117 & 42 \\
GARMI (TUM MIRMI) apartment & 79 & 78 & 11 \\
GARMI robot (cleaned URDF) & 64 & 63 & 26 \\\bottomrule\end{tabularx}\end{table}
\paragraph{The GARMI rig} \texttt{household/garmi\_rig.py} drives the imported GARMI (two FR3 arms with Franka
hands, a two-stage lift, a pan-tilt head, a wheeled base) as one PhysX articulation with the base fixed for
manipulation (the probe found 26 degrees of freedom and 53 links): damped-least-squares Cartesian IK per arm on the PhysX Jacobian with a null-space posture bias,
force-limited grasping (finger drives in force mode, targets ramped -- an instantly set stiff target slams the
fingers shut, a lesson from the egg bench), and the Franka TCP 0.1034~m along the hand axis.

\paragraph{Food in the kitchen} \texttt{scripts/render\_apartment\_kitchen.py} places GARMI and a cutting board on the
EASE kitchen island with seven food assets: apple\_3 whole and half, avocado\_4 whole and half, and three eggs. Only
these egg-set items have been placed; none of the wave-1, wave-2 or wave-3 items is in a kitchen yet, and the GARMI
apartment has been imported but not rendered or furnished with food.
\paragraph{End-to-end task: not yet working} \texttt{scripts/sim\_garmi\_egg.py} is the first
end-to-end task (GARMI picks an egg, cracks it on the bowl rim and opens it with both hands). Its last run
(\texttt{renders/garmi\_egg.log}) stopped during the approach: the arm's joint state and TCP became NaN and the
controller raised \emph{cannot convert float NaN to integer}. No GARMI task log exists yet.
\begin{figure}[htbp]\centering\includegraphics[width=\linewidth,height=0.36\textheight,keepaspectratio]{/home/mbeetz/Documents/HouseholdTwin/report/figures/renders_apartment_ease_kitchen_overview.jpg}\caption{The EASE apartment kitchen with GARMI (existing render, path-traced in Isaac Sim~6.2).}\end{figure}
\begin{figure}[htbp]\centering\begin{minipage}[t]{0.490\linewidth}\centering\includegraphics[width=\linewidth,height=0.2\textheight,keepaspectratio]{/home/mbeetz/Documents/HouseholdTwin/report/figures/renders_apartment_ease_kitchen_island.jpg}\\[-1pt]{\scriptsize island with board and food}\end{minipage}\hfill\begin{minipage}[t]{0.490\linewidth}\centering\includegraphics[width=\linewidth,height=0.2\textheight,keepaspectratio]{/home/mbeetz/Documents/HouseholdTwin/report/figures/renders_apartment_ease_island_food.jpg}\\[-1pt]{\scriptsize close-up: apple, avocado, eggs}\end{minipage}\caption{Food twins on the EASE kitchen island (existing renders).}\end{figure}
\clearpage\section{Data quality}\label{sec:quality}
\begin{figure}[htbp]\centering\includegraphics[width=\linewidth,height=0.4\textheight,keepaspectratio]{/home/mbeetz/Documents/HouseholdTwin/report/figures/register_status.pdf}\caption{Register entries by status, per item and overall (drawn by this report from food/params.py). Numbers in brackets: entries per item.}\end{figure}
Of 564 register entries, 180 (32~\%) are SOURCED, 108 (19~\%) SOURCED-VERIFY, 14 STANDARD, 12 DERIVED, 3 CALIBRATED, 23 TYPICAL and 224 (40~\%) ILLUSTRATIVE. The ILLUSTRATIVE values concentrate where sources end: colours (L*C*h of flesh and skin), internal dimensions (membrane, rib and core sizes), friction coefficients, the peeling and cutting forces, and the outcome constants of the strike and peel models. The sampled instance state (age, storage) is ILLUSTRATIVE by design: it is the scenario, not a property of food.
\begin{table}[htbp]\centering\scriptsize\caption{Register status counts per item (food/params.py).}
\begin{tabularx}{\linewidth}{@{}lrrrrrrrr@{}}\toprule Item & STANDARD & SOURCED & CALIBRATED & DERIVED & SOURCED-VERIFY & TYPICAL & ILLUSTRATIVE & total \\\midrule
egg & 0 & 16 & 2 & 1 & 2 & 3 & 12 & 36 \\
apple & 0 & 19 & 0 & 0 & 4 & 2 & 8 & 33 \\
avocado & 1 & 14 & 0 & 2 & 7 & 2 & 22 & 48 \\
roll & 1 & 10 & 1 & 4 & 8 & 3 & 19 & 46 \\
carrot & 2 & 6 & 0 & 0 & 17 & 1 & 16 & 42 \\
banana & 1 & 16 & 0 & 0 & 8 & 1 & 15 & 41 \\
orange & 0 & 14 & 0 & 1 & 9 & 2 & 12 & 38 \\
loaf & 0 & 8 & 0 & 0 & 6 & 1 & 11 & 26 \\
potato & 1 & 7 & 0 & 0 & 14 & 0 & 14 & 36 \\
tomato & 2 & 9 & 0 & 1 & 8 & 0 & 17 & 37 \\
onion & 1 & 6 & 0 & 0 & 9 & 1 & 15 & 32 \\
lemon & 0 & 19 & 0 & 1 & 1 & 1 & 9 & 31 \\
cucumber & 1 & 7 & 0 & 0 & 7 & 1 & 12 & 28 \\
croissant & 1 & 6 & 0 & 1 & 0 & 0 & 13 & 21 \\
baguette & 2 & 9 & 0 & 1 & 7 & 4 & 19 & 42 \\
pepper & 1 & 14 & 0 & 0 & 1 & 1 & 10 & 27 \\
\textbf{all} & \textbf{14} & \textbf{180} & \textbf{3} & \textbf{12} & \textbf{108} & \textbf{23} & \textbf{224} & \textbf{564} \\\bottomrule\end{tabularx}\end{table}
\paragraph{Consistency checks run by this report} The checks below are computed when the report is built;
they test the repository, not the food.
\begin{table}[htbp]\centering\footnotesize\caption{Repository consistency checks.}
\begin{tabularx}{\linewidth}{@{}Xp{0.45\linewidth}@{}}\toprule Check & Result \\\midrule
part masses sum to the whole & 24/27 instances within 0.5 \%; roll and loaf +1.4 \% \\
egg crack study reproducible with the current register & yes \\
SOURCED / SOURCED-VERIFY / CALIBRATED entries naming a source key of params.SOURCES & 264/291 \\
source keys in params.SOURCES that the register cites & 166/205 \\
egg assets built with the current egg register & yes \\
egg\_1 truth.usda impact\_effective\_mass\_fraction vs register & 0.53 in the asset vs 0.53 (CALIBRATED) in the register \\
docs/ANATOMY.md up to date with the register & yes \\
research notes for wave 3 (docs/wave3) & missing \\
wave-1 and wave-2 designator demos (seed 0) & 20/20 reach their goal \\\bottomrule\end{tabularx}\end{table}
Further known inconsistencies: \texttt{docs/wave1/ROLL.md} gives roll\_0 as 48.1~g (the part sum) while the
asset's whole body is 47.4~g; the \texttt{food/dynamics.py} docstring still calls both impact constants ILLUSTRATIVE
although one is now CALIBRATED; the drop-surface stiffnesses and the strike-model constants live in code
(\texttt{dynamics.SURFACES}, \texttt{egg\_crack.P}) and are mirrored only as a text note in the register.
\section{Limitations and next steps}\label{sec:next}
\paragraph{What is validated, and how far}
\begin{itemize}
\item \emph{Against published measurements, independently:} the tomato Hertz stiffness from one study's modulus
matches another study's measured red-fruit plate firmness (the green fruit does not); the potato peel loss (10.0 vs
about 9~\%); part and cut masses against the register's mass checks.
\item \emph{Fits:} the egg boiling thresholds, the egg impact constant, the roll crust-water curve and the carrot boiling
thresholds are CALIBRATED -- they reproduce the data they were fitted to.
\item \emph{Solver-limited:} the FEM squeezes converge only for soft lanes (ripe avocado, fresh roll, red tomato at
4000~Hz); stiff lanes are lower bounds.
\item \emph{Mechanism only:} banana peeling, carrot planking and the egg crack study show plausible trends, but their
forces or outcome constants are ILLUSTRATIVE.
\item \emph{Not validated at all:} appearance (no comparison with reference photographs yet), wave-3 behaviour, any
real robot run.
\end{itemize}
\paragraph{Gaps in the build}
\begin{itemize}
\item \textbf{Anatomy review gate.} The exploded plates were sent for the user's review; the gate is still pending for
every item, so all part trees are provisional.
\item \textbf{Wave~3 incomplete:} built and rendered (20~October), but the \texttt{docs/wave3/} research
notes, task designators, schema entries and benches are missing.
\item \textbf{Wave~2} has no schema entries; only the tomato has a bench.
\item \textbf{Apple without designators or bench.} The avocado has designators, an emulator and the PR2 scenario
(kinematic: no contact forces; Section~\ref{sec:avocado-reasoning}).
\item \textbf{Kitchens:} only the egg set has been placed (EASE); the GARMI apartment is imported but unused; the
GARMI egg task fails with NaN in the arm controller.
\item \textbf{Bodies are simplified in physics:} homogeneous FEM bodies (no skin, seed or crust), rigid peel segments,
kinematic shell fragments, a rigid yolk.
\item \textbf{Stale artefacts:} rebuild the egg assets and \texttt{docs/ANATOMY.md} with the current register; fix the
bread part-mass overlap.
\end{itemize}
\paragraph{Next steps} (1) Pass the anatomy review; (2) rerun the stiff FEM lanes at higher solver rates (the
memory note: 4000~Hz / 250 iterations was enough for soft lanes only) or replace them by the analytic models; (3) fix
the GARMI IK failure and log the first end-to-end egg task in the EASE kitchen; (4) build wave~3 and render wave~2;
(5) for the avocado: calibrate the photo ripeness likelihood and the peeling monitor on real images, give the PR2
contact physics for the knife, the twist and the spoon (breakable joints carrying the bridge and adhesion strengths),
and estimate the seed's sideways offset; (6) verify the 108 SOURCED-VERIFY entries and calibrate the ILLUSTRATIVE ones that
the benches depend on (peel bonds, strike constants, cutting forces, friction); (7) test the avocado microtheory on
real fruit (penetrometer, squeeze stiffness and photos at known stages to calibrate the bruise ratio and the colour
and stiffness likelihoods) and take a second fruit through the theory-construction loop.
\clearpage\appendix\section{Parameter register}\label{app:register}
{\scriptsize\begin{longtable}{@{}p{0.2\linewidth}p{0.27\linewidth}p{0.12\linewidth}p{0.26\linewidth}@{}}\caption{Register of the hen egg (params.EGG, 36 entries).}\\\toprule Key & Value / distribution & Status & Source and note \\\midrule\endfirsthead
\caption[]{Register of the hen egg (params.EGG, 36 entries). (continued)}\\\toprule Key & Value / distribution & Status & Source and note \\\midrule\endhead
\texttt{length\_\allowbreak{}L} & N(57, 1.8) in [52, 63] mm & SOURCED & kop: 58.05 mm; turjaf: 56.12 mm \\
\texttt{shape\_\allowbreak{}index\_\allowbreak{}B\_\allowbreak{}over\_\allowbreak{}L} & N(0.765, 0.02) in [0.7, 0.81] & SOURCED & kop: 78.25 \%; alt: normal 72-76 \\
\texttt{asymmetry\_\allowbreak{}w} & U(1.6 .. 3.4) mm & ILLUSTRATIVE & narushin2021 parameter w; range assumed \\
\texttt{shell\_\allowbreak{}thickness} & N(0.35, 0.03) in [0.28, 0.42] mm & SOURCED & kop: 0.28-0.41 mm; sb: 0.35 mm \\
\texttt{cuticle\_\allowbreak{}thickness} & N(7.85, 1.5) in [5, 12.5] um & SOURCED & ir: 5.6-12.5 um, mean 7.85 \\
\texttt{mammillary\_\allowbreak{}thickness} & 100 um & SOURCED & fr: $\sim$100 um (palisade $\sim$300, vertical crystal 3-8) \\
\texttt{membranes\_\allowbreak{}thickness} & N(70, 8) in [50, 90] um & SOURCED & fr/mem: $\sim$70 um combined \\
\texttt{outer\_\allowbreak{}membrane\_\allowbreak{}share} & 0.75 & SOURCED & ubc: outer membrane $\sim$3x thicker than inner \\
\texttt{pore\_\allowbreak{}count} & integer 10000 .. 20000 & SOURCED & por: 10,000-20,000; $\sim$2x denser at blunt pole \\
\texttt{shell\_\allowbreak{}E} & U(15 .. 43) GPa & SOURCED & sev: 41-43 GPa nanoindentation; sb model 15 GPa \\
\texttt{shell\_\allowbreak{}poisson} & 0.3 & TYPICAL &  \\
\texttt{shell\_\allowbreak{}density} & 2300 kg/m3 & ILLUSTRATIVE & checked: shell 8-11 \% of mass (br) \\
\texttt{albumen\_\allowbreak{}density} & 1035 kg/m3 & SOURCED-VERIFY & cs: liquid white $\sim$1.03 g/ml (weak source) \\
\texttt{yolk\_\allowbreak{}density} & 1030 kg/m3 & TYPICAL &  \\
\texttt{whole\_\allowbreak{}specific\_\allowbreak{}gravity} & 1.078-1.097 fresh, 1.040-1.059 after >30 d & SOURCED & ro \\
\texttt{yolk\_\allowbreak{}mass\_\allowbreak{}fraction} & N(0.3, 0.015) in [0.27, 0.32] & SOURCED & br: 27-32 \%; usda 31 \% \\
\texttt{albumen\_\allowbreak{}layer\_\allowbreak{}shares} & outer thin 0.23, thick 0.57, inner thin 0.17, chalaziferous 0.03 & SOURCED & usda egg grading manual \\
\texttt{germinal\_\allowbreak{}disc\_\allowbreak{}diameter} & U(2 .. 3) mm & SOURCED & ab: infertile blastodisc 2-3 mm \\
\texttt{yolk\_\allowbreak{}colour\_\allowbreak{}fan} & U(9 .. 14) DSM & SOURCED & dsm YolkFan 1-16; consumers prefer $\ge$ 12 \\
\texttt{storage\_\allowbreak{}temp\_\allowbreak{}C} & one of 6, 22 C & ILLUSTRATIVE & fridge or counter \\
\texttt{age\_\allowbreak{}days} & U(1 .. 21) d & ILLUSTRATIVE & days since lay \\
\texttt{air\_\allowbreak{}cell\_\allowbreak{}height} & 1.5 + k(T) * age; k(6 C) = 0.08, k(22 C) = 0.19 mm/d mm & SOURCED-VERIFY & gr: < 4 mm after 28 d at 6 C; > 6 mm after 25 d at 22 C; graded vs eu589/usda \\
\texttt{F\_\allowbreak{}break\_\allowbreak{}equatorial} & N(35, 8) in [12, 59] N & SOURCED & kop: Hisex 36.43 N short axis, 10-59 N; icf11 34$\pm$6 N \\
\texttt{F\_\allowbreak{}break\_\allowbreak{}axial} & N(37, 9) in [12, 62] N & SOURCED & kop 34.92 N long axis; icf11 41$\pm$8 N; alt: highest along the polar axis at the sharp end \\
\texttt{F\_\allowbreak{}break\_\allowbreak{}thickness\_\allowbreak{}scaling} & F $\sim$ (t / 0.35 mm)\^{}2 & ILLUSTRATIVE & thin-shell scaling assumed \\
\texttt{shell\_\allowbreak{}contact\_\allowbreak{}stiffness} & Reissner k = 4 E t\^{}2 / (R\_local sqrt(3(1-nu\^{}2))) N/mm & DERIVED & reissner1946 + shell\_E + shell\_thickness \\
\texttt{microcrack\_\allowbreak{}band} & 0.85 & ILLUSTRATIVE & F/F\_break in [0.85, 1) may leave an invisible micro-crack \\
\texttt{microcrack\_\allowbreak{}strength\_\allowbreak{}loss} & 0.4 & ILLUSTRATIVE & F\_break of a micro-cracked shell x (1 - 0.40) \\
\texttt{membrane\_\allowbreak{}tear\_\allowbreak{}ratio} & 1.6 & ILLUSTRATIVE & check $\to$ leaker when the load exceeds 1.6 x crack force \\
\texttt{edge\_\allowbreak{}concentration} & F\_crack\_edge = F\_break * sqrt(r\_edge / 5 mm), r\_edge $\ge$ 0.2 mm & ILLUSTRATIVE & sharp tool parts crack at lower force (Morgenstern crackforce depends on parts) \\
\texttt{friction\_\allowbreak{}shell\_\allowbreak{}rubber} & U(0.45 .. 0.75) & ILLUSTRATIVE & dry shell vs gripper pad \\
\texttt{crack\_\allowbreak{}strike\_\allowbreak{}model} & food/egg\_crack.py P: coupled mass 0.05 kg; membrane tears at F/Fc 1.8 (flat), 1.05 (rim), 3.0 (egg-on-egg); shatter at 2.6 / 2.2 / 2.4; shards 0.15 / 0.6 / 0 per F/Fc; yolk puncture 0.06 per shard & ILLUSTRATIVE & structure from atk\_egg, egginfo\_crack, atk\_egg\_on\_egg; numbers to calibrate \\
\texttt{cook\_\allowbreak{}thermal\_\allowbreak{}props} & albumen 0.55/1035/3700, yolk 0.34/1030/2800, shell 0.45/2300/900 W/mK, kg/m3, J/kgK & TYPICAL & egg thermophysical properties, to cite (e.g. Denys et al. 2003) \\
\texttt{cook\_\allowbreak{}yolk\_\allowbreak{}thresholds} & f65 < 0.15 runny, < 0.45 jammy, < 0.85 creamy, else firm & CALIBRATED & alphafoodie timing table (hot start, medium eggs) + Lopez-Alt fridge-egg timings \\
\texttt{impact\_\allowbreak{}effective\_\allowbreak{}mass\_\allowbreak{}fraction} & 0.53 & CALIBRATED & c = sqrt(m\_eff)/rate fitted to cohen2025: side-on drops onto a hard surface crack $\sim$45 \% at 10 mm (model 45 \%), vertical 55 \% at 8 mm (model $\sim$57 \%); the model population is more heterogeneous than one MIT batch (8 mm side-on: model $\sim$28 \% vs 5 \%) \\
\texttt{impact\_\allowbreak{}rate\_\allowbreak{}factor} & 1.3 & ILLUSTRATIVE & dynamic/static crack force; strength rises with loading rate (only sqrt(m\_eff)/rate is identified by the calibration) \\\bottomrule\end{longtable}}
{\scriptsize\begin{longtable}{@{}p{0.2\linewidth}p{0.27\linewidth}p{0.12\linewidth}p{0.26\linewidth}@{}}\caption{Register of the apple (gala-like) (params.APPLE, 33 entries).}\\\toprule Key & Value / distribution & Status & Source and note \\\midrule\endfirsthead
\caption[]{Register of the apple (gala-like) (params.APPLE, 33 entries). (continued)}\\\toprule Key & Value / distribution & Status & Source and note \\\midrule\endhead
\texttt{cultivar} & Gala-like & ILLUSTRATIVE & shape/colour archetype (striped orange-red over yellow) \\
\texttt{diameter\_\allowbreak{}D} & N(75, 3.5) in [66, 86] mm & SOURCED & wiki\_gala: 78 \% of Gala fruit 70-80 mm \\
\texttt{height\_\allowbreak{}over\_\allowbreak{}D} & N(0.88, 0.03) in [0.8, 0.96] & SOURCED & uspto\_pp20028: 71.3 x 81.4 mm \\
\texttt{lobe\_\allowbreak{}amplitude} & U(0.008 .. 0.03) & ILLUSTRATIVE & five-fold relief at the calyx end \\
\texttt{lopsidedness} & U(0 .. 0.03) & ILLUSTRATIVE &  \\
\texttt{skin\_\allowbreak{}thickness} & N(55, 10) in [33, 73] um & SOURCED & konarska2006: 33.3-73.1, mean 55.2 um \\
\texttt{skin\_\allowbreak{}tensile\_\allowbreak{}strength} & 2.56 MPa & SOURCED-VERIFY & jets\_skin\_tensile \\
\texttt{lenticel\_\allowbreak{}density} & U(4.3 .. 12.6) 1/cm2 & SOURCED & ishs\_1099\_97: mean 7.11 /cm2 \\
\texttt{flesh\_\allowbreak{}density} & N(860, 40) in [760, 970] kg/m3 & SOURCED & wiki\_gala 0.86 g/cm3; fuji 955-968 kg/m3 (degruyter) \\
\texttt{cortex\_\allowbreak{}porosity} & U(0.17 .. 0.3) & SOURCED & fmach\_porosity: Jazz 17 \% .. Golden Del. 29.8 \% \\
\texttt{core\_\allowbreak{}density\_\allowbreak{}excess} & 60 kg/m3 & ILLUSTRATIVE & core porosity 14.5 \% vs cortex 26 \% (frontiers2015) \\
\texttt{flesh\_\allowbreak{}E} & lognormal, median 3, factor 1.7 MPa & SOURCED & fuji 5.0 MPa (consensus f4110d16); quasi-static 4.2-7.8; stored 0.5-1.7 (tjaf2005) \\
\texttt{flesh\_\allowbreak{}poisson} & 0.3 & TYPICAL &  \\
\texttt{storage\_\allowbreak{}days} & U(0 .. 120) d & ILLUSTRATIVE & E multiplied by exp(-storage\_days / 180) \\
\texttt{mt\_\allowbreak{}firmness} & N(75, 8) in [50, 95] N & SOURCED & msu2025: Gala $\sim$18 lb (80 N) at harvest \\
\texttt{flesh\_\allowbreak{}shear\_\allowbreak{}strength} & N(0.4, 0.05) in [0.28, 0.5] MPa & SOURCED & fuji shear 0.45 MPa; failure stress 0.32-0.37 MPa (consensus) \\
\texttt{bruise\_\allowbreak{}energy\_\allowbreak{}threshold} & U(0.004 .. 0.012) J & SOURCED-VERIFY & no bruise below 0.25 m/s impact; bruises from 0.011 J pendulum (consensus) \\
\texttt{bruise\_\allowbreak{}volume\_\allowbreak{}per\_\allowbreak{}J} & 0.01 * V\_fruit / 0.41 J cm3/J & SOURCED-VERIFY & fsct2008 (about 1 \% of fruit volume at 410 mJ); linear in energy \\
\texttt{cutting\_\allowbreak{}force\_\allowbreak{}peak} & N(25, 5) in [12, 40] N & ILLUSTRATIVE & no apple cutting force found \\
\texttt{friction\_\allowbreak{}skin\_\allowbreak{}rubber} & U(0.7 .. 0.95) & SOURCED-VERIFY & fuji vs TPU rubber 0.88 (consensus) \\
\texttt{core\_\allowbreak{}line\_\allowbreak{}radius} & N(14, 1.5) in [11, 18] mm & SOURCED & uspto: core diameter 24.7-36.7 mm \\
\texttt{core\_\allowbreak{}length} & N(41, 3) in [34, 48] mm & SOURCED & uspto\_pp29089: 41.1 mm \\
\texttt{locule\_\allowbreak{}length} & N(16, 1) in [13, 19] mm & SOURCED & uspto: 15.2 x 7.9, 16.6 x 9.9 mm \\
\texttt{locule\_\allowbreak{}width} & N(8.5, 0.8) in [7, 10.5] mm & SOURCED & uspto (radial width in the long. section) \\
\texttt{endocarp\_\allowbreak{}thickness} & 0.35 mm & ILLUSTRATIVE & leathery/papery carpel wall \\
\texttt{seeds\_\allowbreak{}per\_\allowbreak{}carpel} & P(0,1,2) = 0.1, 0.4, 0.5 & SOURCED & wiki\_apple: 1-2 seeds per carpel \\
\texttt{seed\_\allowbreak{}length} & N(8.5, 0.6) in [7.5, 9.6] mm & SOURCED & sciRep2021: 7.69-9.43 mm \\
\texttt{seed\_\allowbreak{}width} & N(4.5, 0.3) in [3.9, 5] mm & SOURCED & sciRep2021: 4.0-4.9 mm \\
\texttt{seed\_\allowbreak{}thickness} & N(2.7, 0.2) in [2.2, 3.2] mm & TYPICAL &  \\
\texttt{stem\_\allowbreak{}length} & U(10 .. 28) mm & SOURCED & uspto: 10.5, 12-14, 28.2 mm \\
\texttt{stem\_\allowbreak{}diameter} & N(2.8, 0.15) in [2.5, 3.1] mm & SOURCED & uspto: 2.7-3.0 mm \\
\texttt{vascular\_\allowbreak{}bundle\_\allowbreak{}radius} & 0.45 mm & ILLUSTRATIVE & 10 primary bundles (frontiers2015) \\
\texttt{browning\_\allowbreak{}dE} & 2.38 (0 min), 2.68 (30 min), 4.05 (60 min) dE & SOURCED & pmc7022590 (Golden Delicious) \\\bottomrule\end{longtable}}
{\scriptsize\begin{longtable}{@{}p{0.2\linewidth}p{0.27\linewidth}p{0.12\linewidth}p{0.26\linewidth}@{}}\caption{Register of the hass avocado (params.AVOCADO, 48 entries).}\\\toprule Key & Value / distribution & Status & Source and note \\\midrule\endfirsthead
\caption[]{Register of the hass avocado (params.AVOCADO, 48 entries). (continued)}\\\toprule Key & Value / distribution & Status & Source and note \\\midrule\endhead
\texttt{cultivar} & Hass & SOURCED & avo\_au\_stages; uspto\_pp14239 \\
\texttt{length\_\allowbreak{}L} & N(98, 6) in [80, 120] mm & SOURCED-VERIFY & specialty: 7-15 cm; no Hass mm data held \\
\texttt{diameter\_\allowbreak{}D} & N(66, 3.5) in [56, 76] mm & TYPICAL & chosen so mass $\sim$ 200 g (uspto mean 204 g) \\
\texttt{fruit\_\allowbreak{}mass\_\allowbreak{}check} & 142-340 (CAC), 200-300 (wiki), mean 204 (uspto\_pp14239) g & SOURCED &  \\
\texttt{neck} & U(0.08 .. 0.2) & ILLUSTRATIVE & obovate/pyriform taper of the stem end \\
\texttt{skin\_\allowbreak{}thickness} & N(1.6, 0.15) in [1.2, 2] mm & SOURCED & uspto\_pp14239: Hass skin 1.6 mm \\
\texttt{skin\_\allowbreak{}layers} & cuticle, epidermis (1 layer), hypodermis (1-3 cells), parenchyma, stone-cell layer & SOURCED & cs42 \\
\texttt{pebble\_\allowbreak{}scale} & U(1.2 .. 2) mm & ILLUSTRATIVE & no source for tubercle size \\
\texttt{seed\_\allowbreak{}length} & N(40, 2.5) in [34, 47] mm & DERIVED & seed mass 28.4-30.6 g (uspto) at 1.1 g/cm3 \\
\texttt{seed\_\allowbreak{}width\_\allowbreak{}ratio} & N(0.86, 0.04) in [0.78, 0.95] & SOURCED & uspto: 'variable, mostly ellipsoid' \\
\texttt{seed\_\allowbreak{}offset} & U(0.5 .. 2.5) mm & SOURCED-VERIFY & cs42: seed slightly off-centre \\
\texttt{seed\_\allowbreak{}coat\_\allowbreak{}thickness} & 0.25 mm & ILLUSTRATIVE & two papery testa layers (cs42); 'paper-thin' (ufl) \\
\texttt{endocarp\_\allowbreak{}thickness} & 0.15 mm & ILLUSTRATIVE & a few rows of flattened cells (cs42) \\
\texttt{testa\_\allowbreak{}adheres\_\allowbreak{}to} & one of seed, flesh & SOURCED-VERIFY & cs42: endocarp may adhere to the seed coat; side not settled \\
\texttt{flesh\_\allowbreak{}density} & 990 kg/m3 & ILLUSTRATIVE & whole fruit 970-1018 (clark2007) \\
\texttt{seed\_\allowbreak{}density} & 1100 kg/m3 & ILLUSTRATIVE &  \\
\texttt{skin\_\allowbreak{}density} & 1050 kg/m3 & ILLUSTRATIVE &  \\
\texttt{mass\_\allowbreak{}ratio\_\allowbreak{}seed\_\allowbreak{}skin\_\allowbreak{}flesh} & 16:12:72 (avocadosource) / 14:13:73 (uspto) \% & SOURCED & check \\
\texttt{ripeness\_\allowbreak{}stage} & U(1 .. 4.6) 1..5 & STANDARD & avo\_au\_stages; continuous hidden value \\
\texttt{ripening\_\allowbreak{}rate} & N(0.55, 0.1) in [0.3, 0.8] stage/d & SOURCED-VERIFY & ashton2006: firm-ripe $\sim$day 8, soft-ripe $\sim$9.5 at 20 C; landahl: 304 $\to$ 2 N in 6 d at 18 C \\
\texttt{penetrometer\_\allowbreak{}force} & stage 1..5 = 118, 46, 15, 8.3, 4.9 N (geometric band centres), log-linear N & SOURCED & avo\_au\_stages: hard >10, rubbery 2.2-10, breaking 1.1-2.1, firm-ripe 0.7-1.0, ripe 0.4-0.6 kgf (11.1 mm tip, 8 mm) \\
\texttt{flesh\_\allowbreak{}E} & E = 2.29 (F\_pen / 118 N)\^{}0.837 $\to$ 2.29 .. 0.16 MPa MPa & SOURCED & farinas2021: 2.29$\pm$0.75 $\to$ 0.16$\pm$0.08 MPa over ripening (endpoints); mapping ILLUSTRATIVE \\
\texttt{flesh\_\allowbreak{}poisson} & 0.45 & ILLUSTRATIVE & near-incompressible soft tissue \\
\texttt{bruise\_\allowbreak{}force\_\allowbreak{}ratio} & 0.6 & SOURCED-VERIFY & hia2018: a 10 N thumb press bruises firm-ripe fruit (visible within 48 h); ratio of thumb pressure to penetrometer pressure at stage 4 \\
\texttt{bruise\_\allowbreak{}visible\_\allowbreak{}after} & 24 h & SOURCED & hia2018: bruises take $\sim$24 h to show \\
\texttt{friction\_\allowbreak{}skin\_\allowbreak{}rubber} & U(0.6 .. 0.9) & ILLUSTRATIVE & pebbled skin \\
\texttt{skin\_\allowbreak{}colour\_\allowbreak{}Lab} & L 40$\to$27, C 15$\to$3, h 125$\to$250 over ripening (whole fruit) L*,C*,h & SOURCED & ashton2006 \\
\texttt{flesh\_\allowbreak{}bands\_\allowbreak{}Lab} & three equal bands: dark green L 74$\to$64 h 113; pale green L 83$\to$75 h 105; yellow L 83$\to$75 h 101; C 43$\to$35 L*,h & SOURCED & ashton2006 \\
\texttt{vascular\_\allowbreak{}strands} & 6 & SOURCED & cs42: pedicel cylinder splits into 6 strands, join near the apex and enter the seed coat as one \\
\texttt{squeeze\_\allowbreak{}stiffness\_\allowbreak{}model} & Isaac FEM table (test\_reports/avocado\_squeeze.json): k(4 N) 7.0 / 5.4 / 3.3 / 2.4 / 2.1 / 1.8 for stages 1 / 2 / 3 / 4 / 4.5 / 5 N/mm & DERIVED & convergence (1000 Hz/120 it vs 2000 Hz/200 it): stages $\ge$ 3 within 10 \%; stages 1-2 under-resolved (stage 1: 7.0 vs 11.9 N/mm) $\to$ unripe fruit is stiffer than the table \\
\texttt{stem\_\allowbreak{}button\_\allowbreak{}length} & U(2 .. 4) mm & SOURCED-VERIFY & alder: $\sim$3 mm clipped button \\
\texttt{process\_\allowbreak{}protocol} & cut lengthwise round the seed, twist the halves apart, remove the seed with a spoon (knife strike discouraged) & SOURCED & afm\_howto; farley2019 (knife injuries when pitting) \\
\texttt{skin\_\allowbreak{}puncture\_\allowbreak{}force} & U(6 .. 15) N & ILLUSTRATIVE & GAP: no knife-puncture data for Hass skin held; sharp chef's knife at the stem end \\
\texttt{skin\_\allowbreak{}cut\_\allowbreak{}resistance} & U(1.5 .. 3) N/mm & ILLUSTRATIVE & force per mm of skin thickness while slicing along the perimeter; GAP \\
\texttt{flesh\_\allowbreak{}cut\_\allowbreak{}toughness\_\allowbreak{}unripe} & 300 J/m2 & ILLUSTRATIVE & analogue: fruit-flesh cutting toughness 0.1-0.4 kJ/m2 (apple, potato); scaled with sqrt(E / E\_unripe) over ripening \\
\texttt{seed\_\allowbreak{}adhesion\_\allowbreak{}unripe} & 20 kPa & ILLUSTRATIVE & seed-flesh interface shear strength, scaled with (E / E\_unripe)\^{}0.7; GAP \\
\texttt{clean\_\allowbreak{}twist\_\allowbreak{}stage} & 2.8 1..5 & ILLUSTRATIVE & centre of the clean-separation transition: unripe flesh clings to the seed (twin AvocadoTruth.halve: clean from $\sim$3; afm\_howto: ripe fruit twists apart) \\
\texttt{testa\_\allowbreak{}remnant\_\allowbreak{}fraction} & U(0.4 .. 1) & ILLUSTRATIVE & share of the seed coat left in the cavity when the testa adheres to the flesh (cs42) \\
\texttt{knife\_\allowbreak{}blade\_\allowbreak{}thickness} & 1.6 mm & TYPICAL & chef's knife spine 2-2.5 mm, edge region $\sim$1-2 mm \\
\texttt{storage\_\allowbreak{}temperatures} & counter 20, fridge 4 (ripe fruit 2-4 C; softening fruit 5-8 C up to 10 d) C & SOURCED & ucd\_avocado; marques2011 \\
\texttt{ripening\_\allowbreak{}q10} & 2.5 & ILLUSTRATIVE & temperature factor of the ripening rate per 10 C (respiration-type Q10); no avocado-specific value held \\
\texttt{ripe\_\allowbreak{}window\_\allowbreak{}counter} & U(24 .. 48) h & SOURCED-VERIFY & ripe\_window\_guides: stage 5 lasts $\sim$1-2 d at room temperature \\
\texttt{decline\_\allowbreak{}per\_\allowbreak{}day\_\allowbreak{}counter} & 0.6 1/d & ILLUSTRATIVE & eating-quality loss per day past the ripe window (overripe: soft, stringy, brown spots), slowed by the fridge with ripening\_q10 \\
\texttt{spoil\_\allowbreak{}days\_\allowbreak{}past\_\allowbreak{}window} & U(1 .. 3) d & ILLUSTRATIVE & days past the ripe window at 20 C until the fruit is unusable (rot, flesh browning) \\
\texttt{browning\_\allowbreak{}rate\_\allowbreak{}cut} & 3 dE/sqrt(h) & ILLUSTRATIVE & enzymatic browning of exposed flesh at 20 C, dE $\sim$ k sqrt(t) (PPO, oxygen-limited); apple slices dE 2.4-4 within 1 h (pmc7022590) as an analogue \\
\texttt{browning\_\allowbreak{}visible\_\allowbreak{}dE} & 5 dE & ILLUSTRATIVE & colour difference at which the cut face looks brown \\
\texttt{browning\_\allowbreak{}treatment\_\allowbreak{}factors} & none 1.0, lemon juice 0.45, plastic wrap 0.35, wrap + fridge 0.12, vacuum + fridge 0.03; seed left in: no effect on the cut face & ILLUSTRATIVE & ordering from atk\_browning; factors are assumptions \\
\texttt{peel\_\allowbreak{}coat\_\allowbreak{}success} & 0.85 & ILLUSTRATIVE & chance that one pass with the spoon tip lifts the seed-coat sheets left in the cavity \\\bottomrule\end{longtable}}
{\scriptsize\begin{longtable}{@{}p{0.2\linewidth}p{0.27\linewidth}p{0.12\linewidth}p{0.26\linewidth}@{}}\caption{Register of the wheat roll (weizenbrötchen) (params.ROLL, 46 entries).}\\\toprule Key & Value / distribution & Status & Source and note \\\midrule\endfirsthead
\caption[]{Register of the wheat roll (weizenbrötchen) (params.ROLL, 46 entries). (continued)}\\\toprule Key & Value / distribution & Status & Source and note \\\midrule\endhead
\texttt{product} & Weizenbroetchen, oval, one lengthwise score ('Schnittbroetchen') & STANDARD & leitsaetze\_1993: $\ge$ 90 \% wheat flour; testde\_2021: the score raises the crust share \\
\texttt{length} & N(102, 5) in [90, 115] mm & ILLUSTRATIVE & oval single-score roll; round crispy roll $\sim$90 mm (hirte\_2012), round buns 86-92 mm (musiienko\_2025) \\
\texttt{width} & N(80, 4) in [70, 90] mm & ILLUSTRATIVE & chosen with length/height for $\sim$50 g at $\sim$0.2 g/cm3 \\
\texttt{height} & N(52, 3) in [44, 60] mm & SOURCED & hirte\_2012 $\sim$50 mm; musiienko\_2025 50-55 mm \\
\texttt{mass\_\allowbreak{}check} & 50-90 (bake-at-home rolls, testde\_2021); RMT $\sim$55 g dough per roll (ireks\_rmt) g & SOURCED &  \\
\texttt{base\_\allowbreak{}edge\_\allowbreak{}height} & U(6 .. 10) mm & ILLUSTRATIVE & rounded base edge \\
\texttt{lumpiness} & U(0.01 .. 0.03) & ILLUSTRATIVE & hand-moulded irregularity \\
\texttt{score\_\allowbreak{}length\_\allowbreak{}frac} & U(0.7 .. 0.85) & ILLUSTRATIVE & share of the long axis \\
\texttt{score\_\allowbreak{}width} & U(14 .. 22) mm & ILLUSTRATIVE & opened 'Ausbund' width \\
\texttt{score\_\allowbreak{}depth} & U(4 .. 8) mm & ILLUSTRATIVE &  \\
\texttt{crust\_\allowbreak{}thickness\_\allowbreak{}top} & N(1.3, 0.2) in [0.9, 2] mm & SOURCED & hirte\_2012: cut-off crust 1.0-1.5 mm \\
\texttt{crust\_\allowbreak{}thickness\_\allowbreak{}score} & U(0.6 .. 0.9) mm & ILLUSTRATIVE & thinner, paler Ausbund crust \\
\texttt{crust\_\allowbreak{}thickness\_\allowbreak{}bottom} & U(1.2 .. 1.8) mm & ILLUSTRATIVE & GAP: no bottom-crust measurement held \\
\texttt{transition\_\allowbreak{}thickness} & U(1 .. 2) mm & ILLUSTRATIVE & soleimani\_2014 distinguishes an inner-crust colour band; thickness GAP \\
\texttt{crumb\_\allowbreak{}density} & N(180, 15) in [150, 220] kg/m3 & DERIVED & roll specific volume 5.9-6.4 ml/g (kurek\_2016) $\to$ whole $\sim$0.16-0.17 g/cm3, minus the crust share; crumb 0.088-0.252 (zghal\_1999, SOURCED-VERIFY) \\
\texttt{crust\_\allowbreak{}density} & N(320, 30) in [260, 400] kg/m3 & TYPICAL & crust 50-100 \% denser than crumb \\
\texttt{transition\_\allowbreak{}density} & 240 kg/m3 & ILLUSTRATIVE & between crust and crumb \\
\texttt{crumb\_\allowbreak{}porosity} & U(0.73 .. 0.81) & SOURCED & hirte\_2012 uCT \\
\texttt{cell\_\allowbreak{}diameter\_\allowbreak{}median} & N(0.7, 0.08) in [0.5, 0.9] mm & DERIVED & lagrain\_2012 mean 2D cell area 0.32 mm2 $\to$ equivalent diameter 0.64 mm; 2D under-estimates (x1.5 volume, rathnayake\_2018) \\
\texttt{cell\_\allowbreak{}diameter\_\allowbreak{}gsd\_\allowbreak{}log} & 0.55 & ILLUSTRATIVE & lognormal spread of cell size (2-2189e3 um3 range, hirte) \\
\texttt{cell\_\allowbreak{}number\_\allowbreak{}density} & 1200 1/cm3 & DERIVED & 83 cells/cm2 (lagrain\_2012) / $\sim$0.07 cm mean diameter \\
\texttt{cell\_\allowbreak{}elongation} & U(1.2 .. 1.5) & ILLUSTRATIVE & cells stretched in the rise direction \\
\texttt{crust\_\allowbreak{}L} & N(52, 4) in [42, 62] L* & SOURCED-VERIFY & pan-bread crust 47.6 (chaple\_2023); roll browning index 50.4 $\to$ L* $\sim$50 (kurek\_2016) \\
\texttt{crust\_\allowbreak{}C} & 36 C* & SOURCED-VERIFY & chaple\_2023 a* 17.05, b* 32.2 $\to$ C* 36.4 \\
\texttt{crust\_\allowbreak{}h} & 62 deg & SOURCED-VERIFY & chaple\_2023 $\to$ h 62 \\
\texttt{crumb\_\allowbreak{}L} & 74 L* & SOURCED-VERIFY & pan-bread crumb 72.06 (chaple\_2023); white roll slightly lighter \\
\texttt{crumb\_\allowbreak{}C} & 15 C* & SOURCED-VERIFY & chaple\_2023 a* 0.94, b* 15.37 \\
\texttt{crumb\_\allowbreak{}h} & 86 deg & SOURCED-VERIFY & chaple\_2023 \\
\texttt{variant} & one of plain, dusted & ILLUSTRATIVE & flour-dusted or plain \\
\texttt{slice\_\allowbreak{}height\_\allowbreak{}frac} & 0.45 & TYPICAL & wikihow\_slice, tastingtable\_bagel: horizontally through the middle \\
\texttt{friction\_\allowbreak{}crust\_\allowbreak{}rubber} & U(0.5 .. 0.8) & ILLUSTRATIVE &  \\
\texttt{age\_\allowbreak{}h} & one of 0.5, 2, 6, 20, 30 h & ILLUSTRATIVE & time since baking (fresh .. 'vom Vortag') \\
\texttt{storage} & one of paper\_bag, plastic\_bag, open & ILLUSTRATIVE & oekotest\_storage: paper bag recommended; plastic softens the crust \\
\texttt{crumb\_\allowbreak{}E\_\allowbreak{}kPa\_\allowbreak{}by\_\allowbreak{}age} & 20 @ 5, 70 @ 48, 125 @ 96, 157 @ 168 kPa @ h & SOURCED-VERIFY & lagrain\_2012 pan bread, read off figures; log-time interpolation \\
\texttt{crumb\_\allowbreak{}poisson} & 0.103 & SOURCED & lagrain\_2012 \\
\texttt{crust\_\allowbreak{}water\_\allowbreak{}model} & w(t) = 9.73 t\_h\^{}0.112 (40 \% RH, 22 C) \% & CALIBRATED & power law through hirte\_2012 9.0 \% @ 0.5 h and 13.9 \% @ 24 h; predicts 10.6 \% @ 2.1 h (measured: 10.1 \% critical reached at 2.1 h) $\to$ residual +0.5 points \\
\texttt{crust\_\allowbreak{}water\_\allowbreak{}critical} & 10.1 \% & SOURCED & hirte\_2012: half the sensory crispness lost \\
\texttt{plastic\_\allowbreak{}bag\_\allowbreak{}factor} & 4 & ILLUSTRATIVE & crust softening clock runs x4 in a closed bag (oekotest: 'plastic softens the crust'; no rate held) \\
\texttt{crust\_\allowbreak{}break\_\allowbreak{}force} & 3 N & SOURCED & hirte\_2012: $\ge$ 3 N with a 1 x 10 mm blade \\
\texttt{avrami\_\allowbreak{}firmness} & F = 66.01 - 59.53 exp(-0.186 t\_d) N & SOURCED & cardone\_2020 loaf (n=1, k per day) \\
\texttt{reheat\_\allowbreak{}oven} & 150-170, 5-10 C, min & TYPICAL & utopia\_aufbacken \\
\texttt{toaster\_\allowbreak{}rack} & 2-3 + flip + 1-2 min & SOURCED-VERIFY & utopia\_aufbacken \\
\texttt{toast\_\allowbreak{}times} & 3 light .. 4.5 'starting to burn' (sliced bread, home toaster) min & SOURCED & kaur\_2026 \\
\texttt{saw\_\allowbreak{}setpoint} & -30 N, 0.02 m/s, $\pm$0.02 m strokes (baguette) N, m/s, m & SOURCED & luo\_2026 (robot) \\
\texttt{crumb\_\allowbreak{}squash\_\allowbreak{}strain} & 0.3 & ILLUSTRATIVE & permanent crumb compression when pressing instead of sawing \\
\texttt{squeeze\_\allowbreak{}k\_\allowbreak{}model} & k = 0.29 E / 20 kPa (30 mm pads across the width, 1 $\to$ 3 N) N/mm & DERIVED & Isaac FEM (test\_reports/roll\_squeeze\_summary.json): 20 kPa lane converged (0.276 $\to$ 0.29 N/mm, 1000 $\to$ 2000 Hz); stiffer lanes under-converged $\to$ linear-elastic scaling; crust excluded (GAP) \\\bottomrule\end{longtable}}
{\scriptsize\begin{longtable}{@{}p{0.2\linewidth}p{0.27\linewidth}p{0.12\linewidth}p{0.26\linewidth}@{}}\caption{Register of the nantes carrot (params.CARROT, 42 entries).}\\\toprule Key & Value / distribution & Status & Source and note \\\midrule\endfirsthead
\caption[]{Register of the nantes carrot (params.CARROT, 42 entries). (continued)}\\\toprule Key & Value / distribution & Status & Source and note \\\midrule\endhead
\texttt{type} & Nantes (cylindrical, blunt tip, small core), loose, topped & STANDARD & eu730\_1999; con22 \\
\texttt{length} & N(190, 15) in [150, 230] mm & SOURCED & con22: 182.7 / 210.8 mm (Nantes) \\
\texttt{shoulder\_\allowbreak{}diameter} & N(31, 2) in [25, 38] mm & SOURCED & con22: 30.3 / 31.4 mm \\
\texttt{mass\_\allowbreak{}check} & 110-135 (con22 Nantes); $\ge$ 50 g main crop (eu730\_1999) g & SOURCED &  \\
\texttt{tip\_\allowbreak{}bluntness} & U(0.3 .. 0.9) & ILLUSTRATIVE & Nantes 'blunt tip' (catalogues) \\
\texttt{bend\_\allowbreak{}radius} & U(600 .. 3000) mm & ILLUSTRATIVE & slight curvature \\
\texttt{periderm\_\allowbreak{}thickness} & 150 um & ILLUSTRATIVE & GAP (knowles \& flore 1983 to obtain); $\sim$3 cells of $\sim$50 um \\
\texttt{xylem\_\allowbreak{}radius\_\allowbreak{}fraction} & N(0.33, 0.04) in [0.25, 0.42] & SOURCED-VERIFY & singh17 $\sim$0.38 over mixed genotypes; Nantes bred for a small core; vob24 trend crown $\to$ tip; peng14 4-7 mm core \\
\texttt{primary\_\allowbreak{}xylem\_\allowbreak{}diameter} & 1 mm & ILLUSTRATIVE &  \\
\texttt{petiole\_\allowbreak{}stubs} & integer 4 .. 9 & ILLUSTRATIVE &  \\
\texttt{petiole\_\allowbreak{}stub\_\allowbreak{}length} & U(2 .. 6) mm & ILLUSTRATIVE & topped carrots \\
\texttt{green\_\allowbreak{}shoulder\_\allowbreak{}length} & one of 0, 6, 12 mm & STANDARD & eu730\_1999 Class I $\le$ 20 mm on long roots (instances drawn inside Class I) \\
\texttt{lateral\_\allowbreak{}root\_\allowbreak{}rows} & 4 & SOURCED-VERIFY & diarch\_rows \\
\texttt{lateral\_\allowbreak{}root\_\allowbreak{}scar\_\allowbreak{}spacing} & U(8 .. 15) mm & ILLUSTRATIVE &  \\
\texttt{lenticel\_\allowbreak{}band\_\allowbreak{}spacing} & U(2 .. 4) mm & ILLUSTRATIVE & GAP: no spacing found \\
\texttt{lenticel\_\allowbreak{}band\_\allowbreak{}relief} & U(0.05 .. 0.12) mm & ILLUSTRATIVE &  \\
\texttt{phloem\_\allowbreak{}LCh} & (55, 49, 54) L*,C*,h & SOURCED-VERIFY & pac16\_dry raw flesh 55.35/28.9/40.05 \\
\texttt{xylem\_\allowbreak{}LCh} & (63, 45, 64) L*,C*,h & ILLUSTRATIVE & paler, yellower core (lutein; aub22) -- no L*a*b* \\
\texttt{skin\_\allowbreak{}LCh} & (50, 44, 52) L*,C*,h & ILLUSTRATIVE & duller than the flesh \\
\texttt{soil\_\allowbreak{}residue} & U(0 .. 0.25) & ILLUSTRATIVE & washed supermarket carrots \\
\texttt{friction\_\allowbreak{}skin\_\allowbreak{}rubber} & U(0.5 .. 0.8) & ILLUSTRATIVE & moist skin \\
\texttt{density} & 1040 kg/m3 & SOURCED-VERIFY & jah18 1.04 g/cm3 (1.00-1.17) \\
\texttt{E\_\allowbreak{}phloem} & N(4.6, 0.15) in [4.33, 4.85] MPa & SOURCED-VERIFY & vob24 after Stopa 2017 \\
\texttt{E\_\allowbreak{}xylem} & N(4.75, 0.08) in [4.62, 4.88] MPa & SOURCED-VERIFY & vob24 after Stopa 2017 \\
\texttt{poisson} & 0.475 & SOURCED-VERIFY & vob24 \\
\texttt{E\_\allowbreak{}decay\_\allowbreak{}7d} & 1.32 & SOURCED & vob24: E / 1.32 after 7 d in the fridge (zip bag) \\
\texttt{bend\_\allowbreak{}Fmax} & 48.6 N & SOURCED-VERIFY & jah18 three-point bending (span: see source) \\
\texttt{bend\_\allowbreak{}deflection} & 21.9 mm & SOURCED-VERIFY & jah18 \\
\texttt{blade\_\allowbreak{}shear} & 41.1 N & SOURCED-VERIFY & jah18 straight blade 1.4 mm, 30 deg \\
\texttt{impact\_\allowbreak{}damage\_\allowbreak{}force} & 43-45 N & SOURCED-VERIFY & xia21 \\
\texttt{fracture\_\allowbreak{}energy} & 40-400 (mixed produce) J/m2 & SOURCED-VERIFY & alv00 \\
\texttt{turgor\_\allowbreak{}fresh} & 0.7 MPa & SOURCED-VERIFY & herp99\_calbo00 \\
\texttt{limp\_\allowbreak{}rwc} & 0.95 & SOURCED-VERIFY & herp99\_calbo00: firmness from turgor $\sim$0 below RWC 0.95 \\
\texttt{softening\_\allowbreak{}kinetics} & 2nd-order fractional conversion; k(80 C) 0.035, k(110 C) 1.453; Ea 138.9 kJ/mol; Finf/F0 0.02 1/min & SOURCED & peng14 \\
\texttt{peel\_\allowbreak{}depth} & 2 mm & SOURCED & con22 hand peeler; mass 12.1-14.6 \% \\
\texttt{water\_\allowbreak{}loss\_\allowbreak{}rate} & fridge in bag 0.3, room open 2.0 \%/d & ILLUSTRATIVE & storage\_loss: unwrapped 20 C unacceptable after $\sim$3 d at 5-8 \% loss $\to$ $\sim$2 \%/d; fridge GAP \\
\texttt{age\_\allowbreak{}d} & U(1 .. 14) d & ILLUSTRATIVE & since purchase; hidden \\
\texttt{storage} & one of fridge\_bag, fridge\_open, room\_open & ILLUSTRATIVE &  \\
\texttt{blush\_\allowbreak{}reversible\_\allowbreak{}h} & 1.25 h & SOURCED-VERIFY & sim09\_cz95 \\
\texttt{blush\_\allowbreak{}irreversible\_\allowbreak{}h} & 15 h & SOURCED-VERIFY & sim09\_cz95 \\
\texttt{boil\_\allowbreak{}rounds\_\allowbreak{}6mm} & 4-5 crisp-tender, +1-2 soft min & SOURCED-VERIFY & carrot\_boil \\
\texttt{cut\_\allowbreak{}dims} & rounds 5-6; julienne 3x3x45; brunoise 3; batonnet 6x6x55; dice 6/12/19 mm & TYPICAL & carrot\_cuts \\\bottomrule\end{longtable}}
{\scriptsize\begin{longtable}{@{}p{0.2\linewidth}p{0.27\linewidth}p{0.12\linewidth}p{0.26\linewidth}@{}}\caption{Register of the cavendish banana (params.BANANA, 41 entries).}\\\toprule Key & Value / distribution & Status & Source and note \\\midrule\endfirsthead
\caption[]{Register of the cavendish banana (params.BANANA, 41 entries). (continued)}\\\toprule Key & Value / distribution & Status & Source and note \\\midrule\endhead
\texttt{cultivar} & Cavendish (Musa acuminata AAA; Grand Naine / Valery) & SOURCED & ciro05; mal22 \\
\texttt{length\_\allowbreak{}outer} & N(210, 15) in [170, 250] mm & SOURCED & convex-face length incl. pedicel (eu1333\_2011 definition): 190 (mal22), 233 (ciro05) \\
\texttt{calibre} & N(35, 1.5) in [30, 40] mm & SOURCED & ciro05 34/36/34 mm; eu1333 min 27 mm \\
\texttt{mass\_\allowbreak{}check} & 102 (bant24), 105 (chiw23), 126-148 (polat19) g & SOURCED &  \\
\texttt{curvature\_\allowbreak{}radius} & U(150 .. 260) mm & ILLUSTRATIVE & 'slightly curved' (mal22); GAP \\
\texttt{ridges} & one of 4, 5 & ILLUSTRATIVE & 'slightly ridged' (mal22), 'nearly 4-sided'; count GAP \\
\texttt{angularity} & U(0.04 .. 0.08) & ILLUSTRATIVE & ridges round out as the finger matures (fao) \\
\texttt{peel\_\allowbreak{}thickness} & N(2.6, 0.25) in [2, 3.2] mm & SOURCED & mal22 2.56 mm (Grand Naine); amn12 1.1-3.8 mm, thinning on ripening \\
\texttt{peel\_\allowbreak{}thinning} & 0.2 & ILLUSTRATIVE & fractional thinning stage 1 $\to$ 7 (amn12: thinning, no \% held) \\
\texttt{peel\_\allowbreak{}layer\_\allowbreak{}fractions} & (0.15, 0.6, 0.25) & ILLUSTRATIVE & chlorenchyma / parenchyma / inner aerenchyma share of the peel (amn12 layer order) \\
\texttt{pulp\_\allowbreak{}peel\_\allowbreak{}ratio\_\allowbreak{}check} & 1.69-2.69 (chiw23, bant24, mal22); rises with ripening & SOURCED &  \\
\texttt{central\_\allowbreak{}axis\_\allowbreak{}diameter} & 3 mm & ILLUSTRATIVE & placental axis \\
\texttt{stalk\_\allowbreak{}length} & N(17.3, 1.5) in [13, 22] mm & SOURCED & mal22 pedicel \\
\texttt{stalk\_\allowbreak{}width} & N(10.15, 0.8) in [8, 12.5] mm & SOURCED & mal22 pedicel \\
\texttt{phloem\_\allowbreak{}strings} & integer 3 .. 6 & ILLUSTRATIVE & 'strings' (guapo22); count GAP \\
\texttt{ovules\_\allowbreak{}total} & integer 200 .. 300 & SOURCED & promusa: 200-300 ovules \\
\texttt{ovule\_\allowbreak{}diameter} & 0.6 mm & ILLUSTRATIVE & tiny black specks \\
\texttt{ripeness\_\allowbreak{}stage} & U(2 .. 7) 1..7 & STANDARD & markon / ma22 seven-stage colour scale \\
\texttt{peel\_\allowbreak{}hue\_\allowbreak{}by\_\allowbreak{}stage} & [118, 98.7, 97.3, 92.7, 89.4, 87.5, 82] deg & SOURCED-VERIFY & ringer21 R2..R7; stage 1 from col20 Gros Michel day 0 (h $\sim$122) $\to$ 118 \\
\texttt{peel\_\allowbreak{}L\_\allowbreak{}by\_\allowbreak{}stage} & [48, 55, 62, 68, 72, 72, 64] L* & SOURCED-VERIFY & shape: hou15 L* peaks at stage 5 then falls; values: hega26 45-57 $\to$ 55-73 $\to$ 52-66, shahir14 max 74 \\
\texttt{peel\_\allowbreak{}C\_\allowbreak{}by\_\allowbreak{}stage} & [42, 50, 56, 60, 62, 62, 58] C* & SOURCED-VERIFY & ringer21 b* 56-72 (a* $\sim$0 at R5 $\to$ C* $\sim$ b*); col20 green C* $\sim$48 \\
\texttt{pulp\_\allowbreak{}LCh} & (86, 24, 92) L*,C*,h & ILLUSTRATIVE & 'cream' (mal22); numeric GAP \\
\texttt{sugar\_\allowbreak{}spot\_\allowbreak{}density} & 1 & ILLUSTRATIVE & stage 7 = brown flecks (markon); kinetics GAP \\
\texttt{friction\_\allowbreak{}peel\_\allowbreak{}rubber} & U(0.5 .. 0.8) & ILLUSTRATIVE & waxy peel \\
\texttt{peel\_\allowbreak{}density} & 950 kg/m3 & ILLUSTRATIVE & inner aerenchyma; whole fruit SG 1.023 (arcc, SV) \\
\texttt{pulp\_\allowbreak{}density} & 1050 kg/m3 & ILLUSTRATIVE & whole fruit SG 1.023 (arcc, SV) \\
\texttt{pulp\_\allowbreak{}E\_\allowbreak{}compression} & 1.75 (day 1) $\to$ 1.09 (day 12) MPa & SOURCED & ciro05 cylinders 20 x 40 mm \\
\texttt{pulp\_\allowbreak{}E\_\allowbreak{}bending} & 5.16 $\to$ 2.65 MPa & SOURCED & ciro05 \\
\texttt{pulp\_\allowbreak{}firmness} & 38.0 $\to$ 14.0 over 10 d (2 mm probe) N & SOURCED & jaiswal12; chiw23 44.8 $\to$ 22.1 \\
\texttt{peel\_\allowbreak{}puncture} & 13.6-15.2 (2 mm probe) N & SOURCED & ciro05 \\
\texttt{peel\_\allowbreak{}fracture} & 53.0 $\to$ 2.1 on ripening N & SOURCED-VERIFY & khanal23 (test geometry not in abstract) \\
\texttt{whole\_\allowbreak{}fracture} & 129 bending (14 cm span), 121 compression N & SOURCED & ciro05 \\
\texttt{finger\_\allowbreak{}drop\_\allowbreak{}force} & $\sim$49 on day 1, > 2x lower by day 3 N & SOURCED-VERIFY & imsabai; mbeguie09 \\
\texttt{bruise\_\allowbreak{}threshold} & 200 mJ & SOURCED-VERIFY & bugaud14: no visible bruise up to 200 mJ (Grande Naine) \\
\texttt{peel\_\allowbreak{}bond\_\allowbreak{}ratios} & (0.45, 0.3, 1.5) & ILLUSTRATIVE & peel-to-pulp adhesion / ridge seam / strip tear strength as multiples of the peel fracture force (khanal23 53 $\to$ 2.1 N, SOURCED-VERIFY); adhesion > seam so the peel splits along the ridges (khanal23) instead of sliding off as a sleeve; GAP: no peeling-force study \\
\texttt{stage\_\allowbreak{}rate\_\allowbreak{}20C} & N(0.6, 0.15) in [0.3, 1] stage/d & SOURCED-VERIFY & hou15: 1 $\to$ 7 in 7 d with ethylene (0.86/d); ma22 day bands $\sim$0.4/d; 18 C RS6 at 9 d \\
\texttt{q10\_\allowbreak{}ripening} & 2 & ILLUSTRATIVE & temperature dependence of the stage rate \\
\texttt{chilling\_\allowbreak{}threshold} & 13 C & SOURCED & wflo \\
\texttt{brix\_\allowbreak{}by\_\allowbreak{}day} & 4.5 / 6.8 / 8 / 16.7 / 23.7 at day 1 / 4 / 6 / 8 / 10 Bx & SOURCED & chiw23 \\
\texttt{mass\_\allowbreak{}loss\_\allowbreak{}rate} & 1 \%/d & SOURCED-VERIFY & chiw23 $\sim$2.2 \%/d (22-27 C, outlier-looking); kojima96 $\sim$0 \\
\texttt{rounds\_\allowbreak{}thickness} & 6 mm & TYPICAL & banana\_howto \\\bottomrule\end{longtable}}
{\scriptsize\begin{longtable}{@{}p{0.2\linewidth}p{0.27\linewidth}p{0.12\linewidth}p{0.26\linewidth}@{}}\caption{Register of the sweet orange (params.ORANGE, 38 entries).}\\\toprule Key & Value / distribution & Status & Source and note \\\midrule\endfirsthead
\caption[]{Register of the sweet orange (params.ORANGE, 38 entries). (continued)}\\\toprule Key & Value / distribution & Status & Source and note \\\midrule\endhead
\texttt{cultivar} & one of navel, valencia & ILLUSTRATIVE & German supermarket: navel (winter) / Valencia (summer) \\
\texttt{diameter} & N(75, 6) in [61, 88] mm & SOURCED & ucr\_citrus Washington navel 61-88 mm \\
\texttt{shape\_\allowbreak{}index} & U(1 .. 1.08) & SOURCED & ucr\_citrus L/W 1.0-1.1 \\
\texttt{mass\_\allowbreak{}check} & 162-335 (navel), 113-299 (Valencia) g & SOURCED & ucr\_citrus \\
\texttt{peel\_\allowbreak{}thickness} & N(5, 0.6) in [3.5, 6.4] mm & SOURCED & ucr\_citrus rind 4.1-6.4 / 3.5-6.2 mm \\
\texttt{flavedo\_\allowbreak{}thickness} & U(1 .. 1.6) mm & ILLUSTRATIVE & peel minus albedo (2-5 mm, citrus\_patents; 3.9 mm smith18) \\
\texttt{core\_\allowbreak{}diameter} & U(5 .. 9) mm & ILLUSTRATIVE & central axis; GAP \\
\texttt{segments} & one of 9, 10, 11, 12 & SOURCED-VERIFY & citrus\_patents 9.6-11 \\
\texttt{membrane\_\allowbreak{}thickness} & 0.12 mm & ILLUSTRATIVE & endocarp film; GAP \\
\texttt{seeds} & integer 1 .. 6 & SOURCED & ucr\_citrus Valencia 1.2-6.6; navel $\sim$0 (set in the model) \\
\texttt{seed\_\allowbreak{}length} & 9 mm & TYPICAL &  \\
\texttt{oil\_\allowbreak{}gland\_\allowbreak{}density} & 50 1/cm2 & SOURCED-VERIFY & knight01 8,000-12,000 per fruit $\to$ $\sim$40-60 /cm2 (derived); one cultivar 208 /cm2 (citrus\_patents) \\
\texttt{peel\_\allowbreak{}density} & 746 kg/m3 & SOURCED & jentzsch22 \\
\texttt{albedo\_\allowbreak{}density} & 700 kg/m3 & ILLUSTRATIVE & spongy with air spaces (sadka19) \\
\texttt{pulp\_\allowbreak{}density} & 1045 kg/m3 & DERIVED & juice SG 1.045 (orange\_juice) \\
\texttt{peel\_\allowbreak{}LCh} & (60, 74, 57) L*,C*,h & ILLUSTRATIVE & chosen for CCI $\sim$ +11 (cci market target +10..+20); measured L*a*b* GAP \\
\texttt{albedo\_\allowbreak{}LCh} & (92, 10, 90) L*,C*,h & ILLUSTRATIVE & white; GAP \\
\texttt{pulp\_\allowbreak{}LCh} & (66, 62, 66) L*,C*,h & ILLUSTRATIVE & 'medium orange' (ucr\_citrus); GAP \\
\texttt{membrane\_\allowbreak{}LCh} & (84, 28, 80) L*,C*,h & ILLUSTRATIVE & translucent pale; GAP \\
\texttt{friction\_\allowbreak{}peel\_\allowbreak{}rubber} & U(0.5 .. 0.8) & ILLUSTRATIVE & waxy, pitted peel \\
\texttt{flavedo\_\allowbreak{}E} & 103 MPa & SOURCED & smith18 (30-300 across species) \\
\texttt{albedo\_\allowbreak{}E} & 0.79 MPa & SOURCED & smith18 ($\pm$0.58) \\
\texttt{peel\_\allowbreak{}compressive\_\allowbreak{}E} & 0.02 MPa & SOURCED-VERIFY & jentzsch22 whole peel, compression \\
\texttt{peel\_\allowbreak{}poisson} & -0.05 & SOURCED & jentzsch22 \\
\texttt{punch\_\allowbreak{}force} & 15.5-79 (4.3-14.3 mm punches) N & SOURCED-VERIFY & orange\_mech \\
\texttt{rupture\_\allowbreak{}force} & 140-221 (Valencia) N & SOURCED-VERIFY & orange\_mech \\
\texttt{gland\_\allowbreak{}damage\_\allowbreak{}force} & 10-250 (sphere; oil release) N & SOURCED-VERIFY & orange\_mech \\
\texttt{oil\_\allowbreak{}jet} & 8.5 $\pm$ 4.0, 102 $\pm$ 20 m/s, um & SOURCED & smith18 \\
\texttt{gland\_\allowbreak{}burst\_\allowbreak{}pressure} & 30-140 kPa & SOURCED & smith18 \\
\texttt{juice\_\allowbreak{}fraction} & U(45 .. 54) \% & SOURCED & ucr\_citrus navel 45-54 (Valencia 37-56) \\
\texttt{juice\_\allowbreak{}by\_\allowbreak{}hand} & 60-80 per medium orange mL & TYPICAL & orange\_juice \\
\texttt{brix} & U(9.3 .. 14.1) Bx & SOURCED & ucr\_citrus \\
\texttt{juice\_\allowbreak{}viscosity} & 5 mPa s & SOURCED-VERIFY & orange\_juice 4-6 at 8 C \\
\texttt{peelability} & navel medium-easy; Valencia hard (strips break); cold fruit harder & SOURCED & ritenour13 (no N values) \\
\texttt{peel\_\allowbreak{}strip\_\allowbreak{}force} & navel 4, Valencia 8 per 16 mm strip N & ILLUSTRATIVE & ranking from ritenour13; GAP \\
\texttt{storage} & 3-8, 90-95, $\le$ 3 months; chilling after 3-6 weeks at $\sim$4 C C, \%RH & SOURCED-VERIFY & orange\_storage \\
\texttt{bitter\_\allowbreak{}juice\_\allowbreak{}h} & 3 h & SOURCED-VERIFY & edis891 'within several hours' (navel) \\
\texttt{age\_\allowbreak{}d} & U(2 .. 20) d & ILLUSTRATIVE & since purchase \\\bottomrule\end{longtable}}
{\scriptsize\begin{longtable}{@{}p{0.2\linewidth}p{0.27\linewidth}p{0.12\linewidth}p{0.26\linewidth}@{}}\caption{Register of the sandwich loaf (toastbrot) (params.LOAF, 26 entries).}\\\toprule Key & Value / distribution & Status & Source and note \\\midrule\endfirsthead
\caption[]{Register of the sandwich loaf (toastbrot) (params.LOAF, 26 entries). (continued)}\\\toprule Key & Value / distribution & Status & Source and note \\\midrule\endhead
\texttt{product} & Toastbrot (Buttertoast), lidded tin, pre-sliced, 500 g & SOURCED & test25\_toast; toast\_pack \\
\texttt{length} & N(250, 6) in [235, 265] mm & SOURCED-VERIFY & toast\_pack $\sim$0.25 m \\
\texttt{width} & N(100, 3) in [92, 110] mm & SOURCED-VERIFY & tins 100 / 110 mm \\
\texttt{tin\_\allowbreak{}height} & N(98, 3) in [88, 105] mm & SOURCED-VERIFY & square slices: 25 g x 0.19 g/cm3 / 12.5 mm $\to$ $\sim$100 x 100 mm (derived; tins 80-90 mm deep) \\
\texttt{dome\_\allowbreak{}height} & 0 mm & SOURCED & test25\_toast: baked under a lid (flat top) \\
\texttt{edge\_\allowbreak{}radius} & U(8 .. 12) mm & ILLUSTRATIVE & rounded tin corners \\
\texttt{slices} & 20 & SOURCED-VERIFY & toast\_pack \\
\texttt{heels} & one of 0, 2 & SOURCED & test25\_toast: most packs no heels; organic 2 \\
\texttt{crust\_\allowbreak{}thickness\_\allowbreak{}top} & U(0.8 .. 1.3) mm & ILLUSTRATIVE & thin soft lidded crust; GAP \\
\texttt{crust\_\allowbreak{}thickness\_\allowbreak{}tin} & U(0.5 .. 0.9) mm & ILLUSTRATIVE & GAP \\
\texttt{crumb\_\allowbreak{}density} & N(190, 10) in [170, 215] kg/m3 & SOURCED & nuss90 0.19 g/cm3 \\
\texttt{crust\_\allowbreak{}density} & 300 kg/m3 & ILLUSTRATIVE &  \\
\texttt{cell\_\allowbreak{}diameter\_\allowbreak{}median} & 0.5 mm & ILLUSTRATIVE & finer, more even than a roll; C-Cell values GAP \\
\texttt{cell\_\allowbreak{}number\_\allowbreak{}density} & 2500 1/cm3 & ILLUSTRATIVE &  \\
\texttt{cell\_\allowbreak{}elongation} & 1.15 & ILLUSTRATIVE &  \\
\texttt{crumb\_\allowbreak{}LCh} & (74.7, 15.5, 88.6) L*,C*,h & SOURCED-VERIFY & ishida14 74.73 / 0.37 / 15.51 \\
\texttt{crust\_\allowbreak{}LCh} & (62, 34, 70) L*,C*,h & ILLUSTRATIVE & pale lidded crust; tin loaves bottom out at L* $\sim$64.5 at 200 C (giovanelli\_2021, derived) \\
\texttt{friction\_\allowbreak{}bread} & U(0.6 .. 0.9) & ILLUSTRATIVE &  \\
\texttt{firmness\_\allowbreak{}by\_\allowbreak{}day} & 1.98 / 3.62 / 4.18 / 5.64 at day 1 / 5 / 9 / 13 (25 C, bagged) N & SOURCED & sehn20 lab white pan loaf \\
\texttt{tensile\_\allowbreak{}failure} & 0.31-0.39 (12 mm neck) at 9.2-9.7 mm N & SOURCED & nuss90 \\
\texttt{best\_\allowbreak{}before\_\allowbreak{}d} & 12-20 d & SOURCED & test25\_toast \\
\texttt{fridge\_\allowbreak{}staling\_\allowbreak{}factor} & 6 & TYPICAL & toast\_misc (popular, weak) \\
\texttt{toast\_\allowbreak{}time} & 80 .. > 150 to 'best' (per toaster) s & SOURCED & toaster26 \\
\texttt{butter\_\allowbreak{}load} & 0.068 g/cm2 & SOURCED-VERIFY & toast\_misc \\
\texttt{age\_\allowbreak{}d} & U(0.5 .. 8) d & ILLUSTRATIVE & since baking \\
\texttt{storage} & one of bag\_closed, bag\_open, fridge & ILLUSTRATIVE &  \\\bottomrule\end{longtable}}
{\scriptsize\begin{longtable}{@{}p{0.2\linewidth}p{0.27\linewidth}p{0.12\linewidth}p{0.26\linewidth}@{}}\caption{Register of the potato (params.POTATO, 36 entries).}\\\toprule Key & Value / distribution & Status & Source and note \\\midrule\endfirsthead
\caption[]{Register of the potato (params.POTATO, 36 entries). (continued)}\\\toprule Key & Value / distribution & Status & Source and note \\\midrule\endhead
\texttt{cultivar} & one of Linda, Belana, Gala & STANDARD & bsl21 \\
\texttt{cooking\_\allowbreak{}type} & one of festkochend, vorwiegend festkochend & SOURCED-VERIFY & Belana / Linda festkochend, Gala vorwiegend festkochend \\
\texttt{length} & N(85, 10) in [60, 110] mm & SOURCED-VERIFY & 75 / 65 / 99 mm (Astrix / Diamont / Santana) \\
\texttt{width} & N(55, 4) in [45, 65] mm & SOURCED-VERIFY & 48 / 45 / 58 mm \\
\texttt{thickness\_\allowbreak{}ratio} & U(0.8 .. 0.92) & SOURCED-VERIFY & thickness / width 0.81-0.89 \\
\texttt{mass\_\allowbreak{}check} & 64-167 (examples); min square mesh 35 / 30 mm (repealed HKV) g & SOURCED-VERIFY &  \\
\texttt{lumpiness} & U(0.02 .. 0.05) & ILLUSTRATIVE &  \\
\texttt{eyes} & integer 8 .. 14 & SOURCED-VERIFY & Norgold 10-12 buds; 6-9 apical (mag25) \\
\texttt{eye\_\allowbreak{}crowding} & 1.8 & ILLUSTRATIVE & apical bud density $\sim$7x lateral (mag25) \\
\texttt{phyllotaxis} & 0.4 turn & ILLUSTRATIVE & spiral arrangement sourced; the 2/5 fraction is a GAP \\
\texttt{eye\_\allowbreak{}depth} & U(0.6 .. 1.4) mm & ILLUSTRATIVE & bsl21 'flach' (ordinal) \\
\texttt{eye\_\allowbreak{}width} & U(3 .. 4) mm & ILLUSTRATIVE &  \\
\texttt{periderm\_\allowbreak{}thickness} & N(65, 10) in [46, 84] um & SOURCED & mc20 phellem \\
\texttt{lenticel\_\allowbreak{}density} & U(5 .. 12) 1/cm2 & SOURCED & mag25 \\
\texttt{vascular\_\allowbreak{}ring\_\allowbreak{}depth} & U(3 .. 6) mm & ILLUSTRATIVE & GAP (Reeve 1969 paywalled) \\
\texttt{pith\_\allowbreak{}fraction} & U(0.35 .. 0.5) & ILLUSTRATIVE & GAP \\
\texttt{pith\_\allowbreak{}branch\_\allowbreak{}width} & U(1.5 .. 2.5) mm & ILLUSTRATIVE & star of the pith; GAP \\
\texttt{density} & N(1070, 12) in [1050, 1100] kg/m3 & SOURCED-VERIFY & 1053 (kk18); SG 1.070-1.103 \\
\texttt{skin\_\allowbreak{}LCh} & (62, 30, 75) L*,C*,h & ILLUSTRATIVE & yellow skin, L* > 55 typical (SV); GAP \\
\texttt{flesh\_\allowbreak{}LCh} & (80, 30, 92) L*,C*,h & ILLUSTRATIVE & medium-yellow (bsl21), L* > 75 (SV) \\
\texttt{soil\_\allowbreak{}residue} & U(0 .. 0.3) & ILLUSTRATIVE &  \\
\texttt{friction\_\allowbreak{}skin\_\allowbreak{}rubber} & U(0.5 .. 0.8) & ILLUSTRATIVE &  \\
\texttt{E} & U(3.1 .. 5.6) MPa & SOURCED-VERIFY & potato\_mech \\
\texttt{poisson} & 0.45 & SOURCED-VERIFY & potato\_mech (Zhao 2023) \\
\texttt{fracture\_\allowbreak{}stress} & 1.13-1.40 MPa & SOURCED-VERIFY & potato\_mech \\
\texttt{thermal} & k 0.55, cp 3.43, alpha 0.15e-6 W/mK, kJ/kgK, m2/s & SOURCED & kk18 \\
\texttt{gelatinisation} & onset 61, peak 67.5, end 76-82 C & SOURCED-VERIFY & potato\_cook \\
\texttt{softening} & first order, k(100 C) 0.28-0.40, Ea 103-145 1/min, kJ/mol & SOURCED & potato\_cook (Harada) \\
\texttt{k100} & U(0.28 .. 0.4) 1/min & SOURCED & potato\_cook \\
\texttt{Ea} & 133 kJ/mol & SOURCED & potato\_cook (Harada 133-145) \\
\texttt{boil\_\allowbreak{}times} & diced 15-20, large whole up to 60, whole 20-25 (cold start) min & SOURCED & gritzer \\
\texttt{peel\_\allowbreak{}depth} & 1.1 mm & SOURCED-VERIFY & potato\_peel ($\sim$9 \% waste) \\
\texttt{storage\_\allowbreak{}temp} & 7-13 dark; < 8 sweetens C & SOURCED-VERIFY & potato\_storage \\
\texttt{greening\_\allowbreak{}days\_\allowbreak{}20C} & 1.5 d & SOURCED-VERIFY & potato\_storage 1-2.5 d under retail light \\
\texttt{age\_\allowbreak{}d} & U(2 .. 30) d & ILLUSTRATIVE &  \\
\texttt{light\_\allowbreak{}exposure\_\allowbreak{}h} & one of 0, 24, 48 h & ILLUSTRATIVE & hours under light \\\bottomrule\end{longtable}}
{\scriptsize\begin{longtable}{@{}p{0.2\linewidth}p{0.27\linewidth}p{0.12\linewidth}p{0.26\linewidth}@{}}\caption{Register of the round / truss tomato (params.TOMATO, 37 entries).}\\\toprule Key & Value / distribution & Status & Source and note \\\midrule\endfirsthead
\caption[]{Register of the round / truss tomato (params.TOMATO, 37 entries). (continued)}\\\toprule Key & Value / distribution & Status & Source and note \\\midrule\endhead
\texttt{type} & round / truss tomato, 80-120 g & STANDARD & fdc\_tomato \\
\texttt{diameter} & N(63, 4) in [52, 76] mm & SOURCED & fdc\_tomato 61-66 mm for 91-123 g \\
\texttt{height\_\allowbreak{}ratio} & U(0.75 .. 0.9) & SOURCED-VERIFY & ince16 0.75; GRIN 'slightly flattened' \\
\texttt{mass\_\allowbreak{}check} & 91 / 123 (small / medium) g & SOURCED & fdc\_tomato \\
\texttt{locules} & one of 2, 3, 4 & SOURCED & ucd\_rost: small types 2, globe 3-5 \\
\texttt{skin\_\allowbreak{}thickness} & U(120 .. 250) um & ILLUSTRATIVE & cuticle 9.9-15.7 um (bn05) + epidermis + collenchymatous hypodermis; exocarp 9.6 \% FW (meta19) \\
\texttt{pericarp\_\allowbreak{}thickness} & N(6.5, 0.8) in [5, 9] mm & SOURCED-VERIFY & meta19 5.3 mm (small); 7-12 germplasm \\
\texttt{septum\_\allowbreak{}thickness} & U(2.5 .. 3.5) mm & ILLUSTRATIVE &  \\
\texttt{columella\_\allowbreak{}diameter} & U(10 .. 16) mm & ILLUSTRATIVE & columella + placenta 10 \% FW (meta19) \\
\texttt{placenta\_\allowbreak{}width} & U(3 .. 5) mm & ILLUSTRATIVE &  \\
\texttt{seeds\_\allowbreak{}per\_\allowbreak{}fruit} & integer 50 .. 120 & SOURCED-VERIFY & 68.8 (one cultivar) \\
\texttt{seed\_\allowbreak{}length} & 3.5 mm & ILLUSTRATIVE & GAP \\
\texttt{rib\_\allowbreak{}depth} & U(0.01 .. 0.03) & ILLUSTRATIVE &  \\
\texttt{sepals} & one of 5, 6 & ILLUSTRATIVE & GAP \\
\texttt{sepal\_\allowbreak{}length} & U(12 .. 20) mm & ILLUSTRATIVE &  \\
\texttt{pedicel\_\allowbreak{}length} & one of 0, 12, 15 mm & ILLUSTRATIVE & truss tomatoes keep calyx and pedicel \\
\texttt{ripeness\_\allowbreak{}stage} & U(3.5 .. 6) 1..6 & STANDARD & cfr\_tomato colour classes \\
\texttt{skin\_\allowbreak{}hue\_\allowbreak{}by\_\allowbreak{}stage} & [118, 110, 96, 71, 52, 43] deg & DERIVED & h = atan2(1, a*/b*) at the band centres of ince16; ripe $\sim$40 deg (Radzevicius 2009, SV) \\
\texttt{skin\_\allowbreak{}L\_\allowbreak{}by\_\allowbreak{}stage} & [56, 56, 53, 49, 45, 42] L* & ILLUSTRATIVE & L* falls on ripening (rad08); GAP \\
\texttt{skin\_\allowbreak{}C\_\allowbreak{}by\_\allowbreak{}stage} & [38, 36, 40, 50, 60, 66] C* & ILLUSTRATIVE & GAP \\
\texttt{flesh\_\allowbreak{}LCh\_\allowbreak{}green} & (70, 30, 110) L*,C*,h & ILLUSTRATIVE & GAP \\
\texttt{flesh\_\allowbreak{}LCh\_\allowbreak{}breaker} & (80, 22, 85) L*,C*,h & ILLUSTRATIVE & pale interior at breaker; GAP \\
\texttt{flesh\_\allowbreak{}LCh\_\allowbreak{}red} & (48, 58, 36) L*,C*,h & ILLUSTRATIVE & GAP \\
\texttt{flesh\_\allowbreak{}density} & 1000 kg/m3 & ILLUSTRATIVE & GAP (cherry 1194, suspect) \\
\texttt{gel\_\allowbreak{}density} & 1010 kg/m3 & ILLUSTRATIVE &  \\
\texttt{friction\_\allowbreak{}skin\_\allowbreak{}rubber} & U(0.4 .. 0.7) & ILLUSTRATIVE & smooth waxy skin \\
\texttt{plate\_\allowbreak{}firmness} & initial G/P/R 15.1 / 6.0 / 4.5; average 13.2 / 4.7 / 3.1 N/mm & SOURCED-VERIFY & ince16 \\
\texttt{skin\_\allowbreak{}rupture\_\allowbreak{}2mm} & G/P/R 6.92 / 3.99 / 2.74 N & SOURCED-VERIFY & ince16 \\
\texttt{skin\_\allowbreak{}E} & 252-293 (red); failure $\le$ 25-28 MPa at 9-10 \% MPa & SOURCED & bn05 \\
\texttt{damage\_\allowbreak{}strain} & 0.03 & SOURCED-VERIFY & tomato\_mech (Li 2017 FEM) \\
\texttt{bio\_\allowbreak{}yield\_\allowbreak{}force} & 2.57 N & SOURCED-VERIFY & tomato\_mech \\
\texttt{stage\_\allowbreak{}rate\_\allowbreak{}20C} & U(0.35 .. 0.7) stage/d & SOURCED-VERIFY & breaker $\to$ full colour 7-14 d (patent) \\
\texttt{chilling\_\allowbreak{}threshold} & 12.5 C & SOURCED & wflo\_tomato 12.5-12.8 \\
\texttt{fridge\_\allowbreak{}volatile\_\allowbreak{}loss} & 65 after 8 d at 5 C \% & SOURCED & zhang16 \\
\texttt{water} & 94.5 \% & SOURCED & fdc\_tomato \\
\texttt{slice\_\allowbreak{}thickness} & 6-13 mm & SOURCED & tomato\_howto \\
\texttt{blanch\_\allowbreak{}s} & 30-45 (20-60) s & SOURCED & tomato\_howto \\\bottomrule\end{longtable}}
{\scriptsize\begin{longtable}{@{}p{0.2\linewidth}p{0.27\linewidth}p{0.12\linewidth}p{0.26\linewidth}@{}}\caption{Register of the yellow onion (params.ONION, 32 entries).}\\\toprule Key & Value / distribution & Status & Source and note \\\midrule\endfirsthead
\caption[]{Register of the yellow onion (params.ONION, 32 entries). (continued)}\\\toprule Key & Value / distribution & Status & Source and note \\\midrule\endhead
\texttt{type} & yellow storage onion (Haushaltszwiebel), cured & STANDARD & hb66\_onion \\
\texttt{diameter} & N(65, 7) in [50, 85] mm & SOURCED & hb66\_onion medium 51-83 mm; bah04 51.7 \\
\texttt{shape\_\allowbreak{}index} & U(0.85 .. 1.1) & SOURCED-VERIFY & height / diameter: bah04 57.1 / 51.7; flat-round German types \\
\texttt{mass\_\allowbreak{}check} & 78.7 $\pm$ 24 (51.7 mm) g & SOURCED-VERIFY & bah04 \\
\texttt{n\_\allowbreak{}fleshy} & one of 7, 8 & SOURCED-VERIFY & kim\_heb20 7.6 $\pm$ 0.2 \\
\texttt{scale\_\allowbreak{}thickness\_\allowbreak{}outer} & U(3.5 .. 4.5) mm & ILLUSTRATIVE & kim\_heb20 5.5 mm (bulk mean), wu25 $\sim$3 mm; constrained by count x radius \\
\texttt{scale\_\allowbreak{}thickness\_\allowbreak{}inner} & U(2 .. 2.8) mm & ILLUSTRATIVE & inner scales thinner; GAP \\
\texttt{scale\_\allowbreak{}gap} & 0.05 mm & ILLUSTRATIVE & epidermis interface between scales \\
\texttt{n\_\allowbreak{}dry} & one of 2, 3 & ILLUSTRATIVE & GAP (ordinal skin-thickness classes only) \\
\texttt{dry\_\allowbreak{}thickness} & U(0.08 .. 0.15) mm & ILLUSTRATIVE & GAP \\
\texttt{plate\_\allowbreak{}diameter} & U(15 .. 22) mm & ILLUSTRATIVE & GAP \\
\texttt{plate\_\allowbreak{}thickness} & U(4 .. 6) mm & ILLUSTRATIVE & GAP \\
\texttt{n\_\allowbreak{}roots} & integer 15 .. 30 & ILLUSTRATIVE &  \\
\texttt{root\_\allowbreak{}length} & U(3 .. 10) mm & ILLUSTRATIVE & trimmed root crown \\
\texttt{neck\_\allowbreak{}length} & U(5 .. 20) mm & ILLUSTRATIVE &  \\
\texttt{neck\_\allowbreak{}diameter} & U(7.4 .. 10.9) mm & SOURCED-VERIFY & onion\_misc (cured neck, Indian cultivars) \\
\texttt{flesh\_\allowbreak{}density} & 1040 kg/m3 & SOURCED-VERIFY & bah04 bulb 1.04 g/cm3 \\
\texttt{dry\_\allowbreak{}density} & 400 kg/m3 & ILLUSTRATIVE & papery skin \\
\texttt{plate\_\allowbreak{}density} & 1050 kg/m3 & ILLUSTRATIVE &  \\
\texttt{skin\_\allowbreak{}LCh} & (52, 45, 65) L*,C*,h & ILLUSTRATIVE & golden brown; GAP \\
\texttt{flesh\_\allowbreak{}LCh} & (88, 12, 95) L*,C*,h & ILLUSTRATIVE & white to pale yellow; GAP \\
\texttt{friction\_\allowbreak{}skin\_\allowbreak{}rubber} & U(0.6 .. 1) & SOURCED-VERIFY & bah04 0.67-1.34 \\
\texttt{bulb\_\allowbreak{}E} & 2.03-5.45, nu 0.26-0.45 MPa & SOURCED-VERIFY & jafari16 \\
\texttt{tissue\_\allowbreak{}E} & 2.2 skin on / 1.4 off MPa & SOURCED & wu25 \\
\texttt{epidermis\_\allowbreak{}E} & 3.5-8.0 MPa & SOURCED & wei01 \\
\texttt{crushing\_\allowbreak{}load} & 400-780 N & SOURCED-VERIFY & bah04 \\
\texttt{rolling\_\allowbreak{}angle} & 20-31 on base, 14-23 on side deg & SOURCED-VERIFY & bah04 \\
\texttt{lf\_\allowbreak{}window} & peak $\sim$1, gone $\sim$10 min & SOURCED & lokke12 \\
\texttt{storage} & 0, 65-75, 6-9 months; household 3-4 weeks C, \%RH & SOURCED & hb66\_onion \\
\texttt{dice\_\allowbreak{}spacing} & mince $\le$3, fine 3-6, medium 6-13, coarse 13-19 mm & SOURCED & onion\_howto \\
\texttt{caramelise\_\allowbreak{}min} & 30-45 min & TYPICAL & onion\_misc \\
\texttt{age\_\allowbreak{}d} & U(3 .. 40) d & ILLUSTRATIVE &  \\\bottomrule\end{longtable}}
{\scriptsize\begin{longtable}{@{}p{0.2\linewidth}p{0.27\linewidth}p{0.12\linewidth}p{0.26\linewidth}@{}}\caption{Register of the lemon (params.LEMON, 31 entries).}\\\toprule Key & Value / distribution & Status & Source and note \\\midrule\endfirsthead
\caption[]{Register of the lemon (params.LEMON, 31 entries). (continued)}\\\toprule Key & Value / distribution & Status & Source and note \\\midrule\endhead
\texttt{cultivar} & one of Primofiori, Verna & SOURCED & agu20 \\
\texttt{diameter} & N(61, 3) in [55, 67] mm & SOURCED & agu20 55-67 mm \\
\texttt{shape\_\allowbreak{}index} & U(1.3 .. 1.55) & SOURCED & length / diameter from agu20 ED/PD 0.6-0.8 (body without nipple) \\
\texttt{nipple\_\allowbreak{}height} & U(3 .. 7) mm & ILLUSTRATIVE & stylar mammilla; GAP \\
\texttt{neck\_\allowbreak{}height} & U(1.5 .. 4) mm & ILLUSTRATIVE & stem-end neck; GAP \\
\texttt{mass\_\allowbreak{}check} & 100-184 (agu20); 166 (jen22\_lemon) g & SOURCED &  \\
\texttt{peel\_\allowbreak{}thickness} & N(5.6, 0.5) in [4.9, 7.1] mm & SOURCED & agu20, jen22\_lemon \\
\texttt{flavedo\_\allowbreak{}thickness} & U(1 .. 1.4) mm & SOURCED-VERIFY & jen22\_lemon $\sim$1 mm (ben26: $\sim$46 \% of peel) \\
\texttt{core\_\allowbreak{}diameter} & U(4 .. 7) mm & ILLUSTRATIVE & GAP \\
\texttt{segments} & one of 9, 10 & SOURCED & agu20 8.8-10.2 \\
\texttt{membrane\_\allowbreak{}thickness} & 0.1 mm & ILLUSTRATIVE & GAP \\
\texttt{seeds} & integer 1 .. 8 & SOURCED & agu20 1.3-11.3, lemon\_misc 1.3-2.3 \\
\texttt{seed\_\allowbreak{}length} & 8 mm & TYPICAL &  \\
\texttt{oil\_\allowbreak{}gland\_\allowbreak{}density} & 125 1/cm2 & SOURCED & ben26 1.25 /mm2 \\
\texttt{peel\_\allowbreak{}density} & 818 kg/m3 & SOURCED & jen22\_lemon \\
\texttt{albedo\_\allowbreak{}density} & 390 kg/m3 & SOURCED & ben26 \\
\texttt{pulp\_\allowbreak{}density} & 1030 kg/m3 & DERIVED & juice $\sim$1.03 g/mL from USDA portion weights (lemon\_misc) \\
\texttt{peel\_\allowbreak{}LCh} & (66, 46, 93) L*,C*,h & SOURCED & agu20 L* 64-68, b* 42-49, hue 88-99 \\
\texttt{albedo\_\allowbreak{}LCh} & (92, 12, 95) L*,C*,h & ILLUSTRATIVE & GAP \\
\texttt{pulp\_\allowbreak{}LCh} & (82, 35, 98) L*,C*,h & ILLUSTRATIVE & pale yellow; GAP \\
\texttt{membrane\_\allowbreak{}LCh} & (88, 18, 95) L*,C*,h & ILLUSTRATIVE & GAP \\
\texttt{friction\_\allowbreak{}peel\_\allowbreak{}rubber} & U(0.5 .. 0.8) & ILLUSTRATIVE &  \\
\texttt{juice\_\allowbreak{}fraction} & U(30 .. 45) \% & SOURCED & agu20 31.7-44.3 (EU min 20) \\
\texttt{juice\_\allowbreak{}of\_\allowbreak{}one\_\allowbreak{}lemon} & 48 g & SOURCED & lemon\_misc (USDA) \\
\texttt{acidity} & 56-75 g/L citric & SOURCED & agu20 \\
\texttt{puncture\_\allowbreak{}force} & 41.2 N & SOURCED & ben26 5 mm probe \\
\texttt{cutting\_\allowbreak{}force} & 36.3 N & SOURCED & ben26 craft knife \\
\texttt{compression\_\allowbreak{}force} & 165.8 at 16 mm N & SOURCED & ben26 \\
\texttt{chilling\_\allowbreak{}threshold} & 10 C & SOURCED & lemon\_misc (Eureka) \\
\texttt{wedge\_\allowbreak{}mass} & 5.9 g & SOURCED & lemon\_misc (USDA) \\
\texttt{bitter\_\allowbreak{}juice\_\allowbreak{}h} & 1e+06 h & ILLUSTRATIVE & lemon juice does not turn bitter like navel juice \\\bottomrule\end{longtable}}
{\scriptsize\begin{longtable}{@{}p{0.2\linewidth}p{0.27\linewidth}p{0.12\linewidth}p{0.26\linewidth}@{}}\caption{Register of the greenhouse cucumber (params.CUCUMBER, 28 entries).}\\\toprule Key & Value / distribution & Status & Source and note \\\midrule\endfirsthead
\caption[]{Register of the greenhouse cucumber (params.CUCUMBER, 28 entries). (continued)}\\\toprule Key & Value / distribution & Status & Source and note \\\midrule\endhead
\texttt{type} & greenhouse long cucumber (Schlangengurke), parthenocarpic & STANDARD & dewiki\_sg \\
\texttt{length} & N(320, 25) in [250, 380] mm & SOURCED & unece\_cuc 25-30 cm minimum; sakata11 \\
\texttt{diameter} & N(42, 3) in [36, 50] mm & TYPICAL & derived: 400 g / 32 cm / 950 kg/m3 $\to$ $\sim$41 mm \\
\texttt{mass\_\allowbreak{}check} & 300-500 (dewiki\_sg), typical 350 (cuc\_misc) g & SOURCED &  \\
\texttt{neck\_\allowbreak{}ratio} & U(0.8 .. 0.9) & ILLUSTRATIVE & stem end narrower \\
\texttt{ribs} & 3 & ILLUSTRATIVE & faint ribs over the three carpels \\
\texttt{rib\_\allowbreak{}depth} & U(0.005 .. 0.012) & ILLUSTRATIVE &  \\
\texttt{bend\_\allowbreak{}radius} & U(1000 .. 4000) mm & ILLUSTRATIVE & within Class I ($\ge$ $\sim$130 mm, unece\_cuc) \\
\texttt{skin\_\allowbreak{}thickness} & U(1 .. 1.6) mm & SOURCED-VERIFY & cuc\_misc peelable exocarp 1-2 mm \\
\texttt{cavity\_\allowbreak{}fraction} & U(0.45 .. 0.55) & ILLUSTRATIVE & seed cavity / fruit radius; GAP \\
\texttt{core\_\allowbreak{}diameter} & U(3 .. 5) mm & ILLUSTRATIVE &  \\
\texttt{seeds\_\allowbreak{}per\_\allowbreak{}row} & integer 25 .. 40 & ILLUSTRATIVE & parthenocarpic: soft immature ovules \\
\texttt{seed\_\allowbreak{}length} & U(4 .. 6) mm & ILLUSTRATIVE & < mature 6.4-9.1 mm \\
\texttt{density} & 950 kg/m3 & SOURCED-VERIFY & disect\_cuc; water 95.2 \% (USDA) \\
\texttt{skin\_\allowbreak{}LCh} & (38, 34, 127) L*,C*,h & SOURCED-VERIFY & cuc\_misc a* -20.6, b* 27.5 $\to$ C 34, h 127; L* GAP \\
\texttt{flesh\_\allowbreak{}LCh} & (85, 20, 112) L*,C*,h & ILLUSTRATIVE & pale whitish green; GAP \\
\texttt{friction\_\allowbreak{}skin\_\allowbreak{}rubber} & U(0.4 .. 0.7) & ILLUSTRATIVE & glossy waxy skin \\
\texttt{round\_\allowbreak{}thickness} & 5 mm & SOURCED & disect\_cuc (SliceIt! 5 mm slices) \\
\texttt{puncture\_\allowbreak{}3mm} & flesh 8.7, placenta 2.8, skin 9.7 N & SOURCED & sakata11 \\
\texttt{E\_\allowbreak{}apparent} & 0.85-1.39 (slice) | 2.5 (DiSECt) MPa & SOURCED-VERIFY & jah16; disect\_cuc \\
\texttt{failure} & 0.53-0.85 at 0.45-0.53 strain MPa, 1 & SOURCED-VERIFY & jah16 \\
\texttt{firmness\_\allowbreak{}loss\_\allowbreak{}9d} & 49 \% & SOURCED-VERIFY & jah16 \\
\texttt{storage\_\allowbreak{}optimum} & 10-12.5, 95, < 14 d C, \%RH & SOURCED & ucd\_cuc \\
\texttt{chilling} & < 10 C after 1-3 d C, d & SOURCED & ucd\_cuc \\
\texttt{weight\_\allowbreak{}loss\_\allowbreak{}15d} & wrapped 0.66, unwrapped 11.1 (12 C) \% & SOURCED-VERIFY & cuc\_misc \\
\texttt{salting\_\allowbreak{}min} & 30 min & SOURCED & cuc\_misc (wikibooks) \\
\texttt{age\_\allowbreak{}d} & U(1 .. 10) d & ILLUSTRATIVE &  \\
\texttt{storage} & one of fridge\_wrapped, counter\_wrapped, fridge\_unwrapped & ILLUSTRATIVE & home fridge 4 C is in the chilling range \\\bottomrule\end{longtable}}
{\scriptsize\begin{longtable}{@{}p{0.2\linewidth}p{0.27\linewidth}p{0.12\linewidth}p{0.26\linewidth}@{}}\caption{Register of the butter croissant (params.CROISSANT, 21 entries).}\\\toprule Key & Value / distribution & Status & Source and note \\\midrule\endfirsthead
\caption[]{Register of the butter croissant (params.CROISSANT, 21 entries). (continued)}\\\toprule Key & Value / distribution & Status & Source and note \\\midrule\endhead
\texttt{product} & Buttercroissant (butter the only fat) & STANDARD & leitsaetze\_fb \\
\texttt{length} & N(135, 8) in [115, 160] mm & ILLUSTRATIVE & finished dimensions GAP; chosen for $\sim$180 cm3 = 57 g x 3.13 cm3/g (ste25, USDA medium) \\
\texttt{width} & N(68, 5) in [55, 80] mm & ILLUSTRATIVE & see length \\
\texttt{height} & N(42, 4) in [34, 50] mm & ILLUSTRATIVE & see length \\
\texttt{mass\_\allowbreak{}check} & 42 / 57 / 67 (small / medium / large) g & SOURCED & croissant\_misc (USDA) \\
\texttt{curvature\_\allowbreak{}radius} & one of 70, 90, 100000 mm & ILLUSTRATIVE & crescent or straight (fr: butter = straight, SV) \\
\texttt{steps} & integer 5 .. 7 & ILLUSTRATIVE & visible rolls; rolled 3.5-4 turns (croissant\_misc) \\
\texttt{step\_\allowbreak{}relief} & U(0.05 .. 0.1) & ILLUSTRATIVE &  \\
\texttt{rolling\_\allowbreak{}turns} & U(3.5 .. 4) & SOURCED & croissant\_misc \\
\texttt{butter\_\allowbreak{}layers} & 27 & SOURCED & ste25: 3 single folds \\
\texttt{butter\_\allowbreak{}on\_\allowbreak{}flour} & U(50 .. 64) \% & SOURCED & ste25, recipes; legal min 30 (leitsaetze\_fb) \\
\texttt{crust\_\allowbreak{}thickness} & U(0.8 .. 1.4) mm & ILLUSTRATIVE & GAP \\
\texttt{crust\_\allowbreak{}density} & 450 kg/m3 & ILLUSTRATIVE &  \\
\texttt{crumb\_\allowbreak{}density} & 300 kg/m3 & DERIVED & 1 / 3.13 cm3/g (ste25) = 0.32 g/cm3 whole \\
\texttt{crust\_\allowbreak{}LCh} & (52, 48, 66) L*,C*,h & ILLUSTRATIVE & deep golden, egg-washed; GAP \\
\texttt{crumb\_\allowbreak{}LCh} & (84, 26, 88) L*,C*,h & ILLUSTRATIVE & pale yellow (butter); GAP \\
\texttt{flakes} & integer 40 .. 120 & ILLUSTRATIVE & loose flakes; GAP \\
\texttt{friction\_\allowbreak{}crust} & U(0.4 .. 0.7) & ILLUSTRATIVE &  \\
\texttt{firmness} & 5.16 (1 h) $\to$ 19.2 (24 h), TPA 25 mm 50 \% N & SOURCED & ste25 \\
\texttt{reheat} & 180, 8 from frozen; refresh 160-180, 3-5 (typical) C, min & SOURCED & croissant\_misc \\
\texttt{age\_\allowbreak{}h} & one of 1, 3, 8, 24 h & ILLUSTRATIVE &  \\\bottomrule\end{longtable}}
{\scriptsize\begin{longtable}{@{}p{0.2\linewidth}p{0.27\linewidth}p{0.12\linewidth}p{0.26\linewidth}@{}}\caption{Register of the baguette (params.BAGUETTE, 42 entries).}\\\toprule Key & Value / distribution & Status & Source and note \\\midrule\endfirsthead
\caption[]{Register of the baguette (params.BAGUETTE, 42 entries). (continued)}\\\toprule Key & Value / distribution & Status & Source and note \\\midrule\endhead
\texttt{product} & baguette de tradition francaise (flour, water, yeast/levain, salt) & STANDARD & wp\_baguette \\
\texttt{length} & N(600, 30) in [540, 660] mm & SOURCED & wp\_baguette 55-65 cm \\
\texttt{width} & N(55, 4) in [45, 62] mm & SOURCED & wp\_baguette 4-6 cm \\
\texttt{height} & N(42, 4) in [32, 50] mm & SOURCED & wp\_baguette 3-5 cm \\
\texttt{mass\_\allowbreak{}check} & $\sim$250 (250-300 tradition) g & SOURCED & wp\_baguette \\
\texttt{base\_\allowbreak{}edge\_\allowbreak{}height} & U(4 .. 7) mm & ILLUSTRATIVE &  \\
\texttt{lumpiness} & U(0.01 .. 0.025) & ILLUSTRATIVE &  \\
\texttt{n\_\allowbreak{}scores} & one of 5, 6, 7 & SOURCED & wp\_grigne: 5 classic; 5-7 typical for 55-65 cm \\
\texttt{score\_\allowbreak{}length\_\allowbreak{}mm} & U(100 .. 125) mm & TYPICAL & $\sim$10-12 cm, overlapping $\sim$1/3 \\
\texttt{score\_\allowbreak{}axis\_\allowbreak{}angle} & U(10 .. 20) deg & TYPICAL &  \\
\texttt{score\_\allowbreak{}length\_\allowbreak{}frac} & 1 & ILLUSTRATIVE & per score: set by score\_length\_mm \\
\texttt{score\_\allowbreak{}width} & U(9 .. 14) mm & ILLUSTRATIVE & opened grigne \\
\texttt{score\_\allowbreak{}depth} & U(5 .. 9) mm & SOURCED-VERIFY & wp\_grigne $\sim$1 cm cut \\
\texttt{crust\_\allowbreak{}thickness\_\allowbreak{}top} & N(1.9, 0.3) in [1.4, 3.5] mm & SOURCED & dv2012 1.7 mean (zou2024: 3-4, SV) \\
\texttt{crust\_\allowbreak{}thickness\_\allowbreak{}score} & U(0.8 .. 1.2) mm & ILLUSTRATIVE &  \\
\texttt{crust\_\allowbreak{}thickness\_\allowbreak{}bottom} & U(1.6 .. 2.4) mm & ILLUSTRATIVE &  \\
\texttt{transition\_\allowbreak{}thickness} & U(1 .. 2) mm & ILLUSTRATIVE &  \\
\texttt{crumb\_\allowbreak{}density} & N(200, 20) in [160, 250] kg/m3 & DERIVED & $\sim$250 g / $\sim$0.9 L bulk (dimensions, wp\_baguette) minus a denser crust \\
\texttt{crust\_\allowbreak{}density} & 380 kg/m3 & ILLUSTRATIVE &  \\
\texttt{transition\_\allowbreak{}density} & 260 kg/m3 & ILLUSTRATIVE &  \\
\texttt{crumb\_\allowbreak{}porosity} & U(0.75 .. 0.85) & ILLUSTRATIVE & open irregular crumb; tomography GAP \\
\texttt{cell\_\allowbreak{}diameter\_\allowbreak{}median} & U(1.3 .. 1.9) mm & ILLUSTRATIVE & 'irregular voids of various sizes' \\
\texttt{cell\_\allowbreak{}diameter\_\allowbreak{}gsd\_\allowbreak{}log} & 0.75 & ILLUSTRATIVE & wide spread: open crumb \\
\texttt{cell\_\allowbreak{}number\_\allowbreak{}density} & 260 1/cm3 & ILLUSTRATIVE &  \\
\texttt{cell\_\allowbreak{}elongation} & U(1.1 .. 1.3) & ILLUSTRATIVE &  \\
\texttt{crust\_\allowbreak{}L} & N(52, 3) in [44, 58] L* & SOURCED-VERIFY & zou2024 powder 57.5; intact darker \\
\texttt{crust\_\allowbreak{}C} & 36.4 C* & SOURCED-VERIFY & zou2024 a* 10.5, b* 34.9 \\
\texttt{crust\_\allowbreak{}h} & 73.3 deg & SOURCED-VERIFY & zou2024 \\
\texttt{crumb\_\allowbreak{}L} & 70 L* & SOURCED-VERIFY & zou2024 66.4 (powder) \\
\texttt{crumb\_\allowbreak{}C} & 15.8 C* & SOURCED-VERIFY & zou2024 \\
\texttt{crumb\_\allowbreak{}h} & 96 deg & SOURCED-VERIFY & zou2024 \\
\texttt{variant} & one of plain & STANDARD &  \\
\texttt{slice\_\allowbreak{}height\_\allowbreak{}frac} & 0.42 & TYPICAL & split lengthwise for sandwiches \\
\texttt{friction\_\allowbreak{}crust\_\allowbreak{}rubber} & U(0.5 .. 0.8) & ILLUSTRATIVE &  \\
\texttt{age\_\allowbreak{}h} & one of 0.5, 2, 6, 20 h & ILLUSTRATIVE &  \\
\texttt{storage} & one of paper\_bag, open & ILLUSTRATIVE & baguette\_misc \\
\texttt{crumb\_\allowbreak{}poisson} & 0.103 & SOURCED & lagrain\_2012 (bread) \\
\texttt{crust\_\allowbreak{}water\_\allowbreak{}critical} & 10.1 \% & SOURCED & hirte\_2012 \\
\texttt{plastic\_\allowbreak{}bag\_\allowbreak{}factor} & 4 & ILLUSTRATIVE &  \\
\texttt{crust\_\allowbreak{}break\_\allowbreak{}force} & 3 N & SOURCED & hirte\_2012 \\
\texttt{crumb\_\allowbreak{}squash\_\allowbreak{}strain} & 0.3 & ILLUSTRATIVE &  \\
\texttt{bias\_\allowbreak{}slice\_\allowbreak{}mm} & 12 mm & TYPICAL & crostini 1-1.5 cm at 30-45 deg \\\bottomrule\end{longtable}}
{\scriptsize\begin{longtable}{@{}p{0.2\linewidth}p{0.27\linewidth}p{0.12\linewidth}p{0.26\linewidth}@{}}\caption{Register of the bell pepper (params.PEPPER, 27 entries).}\\\toprule Key & Value / distribution & Status & Source and note \\\midrule\endfirsthead
\caption[]{Register of the bell pepper (params.PEPPER, 27 entries). (continued)}\\\toprule Key & Value / distribution & Status & Source and note \\\midrule\endhead
\texttt{type} & blocky sweet pepper (Paprika) & STANDARD & pepper\_misc (EU 2023/2429) \\
\texttt{colour} & one of red, yellow, green & ILLUSTRATIVE & retail mix \\
\texttt{diameter} & N(78, 4) in [68, 90] mm & SOURCED & he22 77.6 mm \\
\texttt{length} & N(83, 6) in [70, 100] mm & SOURCED & he22 81.3-85.3 mm \\
\texttt{mass\_\allowbreak{}check} & 237 $\pm$ 37 (he22), 230-245 German retail (alt20) g & SOURCED &  \\
\texttt{lobes} & one of 3, 4 & SOURCED & sha26: 2-4, mostly 3 \\
\texttt{lobe\_\allowbreak{}depth} & U(0.04 .. 0.08) & ILLUSTRATIVE &  \\
\texttt{wall\_\allowbreak{}thickness} & N(7.2, 0.6) in [5.5, 8.5] mm & SOURCED & he22 7.5 red / 6.9 orange \\
\texttt{wall\_\allowbreak{}density} & 1000 kg/m3 & TYPICAL & 92-94 \% water \\
\texttt{rib\_\allowbreak{}depth} & U(6 .. 12) mm & ILLUSTRATIVE & GAP \\
\texttt{rib\_\allowbreak{}thickness} & U(1.2 .. 2) mm & ILLUSTRATIVE & GAP \\
\texttt{placenta\_\allowbreak{}diameter} & U(20 .. 28) mm & ILLUSTRATIVE & GAP \\
\texttt{placenta\_\allowbreak{}height} & U(22 .. 32) mm & ILLUSTRATIVE & GAP \\
\texttt{seeds} & integer 40 .. 140 & SOURCED-VERIFY & pepper\_misc 38-142 \\
\texttt{seed\_\allowbreak{}diameter} & 3.5 mm & ILLUSTRATIVE & 1000-seed weight 4-6 g (sha26); size GAP \\
\texttt{stem\_\allowbreak{}length} & U(10 .. 25) mm & ILLUSTRATIVE & cut peduncle (EU: neatly cut) \\
\texttt{red\_\allowbreak{}LCh} & (38.8, 41.9, 41.6) L*,C*,h & SOURCED & bar20 \\
\texttt{yellow\_\allowbreak{}LCh} & (61.4, 50.2, 84.3) L*,C*,h & SOURCED & bar20 \\
\texttt{green\_\allowbreak{}LCh} & (33.7, 19, 116.9) L*,C*,h & SOURCED & bar20 \\
\texttt{friction\_\allowbreak{}skin\_\allowbreak{}rubber} & U(0.35 .. 0.6) & ILLUSTRATIVE & glossy waxy skin; robots report slip \\
\texttt{firmness\_\allowbreak{}whole} & 43.5 N & SOURCED & pepper\_misc (8 mm probe) \\
\texttt{bite\_\allowbreak{}firmness} & red 11.6, orange 6.9; -81 \% in 15 d at 20 C N & SOURCED & he22 \\
\texttt{chilling\_\allowbreak{}threshold} & 7 C & SOURCED & pepper\_misc \\
\texttt{weight\_\allowbreak{}loss} & 2.9-4.8 at 17 C / 55 \% RH; shrivel from 6-8 \% loss g/d & SOURCED & alt20 \\
\texttt{shelf\_\allowbreak{}life\_\allowbreak{}d} & 12-28 d & SOURCED & bar20 \\
\texttt{ring\_\allowbreak{}mass} & 10 g & SOURCED & pepper\_misc (USDA) \\
\texttt{age\_\allowbreak{}d} & U(1 .. 10) d & ILLUSTRATIVE &  \\\bottomrule\end{longtable}}
\section{Sources}
{\scriptsize\begin{longtable}{@{}p{0.16\linewidth}p{0.8\linewidth}@{}}\caption{Sources held in food/params.py (SOURCES).}\\\toprule Key & Source \\\midrule\endfirsthead
\caption[]{Sources held in food/params.py (SOURCES). (continued)}\\\toprule Key & Source \\\midrule\endhead
\texttt{narushin2021} & Narushin V.G., Romanov M.N., Griffin D.K. (2021) Egg and math: introducing a universal formula for egg shape. Ann. N.Y. Acad. Sci. 1505:169-177 (Hügelschäffer egg model). \\
\texttt{eu589\_\allowbreak{}2008} & Commission Regulation (EC) No 589/2008, marketing standards for eggs: weight grades S<53, M 53-63, L 63-73, XL$\ge$73 g; Class A air space height $\le$6 mm, 'extra' $\le$4 mm. \\
\texttt{icf11\_\allowbreak{}eggshell} & ICF11 conference paper on eggshell fracture, \url{gruppofrattura.it/ocs/index.php/ICF/ICF11/paper/download/10244/9618:} failure at 34$\pm$6 N (equatorial compression) and 41$\pm$8 N (axial, across the poles). Species to verify in the full text. \\
\texttt{nedomova2004} & Hen-egg quasi-static compression, Acta Univ. Agric. Silvic. Mendel. Brun. 2004 \url{(actavia.mendelu.cz/pdfs/acu/2004/05/07.pdf):} model with shell thickness 0.35 mm and Young's modulus 15,000 MPa. \\
\texttt{reissner1946} & Reissner E. (1946) Stresses and small displacements of shallow spherical shells. J. Math. Phys. 25: point-load stiffness of a shallow spherical shell k = 4 E t\^{}2 / (R sqrt(3(1-nu\^{}2))). \\
\texttt{apple\_\allowbreak{}E\_\allowbreak{}quasistatic} & Quasi-static equivalent Young's modulus of apple flesh: Jonagold 4.24$\pm$0.96, Joly Red 5.09$\pm$1.27, Kanzi 7.82$\pm$0.41 MPa (impact study, found via \url{researchportal.be/en/node/4347597;} citation to verify). \\
\texttt{tjaf2005} & Turkish J. Agric. For. 29(6) 2005 (tar-29-6-6-0411-8): apple flesh modulus falls with storage, 1.68$\to$0.51 MPa (Williams), 1.45$\to$0.88 (Starkspur Golden Delicious), 1.51$\to$1.1 (Starking). \\
\texttt{fsct2008} & J. Food Sci. Technol. (Iran) 5(18):37-44 (2008): bruise volume at 410 and 1210 mJ impacts; Golden Delicious about 1 \% of fruit volume at 410 mJ; bruise volume rises with impact energy. \\
\texttt{hertz} & Hertz contact theory (Johnson, Contact Mechanics, 1985): sphere-on-flat force, pressure, impact. \\
\texttt{avo\_\allowbreak{}au\_\allowbreak{}stages} & Avocados Australia, 'Avocado Stages of Ripeness chart for Hass' (2017): stages 1 (hard) .. 5 (ripe/eat now); penetrometer (11.1 mm tip, 8 mm into skinless flesh) >10 kgf at stage 1, 0.4-0.6 kgf at 'eat now'; firmness is a better ripeness index than skin colour. \\
\texttt{avo\_\allowbreak{}firmness\_\allowbreak{}12d} & Hass flesh firmness falling from 75.43 to 2.63 N over 12 days of ripening (reported in a ripening study surfaced by literature search; citation to verify). \\
\texttt{wiki\_\allowbreak{}apple} & \url{https://en.wikipedia.org/wiki/Apple} and \url{https://en.wikipedia.org/wiki/Pome:} flesh from the receptacle/hypanthium; usually 5 carpels with 1-2 seeds; endocarp = leathery core case. \\
\texttt{wiki\_\allowbreak{}gala} & \url{https://en.wikipedia.org/wiki/Gala_(apple):} 70-80 mm, striped orange-red over yellow, density 0.86. \\
\texttt{konarska2006} & Hort. Sci. (Prague) 2006 \url{(agriculturejournals.cz/pdfs/hor/2006/03/04.pdf):} apple skin 33.3-73.1 um. \\
\texttt{ishs\_\allowbreak{}1099\_\allowbreak{}97} & ISHS Acta Hortic. 1099\_97: lenticel density 4.3-12.6 per cm2, mean 7.11. \\
\texttt{frontiers2015} & Front. Plant Sci. 2015, 6:679 (fpls.2015.00679): porosity 10.6 \% centre .. 26.4 \% periphery; 10 primary bundles (5 sepal + 5 petal), carpellary 10 ventral + 5 dorsal bundles. \\
\texttt{uspto} & USPTO apple plant patents PP29089, PP20028, PP21710, PP23862, PP26766: core 24.7-36.7 mm, core length 41.1 mm, locules 15.2 x 7.9 / 16.6 x 9.9 mm, stems 10.5-28.2 x 2.7-3 mm, stem cavity 8-25 mm deep, calyx basin 3-17 mm deep. \\
\texttt{sciRep2021} & Sci. Rep. 2021 (s41598-021-02143-z): apple seed length 7.69-9.43 mm, width 4.0-4.9 mm. \\
\texttt{msu2025} & MSU Extension apple maturity report 2025: Gala $\sim$18 lb Magness-Taylor at harvest. \\
\texttt{fmach\_\allowbreak{}porosity} & \url{openpub.fmach.it} 10449/23207: apple porosity 17 \% (Jazz) .. 29.8 \% (Golden Delicious). \\
\texttt{jets\_\allowbreak{}skin\_\allowbreak{}tensile} & J. Eng. Technol. Sci. \url{(journals.itb.ac.id/index.php/jets/article/view/1231):} skin 2.56 MPa. \\
\texttt{pmc7022590} & PMC7022590: Golden Delicious slices browning dE 2.38 / 2.68 / 4.05 at 0 / 30 / 60 min. \\
\texttt{atk\_\allowbreak{}cut\_\allowbreak{}apple} & America's Test Kitchen 'How to cut an apple without coring': four cuts beside the core, slices 6 mm (matchsticks) / 13 mm (dice). \\
\texttt{usda\_\allowbreak{}egg} & USDA AMS Egg-Grading Manual (ams.usda.gov/.../Egg Grading Manual.pdf): cuticle (not 'bloom'), palisade/'spongy' and mammillary layers, two membranes, air cell grades AA $\le$3.2 mm, A $\le$4.8 mm; albumen = outer thin 23 \%, thick 57 \%, inner thin 17 \%, chalaziferous 3 \%; yolk $\sim$31 \%; check vs leaker. \\
\texttt{fr} & Front. Immunol. 2022, 13:838525: cuticle 0.5-12 um, vertical crystal 3-8 um, palisade $\sim$300 um, mammillary $\sim$100 um, membranes $\sim$70 um; pores $\sim$2x denser at the blunt pole. \\
\texttt{ir} & ScienceDirect S2666469024000101: cuticle 5.6-12.5 um, mean 7.85 um (search snippet). \\
\texttt{por} & Rev. Bras. Cienc. Avic. (scielo ZBg3hyp6Swr9HgR9XSQSQsN): 10,000-20,000 pores (snippet). \\
\texttt{ubc} & UBC open collections 1.0104002: outer shell membrane $\sim$3x thicker than inner. \\
\texttt{kop} & J. Process. Energy Agric. (aseestant.ceon.rs 32280): 58.05 x 45.46 mm, SI 78.25, 65.89 g, shell 12.76 \%; compression 32.41 N mean, Hisex Brown 36.43 N short axis / 34.92 N long axis, range 10.03-58.96 N. \\
\texttt{alt} & Effect of egg shape index on mechanical properties (unis.gop.edu.tr): SI classes sharp <72, normal 72-76, round >76; highest breaking force along the polar axis at the sharp end. \\
\texttt{br} & Redalyc 395940112007: shell 8-11 \%, albumen 56-61 \%, yolk 27-32 \% of egg mass. \\
\texttt{sev} & Acta Univ. Agric. Mendel. 2010/01/19: eggshell E 41$\pm$3.9 GPa radial, 43$\pm$4.1 GPa tangential. \\
\texttt{sb} & Acta Univ. Agric. Mendel. 2004/05/07: single egg failed $\sim$49 N, linear to failure, shell 0.35 mm, E 15 GPa. \\
\texttt{gr} & Grashorn, Lohmann Information 49 (2016) egg storage: air cell < 4 mm after 28 d at 6 C; > 6 mm after 25 d at 22 C; mass loss 1 / 2.58 / 4.54 g at 6 / 15 / 22 C over 28 d. \\
\texttt{ro} & Analele Univ. Oradea 2011 (Usturoi): fresh egg specific gravity 1.078-1.097; 1.040-1.059 after >30 d. \\
\texttt{ab} & Alberta Agriculture: infertile blastodisc 2-3 mm, fertile blastoderm 5-9 mm (snippet). \\
\texttt{unsw} & A Laboratory Manual of Vertebrate Embryology (1947), UNSW embryology: concentric yolk layers, latebra, neck of latebra to the nucleus of Pander under the germinal disc; two twisted chalazae. \\
\texttt{eng} & Engormix 'Protoporphyrin in shell and cuticle': 13-20 \% of brown pigment in the cuticle, $\sim$80 \% in shell. \\
\texttt{dsm} & DSM YolkFan (formerly Roche): 15 shades since 1965, 16th added 2016. \\
\texttt{cohen2025} & Cohen et al. 2025, Communications Physics (MIT; \url{physics.aps.org/articles/v18/100):} 180 drops onto a hard surface from 8-10 mm: side-on 5 \% (8 mm) / $\sim$45 \% (10 mm) cracked, vertical 55 \% / $\sim$80 \%; static breaking $\sim$45 N in both orientations. \\
\texttt{bps1976} & British Poultry Science 1976 (doi 10.1080/00071667608416318): some well-formed eggs crack from drops of 3.3 mm onto a stiff flat surface; typically a single hair crack round the equator or to a pole. \\
\texttt{alphafoodie\_\allowbreak{}boil} & \url{alphafoodie.com/how-long-to-boil-eggs:} hot start, medium eggs: 4.5 min white just set/yolk very runny, 6 jammy, 7 barely runny, 8-9 set but creamy, 10-11 hard, >14 overcooked. \\
\texttt{kenji\_\allowbreak{}boil} & Lopez-Alt (NYT 2019, via summaries): fridge eggs into boiling water: 6 min liquid yolk, 8.5 fudgy, 11 just firm; hot start peels cleanly $\sim$90 \% vs $\sim$50 \% cold start. \\
\texttt{egginfo\_\allowbreak{}crack} & British Lion egginfo.co.uk 'How to crack an egg': rounded corner (bowl or pan edge), light tap; too hard on a flat surface and you get egg everywhere. \\
\texttt{atk\_\allowbreak{}egg\_\allowbreak{}on\_\allowbreak{}egg} & America's Test Kitchen 'The best surface for cracking an egg is another egg': only the weaker egg cracks, leaving a clean divot; the membrane holds the shards. \\
\texttt{atk\_\allowbreak{}egg} & America's Test Kitchen 'Egg cracking showdown: bowl vs counter': rim tears the membrane and splits cleanly but leaves shards that can puncture the yolk; counter gives larger pieces, membrane often intact. \\
\texttt{wh\_\allowbreak{}separate} & wikiHow 'Separate an Egg': shell-to-shell transfer $\sim$3x or hand; cold eggs; third bowl per egg. \\
\texttt{cs42} & Cummings \& Schroeder 1942, Calif. Avocado Soc. Yearbook 27:56-64 (avocadosource CAS\_1942\_PG\_056-064): exocarp (cuticle, epidermis, 1-3 hypodermal cells, parenchyma, stone-cell layer), mesocarp with oil idioblasts and 6 vascular strands, thin endocarp, two-layer papery testa, two cotyledons, embryonic axis. \\
\texttt{uspto\_\allowbreak{}pp14239} & USPTO PP14239 (Hass comparison): skin 1.6 mm, fruit 204 g, seed 28.4 g, seed:skin:flesh 14:13:73. \\
\texttt{clark2007} & Clark et al. 2007 Postharvest Biol. Technol. 46:119 (phtnet mb409): fruit density 970-1018 kg/m3. \\
\texttt{landahl} & Landahl \& Terry (Cranfield): Hass whole-fruit compression 304.1 $\to$ 2.1 N over 6 d at 18 C. \\
\texttt{ashton2006} & Ashton et al. 2006 J. Agric. Food Chem. 54:10151: Hass skin L*/C*/h over ripening; flesh in three equal bands dark green / pale green / yellow; cyanidin-3-O-glucoside in ripe skin. \\
\texttt{farinas2021} & Farinas et al. 2021 J. Food Eng. 289:110266: flesh slice modulus 2.29 $\to$ 0.16 MPa with ripening. \\
\texttt{hia2018} & Hort Innovation 2018 'Avocado squeezing is bad': 10 N thumb press bruises firm-ripe fruit; shoppers 3-30 N; bruises show after $\sim$24 h. \\
\texttt{farley2019} & Farley et al. 2019 Am. J. Emerg. Med.: $\sim$50,413 US avocado knife injuries 1998-2017, mostly the non-dominant hand while pitting. \\
\texttt{ucd\_\allowbreak{}avocado} & UC Davis Postharvest Technology Center, avocado produce facts \url{(postharvest.ucdavis.edu/node/7115):} mature-green 5-13 C by cultivar; ripe fruit 2-4 C; ethylene at 20 C $\to$ ripe in 3-6 d. \\
\texttt{marques2011} & Marques et al. (2011) VII World Avocado Congress (avocadosource WAC7 MarquesJose2011b): Hass at the softening stage stored up to 10 d at 5-8 C with low risk of quality loss at eating ripe; avoid < 5 C once past the rubbery stage. \\
\texttt{ripe\_\allowbreak{}window\_\allowbreak{}guides} & Produce guides (axonify produce knowledge; UC Davis-hosted consumer page): a stage-4/5 Hass stays ripe $\sim$24-48 h at room temperature before overripening (weak sources). \\
\texttt{atk\_\allowbreak{}browning} & America's Test Kitchen 'How to keep avocado from turning brown': vacuum sealing kept halves green 5 d; plastic wrap + fridge minimal discolouration; lemon / lime juice or ascorbic acid slow browning but taste sour; keeping the pit in does not protect the cut face (informal kitchen test). \\
\texttt{afm\_\allowbreak{}howto} & Avocados From Mexico / freshfruitportal how-to: cut lengthwise around the seed, twist apart; remove seed with a spoon (knife strike discouraged); score to dice; peel from the small end. \\
\texttt{usda\_\allowbreak{}handling} & USDA ARS egg safety handling (ars.usda.gov .../eggsafety\_handlingeggs.pdf): gently remove a stuck egg; cook cracked eggs thoroughly. \\
\texttt{masterclass\_\allowbreak{}crack} & MasterClass 'How to crack an egg' (snippet): hold firm but gentle. \\
\texttt{cookscountry\_\allowbreak{}storage} & Cook's Country egg storage \url{(cookscountry.com/how_tos/5900-egg-storage):} carton, middle/back shelf 35-40 F; door $\sim$45 F. \\
\texttt{eggs\_\allowbreak{}ca} & Egg Farmers of Canada (eggs.ca/?p=15718): crack on a flat surface; crack into a separate bowl first. \\
\texttt{atk\_\allowbreak{}crack\_\allowbreak{}separate} & America's Test Kitchen 'How to crack and separate eggs' (articles/1126): cold eggs, flat surface, three bowls, half shell lifts yolk specks. \\
\texttt{egginfo\_\allowbreak{}separate} & British Lion egginfo.co.uk 'How to separate egg whites': shell-to-shell, hand, bottle methods. \\
\texttt{foodrepublic\_\allowbreak{}separate} & Food Republic 'Ways to separate eggs' (1769095): methods and separator tools. \\
\texttt{incredibleegg\_\allowbreak{}boil} & Incredible Egg 'How to make hard-boiled eggs': cold start, stand 9/12/15 min (M/L/XL); refrigerate eggs 7-10 d before; never microwave in shell. \\
\texttt{milkstreet\_\allowbreak{}peel} & Milk Street 'How to peel hard-cooked eggs' (2020): steaming 12 min, 5 min ice bath, shake method $\sim$60 \% on warm eggs. \\
\texttt{atk\_\allowbreak{}scrambled} & America's Test Kitchen 'Perfect scrambled eggs' (recipes/6604) and '6 tips' (articles/3118): fork-beat, half-and-half, medium-high then low. \\
\texttt{incredibleegg\_\allowbreak{}scrambled} & Incredible Egg basic scrambled eggs: medium heat, push curds, off while glossy. \\
\texttt{pepin\_\allowbreak{}pbs} & Jacques Pepin classic scrambled eggs (PBS American Masters): stop the cooking with reserved egg and cream. \\
\texttt{loveandlemons\_\allowbreak{}fried} & Love \& Lemons 'Fried egg': 1 tsp fat per egg, lid 2-3 min; over easy 15 s second side. \\
\texttt{incredibleegg\_\allowbreak{}fried} & Incredible Egg basic fried eggs: steam-baste with 1 tsp water and a lid, or butter-baste. \\
\texttt{kenji\_\allowbreak{}baste} & Lopez-Alt fried eggs via Tasting Table (1346566): baste with shimmering oil $\sim$1 min, crisp lacy edges. \\
\texttt{ishs\_\allowbreak{}1119\_\allowbreak{}22} & ISHS Acta Hortic. 1119\_22: avocado firmness devices (Sinclair iQ, firmometer, analogue meter, hand squeeze) are non-linearly related -- readings are device-specific. \\
\texttt{hirte\_\allowbreak{}2012} & Hirte A. (2012) Extending crispness - controlling water migration in bread. PhD thesis, Wageningen (edepot.wur.nl/221072): crispy roll $\sim$90 x 50 mm; crust 1-3 mm (cut-off crust 1.0-1.5, $\sim$1.3 mm); crumb porosity 0.73-0.81 (uCT 6 um), cell walls 20-61 um; crumb aw $\sim$0.97; crust water 9.0 \% at 30 min, 13.9 \% at 1 d (40 \% RH, 22 C); critical crust water 10.1 \% reached after 2.1 h; crust break $\ge$ 3 N (1 x 10 mm blade); bake-off mass loss 2.4-4.1 \%, 4 h storage 1.9-2.9 \%. \\
\texttt{lagrain\_\allowbreak{}2012} & Lagrain et al. (2012) Food Biophysics 7:173 (KU Leuven author manuscript): pan-bread crumb E $\sim$20 kPa (5 h) $\to$ $\sim$70 (48 h) $\to$ $\sim$125 (96 h) $\to$ $\sim$157 kPa (168 h) (read off figures), Poisson 0.103, loss factor 23 $\to$ 10 \%; mean 2D cell area 0.32 mm2, 83 cells/cm2; open porosity 0.81-0.89; Gibson-Ashby E/Es = (rho/rho\_s)\^{}2; crumb moisture 0.46 $\to$ 0.42 g/g. \\
\texttt{cardone\_\allowbreak{}2020} & Cardone et al. (2020) Foods 9:260 (PMC7143436): Avrami crumb firming, control wheat bread, n=1: A0 6.48 N, Ainf 66.01 N, k 0.186 (per day; unit to verify); crust L*$\sim$47, a* 9, b* 25. \\
\texttt{chaple\_\allowbreak{}2023} & Chaple et al. (2023) LWT 174:114326 (PMC9883616): pan bread crust L*a*b* 47.6/17.05/32.2, crumb 72.06/0.94/15.37. \\
\texttt{kurek\_\allowbreak{}2016} & Kurek et al. (2016) Food Technol. Biotechnol. 54(1) (PMC5105627): wheat rolls from 60 g dough, specific volume 6.4 (3 h) $\to$ 5.9 ml/g (48 h); moisture 47.8 \%; crust browning index 50.4. \\
\texttt{musiienko\_\allowbreak{}2025} & Musiienko et al. (2025) Front. Nutr. (PMC12497575): round buns 86-92 x 50-55 mm, 197-263 cm3. \\
\texttt{testde\_\allowbreak{}2021} & Stiftung Warentest 'Aufbackbroetchen' (test.de 4206027): rolls 50-90 g; target crust 'hellgoldbraun bis dunkelgoldbraun'; crispness lost after 4 h, crumb unchanged; score raises the crust share. \\
\texttt{ireks\_\allowbreak{}rmt} & IREKS Kompendium, RMT-Backversuch Broetchen: 1000 g T550 flour, 50 g yeast, 15 g salt, 10 sugar, 10 fat $\to$ 30 rolls; proof 25 min 32 C; bake 20 min at 230 C with steam. \\
\texttt{leitsaetze\_\allowbreak{}1993} & Leitsaetze fuer Brot und Kleingebaeck (1993): Kleingebaeck $\le$ 250 g; Weizenbroetchen $\ge$ 90 \% wheat flour (2021 revision: verify). \\
\texttt{giovanelli\_\allowbreak{}2021} & Giovanelli \& Cappa (2021) Foods 10:417 (PMC7918450): browning starts > 105-120 C at aw < 0.4-0.7; crust 140-180 C at the end of baking; browning index 35.5 (200 C) / 55.9 (225 C). \\
\texttt{soleimani\_\allowbreak{}2014} & Soleimani Pour-Damanab et al. (2014) J Food Sci Technol 51:3439 (PMC4571265): crust thickness grows with bake time; colour bands outer crust L* < 51.6, inner crust 51.6-62.3, crumb 62.3-71.2. \\
\texttt{kaur\_\allowbreak{}2026} & Kaur et al. (2026) Plant Biotechnol. J. (PMC13278543): home toaster 3 / 3.5 / 4 / 4.5 min, 'starting to look burned' at 4.5 min. \\
\texttt{utopia\_\allowbreak{}aufbacken} & \url{utopia.de} 'Broetchen aufbacken': oven 150-170 C for 5-10 min after dampening; toaster rack highest setting 2-3 min, flip, 1-2 min. \\
\texttt{oekotest\_\allowbreak{}storage} & \url{oekotest.de:} keep rolls in a paper bag at room temperature; plastic softens the crust; freeze fresh, do not refreeze. \\
\texttt{wikihow\_\allowbreak{}slice} & wikiHow 'Slice Bread' (via kiwix mirror): serrated knife, roll flat, hand flat on top, saw from the centre; stop short for a hinge. \\
\texttt{tastingtable\_\allowbreak{}bagel} & Tasting Table 'safest way to cut a bagel': flat side down, palm on top, serrated knife horizontally; halfway, stand on the uncut edge and finish downwards. \\
\texttt{luo\_\allowbreak{}2026} & Luo \& Yuan (2026) arXiv 2607.04367, robot food cutting: sawing at -30 N down-force, 2 cm/s, $\pm$2 cm strokes; carrot = serrated-knife food (chopping failed, sawing succeeded). \\
\texttt{zhang\_\allowbreak{}2024} & Zhang et al. (2024) Foods 13:397 (PMC10855756): rolls (80 g dough, 4 C) crumb hardness $\sim$1.5 N (3 d) $\to$ $\sim$4.0 N (7 d). \\
\texttt{eu1333\_\allowbreak{}2011} & Commission Implementing Reg. (EU) 1333/2011 bananas: min length 14 cm along the convex face from the blossom end to where the peduncle joins the crown; min grade 27 mm (transverse thickness); skin defects Extra $\le$ 1, I $\le$ 2, II $\le$ 4 cm2. \\
\texttt{ciro05} & Ciro Velasquez et al. (2005) Rev. Fac. Nal. Agr. Medellin 58:2975 (Cavendish Valery, 24 C, 12 d): outer-curve length 23.3 $\pm$ 1.13 cm; diameter 34/36/34 mm at 25/50/75 \%; peel puncture 13.6-15.2 N (2 mm probe); pulp penetration 7.20 $\to$ 3.41 N; pulp E compression 1.75 $\to$ 1.09 MPa, bending 5.16 $\to$ 2.65 MPa; whole fruit fracture 129 N (bending, 14 cm span), 121 N (compression). \\
\texttt{mal22} & Malikongwa et al. (2022) Plant Science Today 9:364 (Grand Naine): length 19.0 cm; pedicel 17.3 x 10.15 mm; peel 2.56 mm; pulp 62.3 g / peel 23.2 g; IPGRI: slightly ridged, blunt tip, persistent style, slightly curved, cream pulp. \\
\texttt{chiw23} & Chiwate et al. (2023) Curr Agric Res J 11(1) (Cavendish, 22-27 C): day 1/4/6/8/10 mass 105/101/93/87/82 g, pulp:peel 1.69..1.78, Brix 4.5/6.8/8/16.7/23.7, firmness 44.8/36.8/34.2/28.1/22.1 N. \\
\texttt{bant24} & Bantayehu \& Alemayehu (2024) Heliyon 10:e37143: Grand Naine 102 g, pulp:peel 2.18. \\
\texttt{amn12} & Amnuaysin et al. (2012) Trop. Nat. Hist. 12:127 (Hom Thong AAA): epidermis + cuticle, outer ground tissue (chloroplasts, starch), inner ground tissue with air spaces growing on ripening, bundles ringed by laticifers; peel 1.1-3.8 mm, thinning on ripening. \\
\texttt{guapo22} & Guapo \& Ferreira (2022) Cadernos UniFOA 17(50): 3 carpels/locules; 'strings' = innermost mesocarp strands of parenchyma with two small vascular bundles and secretory spaces (not pure phloem). \\
\texttt{promusa} & ProMusa 'Banana flowers': tricarpellate inferior ovary, 200-300 ovules on a central placenta; the black dots in the centre are the ovule remnants. \\
\texttt{khanal23} & Khanal et al. (2023) Postharvest Biol Technol: outer peel layers are the mechanical backbone; peel splits axially in the fruit corners (ridges); peel fracture 53.0 $\to$ 2.1 N on ripening (abstract). \\
\texttt{ringer21} & Ringer \& Blanke (2021) J Food Meas Charact / Tropentag 2021 (German supply chain): peel hue R2..R7 \url{98.7/97.3/92.7/89.4/87.5/82.0} deg, b* 56-72; peel water 90.5 \% ($\le$ R5) $\to$ 82.9 \% (R7); Shore 82 $\to$ 42. \\
\texttt{hou15} & Hou et al. (2015) IJABE 8(6):99: ethylene-treated AAA at 20 C, stage 1 $\to$ 7 in 7 days; peel L*, b* peak at stage 5, a* rises. \\
\texttt{ma22} & Ma et al. (2022) Curr Res Food Sci 5:1808: stage day bands RS1 1-4, RS2 5-6, RS3 7-8, RS4 9-10, RS5 11-12, RS6 13-15, RS7 > 15 d. \\
\texttt{markon} & Markon banana colour chart: 1 green, 2 green trace yellow, 3 more green than yellow, 4 more yellow than green, 5 yellow green tips, 6 all yellow, 7 yellow flecked brown. \\
\texttt{wflo} & WFLO Commodity Storage Manual, Bananas (2018): chilling below 13 C (pulp); ripening rooms 14-18 C; specific heat 3.56 kJ/kg K; moisture 74.3 \%. \\
\texttt{mbeguie09} & Mbeguie-A-Mbeguie et al. (2009) J Exp Bot 60:2021: pedicel rupture force falls > 2x between day 1 and 3 (Grande Naine, 20 C); no abscission zone (ketsa25). \\
\texttt{imsabai} & Thai finger-drop report (TSRI BGJ4380011): finger-drop pull resistance $\sim$5 kg ($\sim$49 N) on day 1 (Hom Thong / Namwa). \\
\texttt{bugaud14} & Bugaud et al. (2014) PBT 87:113: Grande Naine no visible bruise up to 200 mJ; susceptibility rises with ripening (abstract). \\
\texttt{col20} & Colombian Musaceae ripening (PMC7984252), Gros Michel day 0/4/6 L* 39.9/45.0/56.4, a* -26/-12/+13.1, b* 40.9/62.1/85.4. \\
\texttt{kim22} & Kim, Ohmura, Kuniyoshi (2022) arXiv 2203.09749: dual-arm banana peeling, 40.9 \% consecutive success (PeelNeck 95 \%, PeelRight 59 \%, PeelLeft 64 \%); neck too stiff on fresh fruit, pre-cut. \\
\texttt{banana\_\allowbreak{}howto} & Tasting Table (1552037), Chiquita blog: peel from the stem (snap the neck) or from the tip ('monkey method': pinch the tip); rounds $\sim$6 mm (3-6 mm, MyPlate); lemon juice against browning; hang, keep apart from apples/avocados; refrigerate only when ripe. \\
\texttt{sadka19} & Sadka et al. (2019) Front. Plant Sci. 10:1167: hesperidium; flavedo (oil glands), albedo, endocarp segment membrane, juice vesicles on stalks from the endocarp; 3 bundles per segment (1 dorsal, 2 septal) ending in the albedo; central axis (columella). \\
\texttt{smith18} & Smith, Ebrahimi, Ghosh, Dickerson (2018) PNAS 115:E5887 (PMC6042112): oil jets on bending the peel 8.47 $\pm$ 4.03 m/s, 102 um; reservoir $\sim$1 mm3 under an 84 um flavedo cap, burst at 30-140 kPa; flavedo E $\sim$103 MPa (30-300), albedo 0.79 $\pm$ 0.58 MPa; albedo 3.9 mm in the FE model. \\
\texttt{jentzsch22} & Jentzsch et al. (2022) Front. Mater. 9:979151: 'Lane Late' peel 6.2 mm, density 746 kg/m3, compressive E $\sim$0.02 MPa, Poisson -0.05, 0.51 MPa at 50 \% strain. \\
\texttt{ucr\_\allowbreak{}citrus} & UCR Givaudan Citrus Variety Collection data: Washington navel (ripe) 162-335 g, 61-89 x 61-88 mm, rind 4.1-6.4 mm, juice 45-54 \%, TSS 9.3-14.1, acid 0.49-0.77 \%, 0-6 seeds per 10 fruit; Olinda Valencia 113-299 g, L/W 1.0-1.1, rind 3.5-6.2 mm, juice 36.9-56.1 \%, 1.2-6.6 seeds per fruit. \\
\texttt{edis891} & UF/IFAS HS891 Dooryard Citrus Guide: Valencia and navel 0-6 seeds; navel = rudimentary secondary fruit at the blossom end; navel juice turns bitter within several hours. \\
\texttt{citrus\_\allowbreak{}patents} & USPTO PP21535 / PP27829 (sweet orange): 9.6-10.6 and 10-11 segments; PP27829 208 oil cells/cm2; PP21224 albedo 2-5 mm; PP6047 navel peel 5-5.1 mm. \\
\texttt{knight01} & Knight, Klieber, Sedgley (2001) Ann. Bot. 88:1039: 8,000-12,000 oil glands on a mature Washington navel (abstract) $\to$ $\sim$40-60 /cm2 on $\sim$200 cm2. \\
\texttt{eu543\_\allowbreak{}citrus} & Reg. (EU) 543/2011 Annex I B 2 (oranges): min diameter 53 mm (size codes 0-13); min juice 33 \% navel / 35 \% others; light green on $\le$ 1/5 of the surface (superseded by 2023/2429: verify). \\
\texttt{orange\_\allowbreak{}mech} & Khoshnam et al. 2022 Iran. J. Biosyst. Eng. 53:109 (punch 15.5-79 N, 1.5-7.9 N/mm, abstract); Dadvar et al. 2015 (Valencia rupture 140-221 N, abstract); Montero et al. 2012 Sci. Hortic. (oil glands rupture under 10-250 N sphere compression; Valencia no oleocellosis). \\
\texttt{ritenour13} & Ritenour et al. ASHS 2013: peel-removal rig (16 mm strip, 75 mm): Valencia hard-peel (strips break), navel medium-easy; cold fruit (5 C) needs more force and tears more; no N values published. \\
\texttt{orange\_\allowbreak{}juice} & BSDA QUID: juice SG 1.045 at 11.2 Brix; US4938985: single-strength viscosity $\sim$4-6 mPa s at 8 C; alphafoodie: 60-80 mL per medium orange by hand, roll 10-15 s, halve at the equator, ream. \\
\texttt{orange\_\allowbreak{}storage} & UC Davis Produce Facts (Arpaia \& Kader): 3-8 C, 90-95 \% RH, up to 3 months; chilling (pitting) after 3-6 weeks at $\sim$4 C (UF PHNL194). \\
\texttt{orange\_\allowbreak{}howto} & instructables 'Easy Orange Peeling' (score 6-8 mm deep round the navel, spiral to the stem); kitchprep; The Daily Meal / A Couple Cooks: supreme -- peel and pith off, cut inside each membrane; zest the flavedo only (albedo and membranes bitter, limonin). \\
\texttt{cci} & Citrus colour index CCI = 1000 a*/(L* b*) (Jimenez-Cuesta 1981); market CI target $\sim$ +10..+20 (snippet). \\
\texttt{test25\_\allowbreak{}toast} & Stiftung Warentest 'Toastbrot im Test' (23.04.2025, \url{test.de} 4421655): baked $\sim$30 min under a lid (thin soft crust), cooled in a clean room, sliced and bagged; most packs have no heels (organic packs: 2); best-before 12-20 d; no classic preservatives needed. \\
\texttt{toast\_\allowbreak{}pack} & Open Food Facts / REWE / Friedrich Verlag (snippets): 500 g pack, 20 slices of $\sim$25 g, pack $\sim$0.25 m long. \\
\texttt{nuss90} & Nussinovitch, Roy, Peleg (1990) Cereal Chem 67:101: commercial white sliced bread, crumb 0.19 g/cm3; tensile failure 0.31 / 0.39 N (12 mm neck, 10 / 50 mm/min) at 9.2-9.7 mm; compressed bread rarely fails. \\
\texttt{sehn20} & Sehn \& Steel (2020) J Food Sci Technol 57:557 (PMC7016051): white control pan loaf firmness 1.98 / 3.62 / 4.18 / 5.64 N at day 1 / 5 / 9 / 13 (25 C, PE bag); crumb moisture 39.35 $\to$ 36.01 \%. \\
\texttt{ishida14} & Ishida \& Steel (2014) Food Sci Technol Int: Brazilian commercial white pan bread: 4.72 ml/g, crumb 37.0 \% moisture, firmness 276 gf day 1, crumb L*a*b* 74.73 / 0.37 / 15.51 (abstract). \\
\texttt{toaster26} & Stiftung Warentest toaster test (2026): 80 s (fastest) to > 2.5 min to the best toast; higher settings = longer time, not more heat; browning scored as browned area. \\
\texttt{toast\_\allowbreak{}misc} & EU 2017/2158 acrylamide benchmark 50 ug/kg (wheat soft bread); FSA 'Go for Gold'; Daily Meal: 'perfect toast' 216 s at 5/6, 14 mm; butter $\sim$0.068 g/cm2 (Wedzicha); fridge stales $\sim$6x faster (popular, weak); Verbraucherzentrale: no fridge, freeze slices airtight. \\
\texttt{tins} & Runex 500 g bread tin: top inside 280 x 100 mm, depth 90 mm; hobby lidded toast tin 36 x 11 x 8 cm (2 loaves). \\
\texttt{mc20} & Massana-Codina et al. (2020) Front Plant Sci 11:1287 (PMC7468465): periderm = phellem / phellogen / phelloderm; phellem 46-84 um, 6-10 cell layers (five European cultivars). \\
\texttt{mag25} & Magwaza et al. (2025) Front Plant Sci 16:1595828: lenticels $\le$ 12 /cm2 (small tubers); apical buds $\sim$7x the lateral density, 6-9 apical buds; respiration 1.9-2.75 mL CO2/kg h. \\
\texttt{bsl21} & Bundessortenamt Beschreibende Sortenliste Kartoffel 2021: Belana oval yellow-skinned, very flat-flat eyes; Gala round-oval, flat eyes; Linda long-oval, flat-medium eyes; flesh medium-yellow. \\
\texttt{potato\_\allowbreak{}mech} & E 3.12 MPa (Cagnolati 2018; Celik 2019), 3.15 MPa / nu 0.45 / yield 1.18 MPa (Zhao 2023), 4.6-5.6 MPa, fracture 1.13-1.40 MPa (Polish clones, SV); fracture strain $\sim$0.3; toughness $\sim$212 J/m2 (1996, SV); blackspot above $\sim$50 g impact (Hyde 1991, SV). \\
\texttt{kk18} & Krishna Kumar et al. (2018) J Food Meas Charact: raw potato 1053 kg/m3, k 0.56 W/m K, alpha 0.15e-6 m2/s, initial freezing -1.8 C (abstract); ETB: k 0.55, cp 3.43 kJ/kg K. \\
\texttt{potato\_\allowbreak{}cook} & Lee (2017) Food Eng Prog 21:303 review: first-order softening, k(100 C) 0.28-0.40 /min, Ea 103-145 kJ/mol (Harada 1985); gelatinisation onset 61, peak 67.5, end 76-82 C (Witczak 2013, SV); PME firming at 60-70 C (Gonzalez-Martinez 2004). \\
\texttt{gritzer} & Serious Eats (Gritzer): cold start in 1-2 \% salted water; diced $\sim$15-20 min, large whole up to $\sim$1 h; done when a skewer passes without resistance; German guides 20-25 min whole from cold. \\
\texttt{potato\_\allowbreak{}storage} & <8 C cold-induced sweetening; home 7-13 C dark; $\sim$3 \% weight loss in month 1; greening under retail light 1-2.5 d at 20 C (Nofima), from day 4 / 2 at 10 / 15 C (JSHS 1990); BfR: remove green parts and eyes generously. \\
\texttt{potato\_\allowbreak{}peel} & Peeler vs knife study (SV): hand peeler $\sim$1.1 mm, $\sim$9 \% waste; peeled slices under light: glycoalkaloids 2 $\to$ 74 mg/kg in 48 h (OSU EM 9407). \\
\texttt{ince16} & Ince, Cevik, Vursavus (2016) IJABE 9(6):200 (Jadelo F1): 72 x 70 x 53 mm, 165 g; plate firmness initial G/P/R 15.1 / 6.0 / 4.5 N/mm, average 13.2 / 4.7 / 3.1; 2 mm puncture skin rupture 6.92 / 3.99 / 2.74 N; a*/b* bands green -0.59..-0.47, breaker ..-0.27, turning ..0.08, pink ..0.60, light red ..0.95, red ..1.21 (after Batu 2004). \\
\texttt{fdc\_\allowbreak{}tomato} & USDA FoodData Central 170457: small whole 61 mm 91 g, medium 66 mm 123 g; water 94.5 g/100 g; slices 1/4 in 20 g, 1/2 in 27 g; wedge (1/4 medium) 31 g. \\
\texttt{bn05} & Bargel \& Neinhuis (2005) J Exp Bot 56:1049: cuticle 9.9-15.7 um (red); skin E 252-293 MPa, failure $\le$ 25-28 MPa at $\sim$9-10 \% strain (red). \\
\texttt{meta19} & Metabolites 9:93 (2019, Ailsa Craig red): pericarp 5.3 mm; pericarp $\sim$50 \%, exocarp 9.6 \%, columella + placenta 10.0 \%, locular tissue 23.3 \% of FW. \\
\texttt{ucd\_\allowbreak{}rost} & UC Davis Rost lab tomato anatomy: skin, pericarp, septa, columella, placenta, locular gel; most cultivated non-cherry types 4-5 locules; small types often 2 (Wikipedia). \\
\texttt{cfr\_\allowbreak{}tomato} & 7 CFR 51.1860 colour classes: green, breakers ($\le$10 \%), turning (10-30 \%), pink (30-60 \%), light red (60-90 \%), red (> 90 \%). \\
\texttt{wflo\_\allowbreak{}tomato} & WFLO Commodity Storage Manual, Tomatoes: ripen 18.3-20 C; chilling < 12.8 C; ripe 7-10 C 3-5 d; freezing -0.5 C; cp 4.08 kJ/kg K; mature-green seeds pushed aside by the knife (gel liquefied). \\
\texttt{zhang16} & Zhang et al. (2016) PNAS 113:12580: 8 d at 5 C cuts volatiles by 65 \%, not recovered after 3 d at 20 C. \\
\texttt{tomato\_\allowbreak{}mech} & Sirisomboon 2012 (whole stiffness 2000-10000 N/m, peel 70-90 \% of firmness); Safitri 2022 (Hertz E 0.58-1.97 MPa, SV); Li 2017 (FEM: damage from $\sim$3 \% strain; 6.5 N at 10 \%); Kultongkham 2021 (bio-yield 2.57 N, SV); Lana 2004 (slice firmness Ea 87.5 / 94.8 kJ/mol). \\
\texttt{tomato\_\allowbreak{}howto} & Wikibooks Tomato Concasse: cone-core, X 2.5 cm, blanch 30-45 s, ice water, peel, seed, dice 0.75 cm; Jessica Gavin: serrated or very sharp knife, slices 6-13 mm, 8 wedges; room temperature. \\
\texttt{hb66\_\allowbreak{}onion} & USDA AH-66 onion chapter (Adamicki): bulb = enlarged leaf bases + compact stem; 0 C, 65-75 \% RH, 6-9 months; rest 4-6 weeks; curing 3-5 \% loss; respiration 3 (0 C) .. 8 (20 C) mg CO2/kg h; freezing -0.8 C; medium 5.1-8.3 cm. \\
\texttt{bah04} & Bahnasawy et al. (2004) J Food Eng 62:255 (yellow Giza 20): 51.7 x 57.1 mm, 78.7 g, density 1.04 g/cm3, crushing 400-780 N, puncture 27.6-45.5 N, friction 0.67-1.34, rolling angle 20-31 deg (base). \\
\texttt{kim\_\allowbreak{}heb20} & Hortic Environ Biotechnol (2020), 45 cultivars: 7.6 $\pm$ 0.2 fleshy scales, 2.3 growing centres, scale thickness 5.5 $\pm$ 0.2 mm (abstract). \\
\texttt{wu25} & Wu et al. (2025) arXiv 2505.06016 'Droplet outbursts from onion cutting': layer $\sim$3 mm; tissue E 2.2 MPa (skin on) / 1.4 (off); blunt / fast knives up to $\sim$40x / $\sim$4x more droplets; chilling did not change droplet speed. \\
\texttt{wei01} & Wei, Lintilhac, Tanguay (2001) Plant Physiol 126:1129: adaxial epidermal peel 120 um, E 3.5-8.0 MPa (turgor 0.14-0.54 MPa), nu 0.18-0.30; strips break at $\sim$40 / $\sim$20 gf (long./trans.). \\
\texttt{jafari16} & Jafari Malekabadi et al. (2016) J Agric Machinery 6:114: whole yellow onion E 2.03-5.45 MPa, nu 0.26-0.45. \\
\texttt{lokke12} & Lokke et al. (2012) Sensors 12:16060: lachrymatory factor highest at $\sim$1 min after chopping, gone by $\sim$10 min; Imai 2002 Nature 419:685 (LF synthase). \\
\texttt{onion\_\allowbreak{}howto} & America's Test Kitchen 'How to cut an onion' (2016): halve through the root, peel, keep the root on, horizontal then vertical cuts, claw grip; mince $\le$ 3 mm, fine 3-6, medium 6-13, coarse 13-19 mm; rings: trim a flat base; HWQ: 2-3 horizontal cuts, remove tough layers; UNL: chill 30 min. \\
\texttt{onion\_\allowbreak{}misc} & Gracie 2012 (6th true leaf forms the outer skin; skins tear easily); Karate 2025 (cured neck 7.4-10.9 mm); Alamar 2020 (lateral buds up to 1/3 of mass; dormancy break week 8 at 1 C); frying colour a* 5.7 $\to$ 17.9, b* 14.2 $\to$ 29.3 (SV); caramelising 30-45 min (typical). \\
\texttt{agu20} & Aguilar-Hernandez et al. (2020) Foods 10:62 (PMC7824522), Fino / Verna: 100-184 g, ED 55-67 mm, PD 72.7-91.9 mm, ED/PD 0.6-0.8, peel 4.9-6.8 mm, 8.8-10.2 segments, 1.3-11.3 seeds, juice 31.7-44.3 \%, TA 56-75 g/L, TSS 7.8-10.5, peel L* 64-68, a* -7.5..3.9, b* 42-49, hue 88-99 deg. \\
\texttt{jen22\_\allowbreak{}lemon} & Jentzsch et al. (2022) Plants 11:991 (PMC9002614), Primofiore: 91 mm, 166 g, peel 5.26 mm, flavedo $\sim$1 mm (20 \%), peel 36.4 \% of mass, peel density 818 kg/m3, drop impact COR 0.31. \\
\texttt{ben26} & Bengisu et al. (2026) Biomimetics 11:451 (PMC13406879), retail lemon: puncture 41.2 N (5 mm), knife cutting 36.3 N, compression 165.8 N at 16 mm; peel 0.68 g/cm3, albedo 0.39 g/cm3; oil glands 1.25 /mm2. \\
\texttt{eu2023\_\allowbreak{}2429} & Delegated Reg. (EU) 2023/2429 citrus: lemons min juice 20 \%, min diameter 45 mm, size codes 0 (79-90) .. 7 (45-52 mm). \\
\texttt{lemon\_\allowbreak{}misc} & MAR22 Verna (PMC9407548): 93.7-98.1 mm, peel 6.7-7.1, seeds 1.3-2.3; LIA22 Eureka chilling < 10 C, 3.5 \% loss in 80 d; MIT20 degreening CCI -16.3 $\to$ -1.1; USDA FDC 167747 juice of 1 lemon 48 g, wedge 5.9 g; Wikipedia citric acid 5-6 \%, pH 2.2; zest = flavedo only (albedo bitter). \\
\texttt{dewiki\_\allowbreak{}sg} & de.wikipedia 'Salatgurke': normally three locules; trade classes 300-400 / 400-500 g. \\
\texttt{unece\_\allowbreak{}cuc} & UNECE FFV-15 (2017): Extra / I max 10 mm arc height per 10 cm, II 20 mm; weight spread 150 g ($\ge$ 400 g); EU 1677/88 (expired 2009): $\ge$ 250 g, $\ge$ 25 cm (250-500 g), $\ge$ 30 cm ($\ge$ 500 g); seeds soft. \\
\texttt{sakata11} & Sakata et al. (2011) J Japan Soc Hort Sci 80:420: greenhouse type puncture (3 mm) flesh 8.7 N, placenta 2.8 N, skin 9.6-9.8 N; 25-30 cm, $\sim$260 g. \\
\texttt{jah16} & Jahangiri et al. (2016) CIGR J 18:323: apparent modulus 0.85-1.39 MPa, failure stress 0.53-0.85 MPa, strain 0.45-0.53; firmness -49 \% in 9 d; stem end 7.53 vs blossom end 5.01 N. \\
\texttt{disect\_\allowbreak{}cuc} & Heiden et al. (2021) DiSECt (arXiv 2105.12244) Table A1: cucumber E 2.5 MPa, nu 0.37, 950 kg/m3; SliceIt! (2404.02569): 5 mm slices, stiffness jump cucumber $\to$ board. \\
\texttt{ucd\_\allowbreak{}cuc} & UC Davis Produce Facts, Cucumber: 10-12.5 C, 95 \% RH, < 14 d; chilling below 10 C after 1-3 d (7.2 C: 2-3 d); ethylene 1-5 ppm yellows; freezing -0.5 C; respiration 12-15 mL CO2/kg h at 10 C. \\
\texttt{cuc\_\allowbreak{}misc} & rett24 cuticle 4.0 um (1.1-10.7); ando15 peel layer 1-2 mm; wang22 skin a* -20.6, b* 27.5 (yellowing to b* 35 in 20 d); dhall10 wrapped 0.66 vs unwrapped 11.1 \% loss in 15 d at 12 C; shoji22 typical 350 g; wikibooks: fork scrape for wax, zebra peel, spoon out seeds, salt 30 min. \\
\texttt{leitsaetze\_\allowbreak{}fb} & Leitsaetze fuer Feine Backwaren (2025) 2.2.3.4: Plunder $\ge$ 30 kg butter (or equivalent fat) per 100 kg flour, flaky structure; croissant = rolled Plunder, often crescent-shaped; 2.1.1: butter in the name $\to$ no other fats. \\
\texttt{ste25} & Steinkellner, Kroll, Franke (2025) Curr Res Food Sci 10:101042 (PMC11995129): butter croissant $\sim$50 \% butter on flour, 3 single folds = 27 fat layers, specific volume 3.13 $\pm$ 0.66 cm3/g, pore surface 6.62 /cm; TPA firmness 5.16 N (1 h) $\to$ 19.2 N (24 h); bake 200/220 C 20 min; lamination 13-15 C. \\
\texttt{croissant\_\allowbreak{}misc} & Wikipedia en/de/fr Croissant: rolled 3.5-4 turns; optimum 16-50 fat layers (Ooms 2016); butter melts 32-34 C; final sheet 3.5 mm; bake 210-220 C. USDA FDC 174987: 28 / 42 / 57 / 67 g; water 23.2, fat 21.0 g/100 g. Weekend Bakery: egg wash 2 coats; reheat from frozen 180 C 8 min. \\
\texttt{wp\_\allowbreak{}baguette} & Wikipedia fr/en 'Baguette': $\sim$65 cm, 4-6 cm wide, 3-5 cm high, $\sim$250 g; tradition = flour, water, yeast/levain, salt (decret 93-1074); scores at 45 deg with a razor; stales faster than a boule; Grand Prix de Paris: 55-65 cm, 250-300 g, salt 18 g/kg. \\
\texttt{wp\_\allowbreak{}grigne} & Wikipedia fr 'Grigne': coupe parisienne = 5 medium overlapping cuts along the axis, $\sim$1 cm deep; the grignes open into raised lips ('oreilles'). \\
\texttt{dv2012} & Della Valle et al. (2012) J Cereal Sci 56:440: French breads (64), mean crust 1.7 mm; crust ends $\sim$160 C, $\sim$5 \% moisture. \\
\texttt{zou2024} & Zou et al. (2024) npj Sci Food (PMC11390739): baguette crust L*a*b* 57.5 / 10.5 / 34.9, crumb 66.4 / -1.7 / 15.7 (powders, SV); bake 230 C 30 min with steam. \\
\texttt{baguette\_\allowbreak{}misc} & King Arthur baguettes (232 C, 24-28 min, steam; paper bag overnight, freeze); van Vliet 2006 crispness lost within 2 h at 80 \% RH; Wikipedia Staling: fastest at $\sim$4 C, partly reversed at 60 C. \\
\texttt{he22} & He et al. (2022) Plants 11:985 (PMC9003265): blocky red cv. Gina 237 $\pm$ 37 g, 77.6 x 81.3 mm, wall 7.5 $\pm$ 0.7 mm (orange 6.9); cuticle 10.4 / 8.7 um; bite firmness 11.6 / 6.9 N, -81 \% after 15 d at 20 C. \\
\texttt{alt20} & Althaus \& Blanke (2020) J Imaging 6:122: German retail peppers 230-245 g; no stomata; at 17 C / 55 \% RH 2.9-4.8 g/d loss, shrivel from 6-8 \% (4-5 d); roughness Sa 2 / 7 / 24 um; gloss 450-500 $\to$ 80-90; non-climacteric. \\
\texttt{bar20} & Barbosa et al. (2020) Molecules (PMC7287789): pericarp L*a*b* green 33.7 / -8.6 / 16.9, red 38.8 / 31.4 / 27.8, yellow 61.4 / 5.0 / 50.0; whole-fruit shelf life 12-28 d. \\
\texttt{sha26} & Sharma et al. (2026) Sci Rep (DUS): locules 2-4, mostly 3; 1000-seed weight 4 to > 6 g. \\
\texttt{pepper\_\allowbreak{}misc} & Mensah 2026: 38-142 seeds per fruit; Afolabi 2023: whole-fruit 8 mm probe 43.5 N; Darre 2024 / Hanaei 2022: chilling below 7 C (7-10); EU 2023/2429 sweet peppers (classes, sizing); USDA ring (3 in x 1/4 in) 10 g; Harvey / Sweeper harvesting robots (36.9 s, 24 s cycles); Wikibooks: remove stem, core, seeds and white veins; char to loosen the skin. \\
\texttt{eu730\_\allowbreak{}1999} & Commission Reg. (EC) 730/1999 carrot marketing standard (mirrors UNECE FFV-10): main crop $\ge$ 20 mm or 50 g, Extra $\le$ 45 mm / 200 g; green/violet tops: Extra none, Class I $\le$ 1 cm ($\le$ 10 cm long) / $\le$ 2 cm (longer), Class II $\le$ 2 / 3 cm; firm, not woody, not forked. \\
\texttt{con22} & Conversa et al. (2022) Foods (PMC8750513): Nantes 'Presto' 135.4 g, 210.8 mm, 31.4 mm; 'CPT' 110.4 g, 182.7 mm, 30.3 mm; hand peeler removes $\sim$2 mm = 12.1 / 14.6 \% of fresh mass. \\
\texttt{jah18} & Jahanbakhshi et al. (2018) Food Sci Nutr 6:1898 (PMC6189625): 154.6 mm, 72.7 g; density 1.04 g/cm3; moisture 87.2 \%; 3-point bending Fmax 48.6 N at 21.9 mm, 0.50 J; blade shear 41.1 N; compression 71.9 N. \\
\texttt{vob24} & Vo-Bui, Rogers, Pegg (2024) R Soc Open Sci 11:230420 (PMC10805594): Nantes orthotropic constants after Stopa et al. 2017: xylem E 4.62-4.88 MPa, phloem 4.33-4.85 MPa, nu 0.47-0.48, G 1.47-1.66 MPa; E falls 1.32x in 7 d (fridge, zip bag); core/root diameter ratio falls crown $\to$ tip. \\
\texttt{peng14} & Peng, Tang, Barrett, Sablani, Powers (2014) J Food Eng 122:84: carrot softening, 2nd-order fractional conversion, k 0.035 (80 C) $\to$ 1.453 1/min (110 C), Ea 138.9 kJ/mol, Finf/F0 0.037-0.008; raw dice (12.7x12.7x6 mm, 70 \% strain) 274 N; xylem core 4-7 mm in dice. \\
\texttt{singh17} & Singh et al. (2017) Indian J Hort 74:56: 40 genotypes, root 30.55 mm, core 11.74 mm (ratio $\sim$0.38). \\
\texttt{hb66} & USDA Agric. Handbook 66 (2016), carrot chapter: 0-1 C / 98-100 \% RH 7-9 months; bagged 2-3 weeks at 3-5 C; freezing point -1.2 C; respiration 46-95 mg CO2/kg h at 20 C; Nantes prone to cracking and tip breakage; parenchyma cells $\sim$50 um; softening = turgor loss then pectin beta-elimination. \\
\texttt{herp99\_\allowbreak{}calbo00} & Herppich et al. (1999) PBT; Calbo (2000) Bragantia: fresh turgor $\sim$0.7 MPa; turgor firmness $\sim$0 below relative water content 0.95 ($\sim$5 \% water loss) -- the limp carrot. \\
\texttt{storage\_\allowbreak{}loss} & Snippets: unwrapped 20 C / 50 \% RH unacceptable after $\sim$3 d (US4919948); > 5 \% loss in 4 d bulk unsaleable (Semina 2010); 5-8 \% loss unsaleable. \\
\texttt{sim09\_\allowbreak{}cz95} & Simoes et al. (2009) PBT; Cisneros-Zevallos et al. (1995) J Food Sci 60:320: white blush on cut carrot = surface dehydration (reversible within $\sim$1-1.5 h) then lignification (irreversible after $\sim$15 h). \\
\texttt{pac16\_\allowbreak{}dry} & Paciulli et al. (2016) J Food Sci Technol; drying study IJFP 10.1080/10942912.2010.489209: raw carrot L*a*b* 50-55 / 25-29 / 30-40 (hue $\sim$50-54 deg). \\
\texttt{alv00} & Alvarez et al. (2000) J Texture Stud: SENB on carrot/celery/apple, K\_IC 10-50 kPa m\^{}0.5, fracture energy 40-400 J/m2. \\
\texttt{morpheus\_\allowbreak{}2024} & Ye et al. ICRA 2024 MORPHeus (arXiv 2404.06570): instrumented peeler, carrot 90.1 \% area peeled; skin and flesh colours hard to tell apart $\to$ force/vibration sensing. \\
\texttt{carrot\_\allowbreak{}cuts} & ATK 'How to cut carrots' (articles/3368), Leiths, Wikipedia culinary cuts: plank first so it does not roll; 5-7.5 cm lengths; julienne 3x3x40-50 mm, brunoise 3 mm, batonnet 6x6x50-60 mm, dice 6/12/19 mm; peel from the thick end, rotate after each stroke. \\
\texttt{carrot\_\allowbreak{}boil} & Evolving Table / Food \& Home snippets: 1/4 in rounds 4-5 min crisp-tender (+1-2 min soft); whole 10-15 min; not much more than 10 min for slices. \\
\texttt{diarch\_\allowbreak{}rows} & Botany: lateral roots arise beside each xylem pole $\to$ a diarch root has 4 rows of lateral roots (general text; carrot-specific confirmation in Esau 1940 Hilgardia 13:175 to obtain). \\\bottomrule\end{longtable}}
\section{Files and rebuild}
All paths are relative to \texttt{\textasciitilde/Documents/HouseholdTwin}.
\begin{Verbatim}[fontsize=\scriptsize,frame=lines]
food/params.py               parameter register (status per value) and SOURCES
food/<item>.py               analytic anatomy, part tree, asset spec per item
food/<item>_models.py        hidden-state models and beliefs (wave 1-3)
food/dynamics.py, sdt.py     egg/apple/avocado truth dynamics and beliefs
food/egg_crack.py, egg_cook.py, egg_open.py   egg crack, boil and opening models
food/schema.py               discrete SDT schema (egg, apple, avocado, roll, banana, carrot)
food/partgraph.py, usd_partgraph.py           closed meshes, cuts, USD part graph
food/egg_tasks.py, wave1_tasks.py, wave2_tasks.py   task designators
household/garmi_rig.py       GARMI articulation rig (IK, force-limited grasp)
scripts/build_food_assets.py assets/food/<kind>_<seed>/ (whole + cut states, sdt.usda, truth.usda)
scripts/import_apartments.py assets/apartments/ (EASE, GARMI apartment, GARMI robot)
scripts/sim_*.py             Isaac Sim benches -> test_reports/*.json
scripts/render_*.py          path-traced renders -> renders/
tools/*.py                   CPU studies, plots and generated docs
food/avocado_vision.py, avocado_belief.py    photo -> node estimates; particle-filter belief of the avocado node
food/avocado_peel.py, peel_monitor.py        peeling outcomes as USD states; vision-based peeling monitor
food/avocado_process.py, avocado_emulator.py, avocado_tasks.py   halving / pitting mechanics, emulator, designators
scripts/avocado_from_image.py, avocado_observe.py, avocado_process_demo.py, *peel*.py   CLI and studies
cram/avocado_sdt.py          node, competence questions, mental simulation, value of information (CRAM side)
cram/avocado_world.py, avocado_motions.py   Avocado(Fruit) annotation, pieces, PR2 motions
cram/avocado_pr2_demo.py     scripted PR2 scenario; cram/record_run.py -> cram/results/pr2_run.json
cram/avocado_kitchen.py      interactive kitchen for the CRAMERA viewer
report/build_report.py       this document (CPU only; cache in report/cache, figures in report/figures)
report/avocado_reasoning.py  the section on the avocado as the robot's belief state
cram/avocado_agent.py        goal-directed agent: effects to enforce / prevent, missions, lifetime-based decisions
food/avocado_lifetime.py     ripening, ripe window, decline, spoilage, browning under storage treatments
theory/avocado_bws.py, avocado_policies.py, axioms.py   research SDT, task programs, relational axioms and audit
theory/run_microtheory.py    held-out evaluation, ablations, misspecified plants -> theory/results/
theory/package.py, learning.py, prl.py   theory packages and gate, theory construction, programmable RL
report/avocado_theory.py     the section on the agent, the microtheory, theory construction and RL
report/build_avocado_microtheory.py   the standalone microtheory (report/avocado_microtheory.pdf)

E=~/Documents/LabAutomationTwin/env_cad/bin/python
I=~/Documents/LabAutomationTwin/env_isaacsim/bin/python
$E scripts/build_food_assets.py
$E tools/egg_crack_study.py && $E tools/egg_cook_study.py && $E tools/anatomy_doc.py
OMNI_KIT_ACCEPT_EULA=YES $I scripts/sim_egg_bench.py        (GPU, one job at a time)
OMNI_KIT_ACCEPT_EULA=YES $I scripts/render_food.py
$E report/build_report.py [--refresh]
\end{Verbatim}
\end{document}
